\chapter{Every change}\label{ch:changes}
All 392 commits in order, oldest first, each with its own account of what it changed and why. Messages in this repository are written as findings rather than as labels, so this chapter doubles as the project diary. Bodies are condensed to their opening argument; the full text is in \code{git log}.

\section*{2026-07-14 \textnormal{\small\textcolor{srlgrey}{--- 1 commit}}}
\addcontentsline{toc}{section}{2026-07-14}
\begin{longtable}{@{}p{0.055\textwidth}p{0.885\textwidth}@{}}
\endfirsthead\endhead
{\scriptsize\ttfamily df8cc54} & \textbf{Initial commit} \\
\addlinespace[3pt]
\end{longtable}

\section*{2026-08-08 \textnormal{\small\textcolor{srlgrey}{--- 22 commits}}}
\addcontentsline{toc}{section}{2026-08-08}
\begin{longtable}{@{}p{0.055\textwidth}p{0.885\textwidth}@{}}
\endfirsthead\endhead
{\scriptsize\ttfamily 2b4f133} & \textbf{Restructure into six packages; vendor serial; add autonomy, VR and GUI}\newline{\small Six packages with a one-way dependency arrow. srl\_teleop depends on nothing in-repo, because it is the baseline condition of every experiment and a dependency on srl\_autonomy would make the comparison circular. FILES MOVED, NOT DELETED The restructure relocated six modules and two launch files.} \\
\addlinespace[3pt]
{\scriptsize\ttfamily 844e19e} & \textbf{README: fix the patch command; record a verified clean-clone build}\newline{\small The documented invocation was wrong and a fresh clone could not be rebuilt from it. \code{git apply --directory=src} doubles the path, because the patches already begin with src/ -- every hunk failed with "src/src/ros2\_kortex/...: No such file or directory". Correct form is \code{patch -p1 --forward --fuzz=3} from the repo root. The fuzz is needed: 0003 applies at "fuzz 3 (offset -10 lines)" because ros2\_robotiq\_gripper main has moved since the patch was cut.} \\
\addlinespace[3pt]
{\scriptsize\ttfamily bcdf034} & \textbf{NEXT\_SESSION: add a CHECKPOINT so a fresh session starts oriented}\newline{\small Records what works with its evidence, what is blocked and on what, the lab order, and the standing caveat that every protocol number is IK feasibility in simulation -- nothing has been driven by a human through the master arm, and the DIRECT condition of every task depends on channels that are currently incoherent.} \\
\addlinespace[3pt]
{\scriptsize\ttfamily 4b019dc} & \textbf{T2 scenarios: four verified, and the documented geometry was wrong}\newline{\small The bimanual runner refused t2 by name because verify\_scenarios.py never generated any. scripts/verify\_t2\_scenarios.py now generates four graded scenarios and verifies the COMPLETE fill path -- pick, lift, transit, above-opening, release, retreat -- at every waypoint, while the holding arm is simultaneously solvable at the hold pose. Endpoint-only checking would have passed paths that carry the block through the wearer, which is how the T3 respec failed the first time. 4 of 4 verified, path 6/6 each.} \\
\addlinespace[3pt]
{\scriptsize\ttfamily 398f8b7} & \textbf{Audit every protocol coordinate at N=10 over the whole path; instrument T2}\newline{\small AUDIT ----- T2's coordinates passed a single IK call and failed a repeated one, so every figure in every protocol was re-verified. scripts/audit\_scenario\_reachability.py grades each pose SOLID (N/N) / MARGINAL (1..N-1) / FAIL (0) over the whole DENSIFIED path rather than its endpoints. Result: 17/17 scenarios and 17/17 protocol-table coordinates SOLID, 0 MARGINAL. Getting there changed a lot. WHICH FIGURES CHANGED} \\
\addlinespace[3pt]
{\scriptsize\ttfamily cad57ac} & \textbf{Two-person reframing: literature, hypotheses, T8/T9, measures, session order}\newline{\small The arms are worn by one person and driven by a DIFFERENT person. That is not an emphasis change -- it changes who is a participant, what is measured, and which literature the work sits in. THE GAP, CORRECTED TWICE ------------------------ The brief's reading was "Fusion established the operator/wearer architecture but did not study shared autonomy within it." Correct, and not the whole story. Two corrections that must survive into any write-up:} \\
\addlinespace[3pt]
{\scriptsize\ttfamily efa280c} & \textbf{T8/T9 analysers, and a known-answer test that hardened T9's manipulation check}\newline{\small T8's analyser reports ANTICIPATION first and its SIGN before its magnitude: negative means the wearer began moving before being asked, which is a predictive model of another person's intent built by someone with no control over the limbs and no proprioception of them. It also reports repositions per trial alongside latency, because a dyad can cut latency by the wearer shuffling constantly, which is not learning and is worse for the wearer. Both analysers refuse to invent numbers from empty columns.} \\
\addlinespace[3pt]
{\scriptsize\ttfamily 0d471b7} & \textbf{Task-band clearance probe: which link, and how far the band must move}\newline{\small Recording every task showed every chest-band scenario at 48-67 mm of clearance -- IK-valid, contact-free, and below the 120 mm real\_robot floor. That is a property of the BAND, not of any one scenario, so it gets its own measurement rather than an inference from 87 clips.} \\
\addlinespace[3pt]
{\scriptsize\ttfamily 8c692e7} & \textbf{T8/T9 protocols}\newline{\small T8 is well-posed by construction: the operator cannot reach (verified IK failure at nominal stance) and the wearer cannot manipulate (no control input at all), so a trial only succeeds through coordination. Neither is told whose move it is -- who initiates is a dependent variable and instructing it would destroy the measure. The operator holds still while the wearer repositions, because an arm moving toward a person who is themselves moving is the configuration with the least margin and the wearer cannot see it.} \\
\addlinespace[3pt]
{\scriptsize\ttfamily 1e9001a} & \textbf{setup.py: the bimanual analysers were never installed}\newline{\small EXP\_SCRIPTS globbed experiments/*/run\_*.py, one level deep. The bimanual set nests one deeper -- experiments/bimanual/t3\_coordinated\_carry/analyse\_t3.py -- so the glob matched run\_bimanual.py and NONE of the seven analysers. Every 'ros2 run srl\_experiments analyse\_t3.py' has failed since the restructure and looked like an installation problem rather than a packaging bug.} \\
\addlinespace[3pt]
{\scriptsize\ttfamily 036e4a6} & \textbf{87 verification clips, and the clearance finding resolved by measurement}\newline{\small Every task x scenario x autonomy condition driven in sim and recorded: 87 clips, 7 tasks. Zero pass through the wearer. Zero show a stationary robot. Task-phase tracking 0.0-4.7 mm RMS across all of them. recordings/verification/INDEX.md lists every clip with WHAT IT SHOULD SHOW, because an index of filenames tells a reviewer nothing -- they still have to work out what correct looks like before they can tell whether they see it.} \\
\addlinespace[3pt]
{\scriptsize\ttfamily a47dd45} & \textbf{RViz screen captures: 87 clips, objects attached and carried, 0 static}\newline{\small REAL SCREEN RECORDINGS, not TF plots. All 87 task x scenario x condition runs now carry rviz.mp4 alongside the TF clip, with the task objects visible as markers at their real dimensions and posed from the LIVE gripper pose every frame. x11grab ON WSLg's :0 RECORDS BLACK, and this is now in docs/ENGINEERING_LOG.md so it cannot be lost again. A full-screen grab with RViz plainly visible measures mean pixel value 0.0 -- XWayland windows are composited by Wayland and their pixels never reach the X root window x11grab reads.} \\
\addlinespace[3pt]
{\scriptsize\ttfamily 2747eea} & \textbf{thesis: Imperial MSc template scaffolded, evidence audited, first full draft}\newline{\small PHASE 0-1. Template extracted to thesis\_report/\_template (kept pristine) and read. thesis\_report/EVIDENCE.md maps every claim the thesis might make to the file that backs it, and lists nine gaps. Three findings from the evidence pass that shape the whole document:} \\
\addlinespace[3pt]
{\scriptsize\ttfamily a36590c} & \textbf{Four-view clips with the gripper actually operating: 87/87, fingers cycle 45/45}\newline{\small THE GRIPPER NEVER OPENED, and that was a real bug, not a rendering artefact. mock\_components boots the Robotiq driven knuckle at 0.793 rad -- fully closed on nothing -- and nothing in the recording path ever commanded it. So every previous clip showed a shut hand while the "grasp" existed only in the marker layer. The object-attachment check did not catch it and could not: the marker genuinely was following the gripper. It was following a gripper that never opened.} \\
\addlinespace[3pt]
{\scriptsize\ttfamily 3a4c160} & \textbf{Remux two oversized gripper captures (62 MB -> \textasciitilde 1 MB)}\newline{\small Their x11grab was never finalised, so the container carried the whole recording rather than the trial. GitHub warns above 50 MB and the content is identical after a stream copy.} \\
\addlinespace[3pt]
{\scriptsize\ttfamily 52727a7} & \textbf{thesis Phase 3: artefacts ingested, criteria set E'1-E'9, Chapter 1 drafted}\newline{\small INGESTED the external design record (45 files, 47 MB) into thesis\_report/\_source/artefacts/ with PROVENANCE.md: kinematics.cpp (the C++ master\_pose\_node.py cites by name), J1-J7 printed links, PotArm.step, 8 BOM PDFs, backpack and stand prints with sliced gcode, the WH148 PH1 pot CAD, the MuJoCo log dated 7 Jun 2026, and the 12 master trajectories of 27 Jul 2026. Originals untouched.} \\
\addlinespace[3pt]
{\scriptsize\ttfamily 66ac1ac} & \textbf{Real grasp sequence wired to the existing generator; thesis draft committed}\newline{\small THE GRASP WAS FAKE. The arm moved its end effector to a point, the marker attached, and it carried on: no approach vector, no wrist alignment, no pre-grasp standoff, and fingers that never moved. The marker layer was substituting for a grasp rather than following one.} \\
\addlinespace[3pt]
{\scriptsize\ttfamily abd0c74} & \textbf{Five-task consolidation, real grasp, figures, and the thesis in full draft}\newline{\small PART 1 -- FIVE TASKS. Nine families consolidated to five. The criterion for the bimanual three was strict: one arm alone must be unable to complete the task GIVEN UNLIMITED TIME. Because the reachable sets are disjoint that can only be met two ways, and both are used: a rigid object spanning the dead band (Tasks 3 and 4) and two targets requiring simultaneous satisfaction (Task 5).} \\
\addlinespace[3pt]
{\scriptsize\ttfamily 502a473} & \textbf{Five-task verification reproduced from home; single-owner gripper}\newline{\small N=10 full-path verification re-run with both arms verified at 0.0000 rad from home: 0 failures, and byte-identical to the previously committed run. That is the first reachability figure in this project to reproduce across runs rather than be re-asserted; it reproduces because the 500 mm respec moved every coupled waypoint off the feasibility boundary. Two real defects in the grasp, both found in the RECORDED TRACE rather than by inspection: 1. Two owners wrote the gripper on the same cycle.} \\
\addlinespace[3pt]
{\scriptsize\ttfamily da9950b} & \textbf{Rewrite the thesis: reframe as graceful degradation, add two chapters}\newline{\small FRAMING. Channel loss is presented as an operating condition rather than as a list of specific faults. Any body-worn master routing fourteen analog channels through rotating joints will lose channels, so the architecture is described in those terms: per-channel health verdicts computed over distinct sensor updates, validation scoped to the channels the active mode consumes, run-time channel disabling with published capability loss, and a reduced-degree-of-freedom fallback.} \\
\addlinespace[3pt]
{\scriptsize\ttfamily 92d9834} & \textbf{Write the abstract and acknowledgements}\newline{\small Abstract: opens on why the problem is specific to this class of system (the base is a person), states the two design premises, and carries the numbers that matter -- 63.8 mm cost at the last degradation rung, 6 of 14 channels healthy, 54.5 mm of commanded motion injected by the failing ones against a 272 mm master reach, 137.7 deg measured where 90 deg was intended, 14/14 faults handled, 0 of 63 shared workspace cells with a 0.25 m dead band, and the buildable mount fix at 22 of 63.} \\
\addlinespace[3pt]
{\scriptsize\ttfamily ae68df1} & \textbf{Part 1: measure whether the grasp is real. Three instrument bugs found first}\newline{\small Every figure below comes from a bare sim stack with mock hardware, and each carries how it was obtained. INSTRUMENT, VALIDATED BEFORE USE. Reading the two finger-tip frames out of TF gives the gap between LINK ORIGINS, not between pad faces. Taken raw it says the gripper never closes below 50.7 mm and that grip\_for()'s linear model is 52 mm wrong. Both are artefacts.} \\
\addlinespace[3pt]
\end{longtable}

\section*{2026-08-09 \textnormal{\small\textcolor{srlgrey}{--- 60 commits}}}
\addcontentsline{toc}{section}{2026-08-09}
\begin{longtable}{@{}p{0.055\textwidth}p{0.885\textwidth}@{}}
\endfirsthead\endhead
{\scriptsize\ttfamily e22b59c} & \textbf{Part 2: the capability ladder, selected live from whatever channels survive}\newline{\small Not a workaround for specific broken joints. Any body-worn master routing 14 analog channels through rotating joints loses channels, so the position mapping is SELECTED every cycle from whatever passes health, and the selection is published with what it cost. srl\_teleop/capability.py is pure (no ROS handles, no I/O) so the ladder is testable against known channel sets without a stack.} \\
\addlinespace[3pt]
{\scriptsize\ttfamily a870374} & \textbf{Part 3: scene fingerprinting, so calibration runs only when the bed changed}\newline{\small srl\_perception/scene\_fingerprint.py is the matcher and is PURE -- no ROS handles, no camera, no I/O beyond an explicit path -- because that is the part testable against CONSTRUCTED ground truth. Two points 30 mm apart really are 30 mm apart, which is the one kind of synthetic data the standing rule permits, and it is why these figures can be trusted where a rendered-image test of the same logic could not.} \\
\addlinespace[3pt]
{\scriptsize\ttfamily 9b7862b} & \textbf{Part 4: verify all four mode paths. The VR path was broken; fixed}\newline{\small METHOD. Every hop confirmed from the LIVE ROS graph: a hop counts as CONNECTED only when the topic has both a publisher and a subscriber. Reading an import and concluding the path exists is the mistake this script exists to avoid, and it is how the defect below survived. THE VR COMMAND PATH DID NOT REACH THE ARMS. quest\_vendor\_bridge -- the transport the project adopted -- published /vr\_pose\_<side> and /vr\_gripper\_<side>. Nothing in the workspace subscribes to either.} \\
\addlinespace[3pt]
{\scriptsize\ttfamily c4332b3} & \textbf{Part 5: free-form language. Opens the noun, not the verb}\newline{\small THE TENSION WITH THE RECORDED DECISION, RESOLVED RATHER THAN IGNORED. This project chose a deterministic grammar over a language model, and the reason holds: a grammar that does not recognise an utterance returns unparsed and the robot asks again, whereas a model asked for structured output can return a different, plausible verb that nothing downstream can distinguish from a correct parse. Free-form input does not change that. It changes WHERE the openness lives. VERBS stay closed.} \\
\addlinespace[3pt]
{\scriptsize\ttfamily 485fad5} & \textbf{Part 6: Qt operations GUI with RViz embedded, Helvetica, everything indicated}\newline{\small RVIZ EMBEDDING: INVESTIGATED, AND ROUTE 1 VERIFIED FROM PIXELS. 1. X11 reparenting. QX11EmbedContainer is gone in Qt5; the equivalent is QWindow.fromWinId() wrapped by QWidget.createWindowContainer(). WORKS. Proven under Xvfb by screenshot: the host label renders above and the complete RViz UI renders inside the container at 31 fps.} \\
\addlinespace[3pt]
{\scriptsize\ttfamily 5b5abe4} & \textbf{Part 7: one-command Quest connection, supervised, plus a drawn control diagram}\newline{\small scripts/vr\_connect.sh finds adb, waits for the headset, opens the reverse tunnel, starts the bridge, and then SUPERVISES the tunnel for as long as it runs. WHY SUPERVISED RATHER THAN OPENED ONCE. \code{adb reverse} dies with the USB connection. Opened once at start-up, a nudged cable leaves the bridge running and the headset silently unable to reach it, which from the operator's side is indistinguishable from a frozen application.} \\
\addlinespace[3pt]
{\scriptsize\ttfamily 0264277} & \textbf{Part 8: six operator aids evaluated, one built, one honestly unfinished}\newline{\small EVALUATION of all six against what this rig has actually measured, in docs/system/08\_operator\_aids.md. Four rejected for reasons specific to this hardware rather than general ones:} \\
\addlinespace[3pt]
{\scriptsize\ttfamily a5ae397} & \textbf{Part 9: audit every reported number; re-recording not done}\newline{\small AUDIT, in docs/system/09\_results\_audit.md. For each figure in Parts 1-8: how it was measured, whether the instrument was validated against an independently known answer, and a verdict. VALIDATED 14 SOUND 15 UNVERIFIED 4 STRUCTURAL 1 The reason this exists: across this project the measuring tool has been the fault EIGHT times, three of them in this run of work. The base rate is high enough that an unaudited number is not evidence.} \\
\addlinespace[3pt]
{\scriptsize\ttfamily a63cebb} & \textbf{Resolve the boundary-node contradiction: it was the harness, twice}\newline{\small VERIFIED: the warning arrives with 88 mm of travel still remaining, after 334 mm of travel, identical in 3 of 4 repeats. THE CONTRADICTION WAS MINE, NOT THE NODE'S, IN TWO SEPARATE PLACES. First, the drive test commanded an IDENTITY orientation. Measured: identity is unreachable at 0 of 6 points along that line while the home orientation reaches 6 of 6. "Wall at 0 mm" was the correct answer to the question actually being asked.} \\
\addlinespace[3pt]
{\scriptsize\ttfamily 48dcdfa} & \textbf{Task 2 loses DIRECT and VR as a finding; cut the boundary node; adb; methods}\newline{\small 1. TASK 2 HAS NO DIRECT OR VR CONDITION, and the objects were NOT made easier. Teleoperated grasping is not achievable with this master: a top-down grasp needs 169.7 deg (left) / 164.6 deg (right) from the anchor orientation\_mode "fixed" pins the wrist to, and nothing measures the wrist to command it. VR inherits the same limit because it drives the same follower through the same pinned path.} \\
\addlinespace[3pt]
{\scriptsize\ttfamily 2d3f9dc} & \textbf{Item 1 part one: five-task sweep built, verifier validated, 4 clips recorded}\newline{\small THE SWEEP HAS NOT BEEN RUN. 4 of 38 clips exist. Stopping at the item boundary with the infrastructure finished and the verifier cleared, rather than leaving a half-recorded sweep. BUILT * scripts/final5\_runs.py drives the recorder from the CURRENT five-task spec: 14 runs, 38 clips, 266 files at 7 angles, under f1..f5 so nothing collides with the 87 stale nine-task clips. * record\_rviz.py --final5 --resume.} \\
\addlinespace[3pt]
{\scriptsize\ttfamily 9cd2728} & \textbf{Item 1 complete: all 38 clips re-recorded, 39/39 pass a validated verifier}\newline{\small Five-task geometry (500 mm tray, 540 mm sling) with the single-owner gripper fix. 7 angles per clip, 266 files, 132 MB, largest file 5.3 MB. INDEX.md carries the Task 2 caveat at the top and the caveat is burned into every f2 overlay in red. A REAL STALE-GEOMETRY BUG, which is what the re-record existed to find. record\_rviz.SLING\_L was still the retired 350 mm while the spec is 540 mm.} \\
\addlinespace[3pt]
{\scriptsize\ttfamily 54ec2bc} & \textbf{Track every angle for the five-task clips}\newline{\small The brief asked for one file per angle. The ignore rule that kept only the quad tile was written for the retired nine-task sweep, whose rushes were interrupted mid-encode and could not be stream-copied back to a sane size. The five-task clips are recompressed at crf 30: 266 files, 132 MB, largest 5.3 MB, which is small enough to track. A single angle cannot answer every question.} \\
\addlinespace[3pt]
{\scriptsize\ttfamily e5d81a4} & \textbf{Item 2: wrist camera relay. Logic measured; GUI widget NOT visually confirmed}\newline{\small srl\_teleop/camera\_relay.py holds the staleness rule, shared by the GUI widget and the VR publisher so the two cannot drift apart. MEASURED against a synthetic 15 Hz camera on both arms: measured rate 14.2 Hz, reported from ARRIVALS not from configuration transport latency median 2.7 / 3.1 ms, p95 7.6 ms (header stamp to receipt) after the source stops: live -> NO SIGNAL, painting disabled IT SUBSCRIBES, IT NEVER OPENS A DEVICE.} \\
\addlinespace[3pt]
{\scriptsize\ttfamily d4187c9} & \textbf{Item 3: the instrument-validation chapter, the thesis's methods contribution}\newline{\small CHECKED WHAT EXISTED BEFORE WRITING. The master arm chapter (2799 words), the MuJoCo-to-ROS development path (1631), and all four control-mode flowcharts were already written in the earlier rewrite. Duplicating them would have been the easy way to look productive. The genuinely missing piece was the methods contribution, now Chapter 6: "When the instrument is the fault".} \\
\addlinespace[3pt]
{\scriptsize\ttfamily 1394a05} & \textbf{Item 4: the lab list, ordered so each step decides whether the next is worth it}\newline{\small docs/NEXT\_SESSION.md now opens with the hardware session. Everything above the line needs the lab; nothing below it does. 0. Exactly ONE stack, and clear stale /dev/shm segments with the stack stopped. Those segments silently kill newly launched nodes -- empty logs, no processes, no error -- which cost hours on 08-09. Kill by explicit PID: pkill -f <pattern> matches the shell running it and killed the working shell three times in one session. 1. Repair LEFT j2 and j4, and nothing else.} \\
\addlinespace[3pt]
{\scriptsize\ttfamily 1288444} & \textbf{Confirm the GUI camera panel renders, from pixels, in both states}\newline{\small Both states screenshotted and committed to docs/img/: live both feeds painting, labelled LEFT / RIGHT, captions "live 13.8 Hz 320x240" in green source stopped both panels black, "NO SIGNAL" in red, and "last frame 5.9 s ago" That second shot is the one that matters. It is the visible-degradation claim made good: the panel does not dim or freeze the last frame, it stops painting and says why, because a frozen picture of a workspace cannot be told from a live picture of a workspace that is not moving.} \\
\addlinespace[3pt]
{\scriptsize\ttfamily 9d72c6d} & \textbf{f4 clip binaries from the post-sling-fix pass}\newline{\small Six left-behind re-encodes of the S1\_short\_lift/direct angles. Same clips, committed so the working tree is clean before consolidating branches rather than discarded by a checkout.} \\
\addlinespace[3pt]
{\scriptsize\ttfamily 75d3800} & \textbf{Final f4 clip binaries, left changing by orphaned capture processes}\newline{\small The sweep left ffmpeg x11grab processes running against the f4 clips, so those files kept changing under git and blocked a branch switch. Killing them made the tree stable. Same orphan class as Part 4, where 183 stale /dev/shm segments from an earlier sweep silently killed newly launched nodes.} \\
\addlinespace[3pt]
{\scriptsize\ttfamily 15cbbe4} & \textbf{Job A: restructure recordings by control mode, and five silent breakages}\newline{\small recordings/verification/<mode>/<task>/<scenario>/<condition>/ Control mode and autonomy condition are independent axes and the old tree conflated them: "direct" under the mannequin and "direct" under VR shared one folder name with nothing in the path to tell them apart. Six numbered mode directories (01\_master\_teleop .. 06\_full\_autonomy) so they sort by capability rather than alphabetically, which would put full\_autonomy between direct and shared. THE FINDING THAT SHAPED THE MIGRATION.} \\
\addlinespace[3pt]
{\scriptsize\ttfamily c29e9be} & \textbf{Job B: workspace audit. Report only, nothing fixed}\newline{\small scripts/audit\_workspace.py, repeatable, 232 python files. Every check exists because reading did not catch that bug in this project. CLEAN. The five recurring bug classes are essentially gone: wall-clock intervals 0, block()-without-clear 0, second-stack starts 0, unresolved entry points 0, missing referenced paths 0, and rejected task args 0 -- that last being the class that once made five GUI buttons do nothing.} \\
\addlinespace[3pt]
{\scriptsize\ttfamily adaaef8} & \textbf{Fix the four vacuous checks; wire the post-hoc stage holm was written for}\newline{\small THE FOUR VACUOUS CHECKS: 0 remaining. verify\_final5 now RAISES if densify returns no waypoints instead of reporting "TOTAL FAILURES: 0" on a path it never checked -- that check sits under the result this project leans on hardest. Same guard on the place path, on verify\_scenarios's centres dict (all() over an empty dict is True), and on verify\_autonomy's expected-words list (an empty expectation passes anything the robot says, including nothing). HOLM: THE GAP WAS NOT A MISSING CALL. holm() had nothing to correct.} \\
\addlinespace[3pt]
{\scriptsize\ttfamily 64fe28a} & \textbf{Job C: /mnt/c works, first CAD measurement, no library needed}\newline{\small C1: /mnt/c RECOVERED ON ITS OWN. It reads fine and so does the project directory; the mount entry is the same 9p one that was returning EIO earlier, so that was transient. No Windows-side command needed. Prior work confirmed present: mujoco\_menagerie, MUJOCO\_LOG.TXT, singlekinova.py, srl\_teleop.py, testjoints.py, basic\_control/, 3dprint/, BOM/. C2, PARTLY. Recommendation: NO CAD LIBRARY for this part of the question.} \\
\addlinespace[3pt]
{\scriptsize\ttfamily 59908d0} & \textbf{C2 finished: the CAD validates the FK link lengths; the mount angle does not}\newline{\small LIBRARY. pythonocc-core is NOT on PyPI -- it ships through conda and this workspace has none. cadquery bundles OCP, a binding to the same OpenCascade kernel, and installs with pip: identical kernel, different packaging. For overall bounding extents no kernel is needed at all, since STEP is a text format, and measure\_cad\_step.py does that with a regex.} \\
\addlinespace[3pt]
{\scriptsize\ttfamily a7e5204} & \textbf{C3: prior work recovered and read; two corrections to the development chapter}\newline{\small The development chapter was written from reconstruction. The actual programs have now been recovered and archived at thesis\_report/\_source/prior\_work/, and file dates give the calendar directly: 7 June to 27 July 2026. The chapter cites the source instead of supposing. WHAT THE PRIOR WORK ACTUALLY WAS. Five programs import MuJoCo and NONE contains an inverse-kinematics call of any kind -- searching the whole set for "inverse" or "jacobian" returns nothing.} \\
\addlinespace[3pt]
{\scriptsize\ttfamily 12fa17f} & \textbf{C4 part one: the first real RViz figure, and the participant-Results skeleton}\newline{\small THE REPORT HAD NO RVIZ VISUALS AT ALL. It now has one, captured from the running system on Xvfb rather than drawn, cropped to the 3-D view and annotated: left arm at +x, right arm at -x, gripper, backpack mount, and the wearer marked as being IN the collision model. It sits in the workspace chapter, where it SHOWS the disjoint-workspace result instead of only asserting it: both hands finish on the same side of the body, 1.46 m apart.} \\
\addlinespace[3pt]
{\scriptsize\ttfamily 6789d2f} & \textbf{C5: state that the CAD is not the kinematic source, with each claim checked}\newline{\small A new section in the master arm chapter, so nobody later mistakes the mechanical drawings for the model the software uses. Three statements, each verified rather than asserted: 1. THE CAD CARRIES NO ARTICULATED JOINTS. Both STEP files declare AP214 (AUTOMOTIVE\_DESIGN), a schema that IS capable of carrying kinematic structure, so the absence is meaningful rather than a format limitation.} \\
\addlinespace[3pt]
{\scriptsize\ttfamily 3c5df49} & \textbf{C4 part two: CAD renders as vector; and a correction to my own C4 claim}\newline{\small CORRECTION FIRST. The previous C4 commit said the RViz capture recipe was "one command per view". It is not. shoot.sh accepts yaw, pitch and distance arguments and IGNORES all of them -- the camera pose comes from the .rviz config, so every invocation renders the same view. Two "different" captures came out byte-identical, which is how it was caught. The fix is a per-view .rviz config with the camera pose written into it, which record\_rviz.py:ensure\_display() already does and shoot.sh should reuse.} \\
\addlinespace[3pt]
{\scriptsize\ttfamily 77f3394} & \textbf{Job A: audit clean, but the pot repair does not take effect yet}\newline{\small THE HEADLINE, and it changes what Job B can claim. degraded\_mode.decide() reads a STORED BASELINE FILE, not live hardware. channels\_20260806.json still says 6/14 coherent, so master\_pose\_node engages degraded mode at startup and freezes eight now-working channels, silently. stale baseline 6/14 coherent -> degraded engages: True repaired arm 14/14 coherent -> degraded engages: False Both directions verified.} \\
\addlinespace[3pt]
{\scriptsize\ttfamily cafaa5c} & \textbf{Job B.1: wrist orientation capability after the pot repair}\newline{\small Master can now express 180 deg of tip reorientation against the 169.7/164.6 needed for a top-down grasp, so the repair restores the input side. The robot is now the binding constraint: over a 27-pose grid, fixed 74.1/85.2\%, tilt 73.5/81.2\%. Tilt costs 0.6-4.0 pp against fixed but neither clears the 90\% bar, so tilt is not shipped.} \\
\addlinespace[3pt]
{\scriptsize\ttfamily 7f17263} & \textbf{Job B.2: capability delta from the pot repair, and what is lab-blocked}\newline{\small Left arm's commandable set goes from a surface (radial extent 0.000 m, radius not commandable at all) back to a shell (0.139 m). Right arm unchanged, its frozen channels being rolls that carry no radial information. Job B's azimuth, position-ladder and smoothing questions are fits to recorded frames and every recording predates the repair, so they are named as lab-blocked with the exact capture that unblocks them rather than re-run on pre-repair data.} \\
\addlinespace[3pt]
{\scriptsize\ttfamily e12efe0} & \textbf{Job C: fix grasp poses driving the fingers through the object}\newline{\small candidates() returned the bare centroid while the fingertips sit 0.1118 m along the tool +z and the grasp quaternion aims it down, so every object in the catalogue was penetrated by 51.8-105.8 mm. grasp\_offset() now raises the tool by half-height plus pad. Root cause of the missed check: nothing in the autonomy path publishes the perceived object as a CollisionObject, so avoid\_collisions=True ran against a world without it and passed 100\%.} \\
\addlinespace[3pt]
{\scriptsize\ttfamily ca68cc9} & \textbf{Job D: mode command paths verified from the live graph}\newline{\small Modes 4, 5 and 6 cannot command the arm. handover\_arbiter publishes /autonomy/assist\_pose\_<arm> and nothing anywhere subscribes to it; and autonomy\_executive's complete publisher list is speech, stage, decision and estop, with no pose, trajectory or action client. Every stage before the last is real; the last hop was never connected.} \\
\addlinespace[3pt]
{\scriptsize\ttfamily ba22676} & \textbf{Job E: publish the calibration sweep to RViz as live markers}\newline{\small scene\_markers.py renders each fingerprint verdict as a coloured box plus a label carrying name, verdict, confidence and displacement, on /scene/markers. SAME is deliberately desaturated so the verdicts needing attention carry the colour; a move is drawn as an arrow from its stored pose; a dropped object is held for drop\_hold\_s at its stored pose before being released, because an instant removal is indistinguishable from an object that was never there.} \\
\addlinespace[3pt]
{\scriptsize\ttfamily a91f8bc} & \textbf{Job F: precision/speed dial with an enforced safety invariant}\newline{\small One dial in [0,1] moves scale, smoothing, velocity cap and step limit together; scale is geometric so the low end gets the resolution. SAFETY\_PARAMS names the nine things it must never touch and assert\_safety\_unconditional raises on any of them, checked across all 101 dial positions so a future key added to RANGE fails the test rather than shipping. rebase\_anchor keeps the command continuous to 1e-12 through a dial change, including far from the engage point where the uncorrected jump is worst.} \\
\addlinespace[3pt]
{\scriptsize\ttfamily 6bab268} & \textbf{Blocker 1: connect the autonomy command path through the teleop safety stack}\newline{\small ik\_follower\_node now subscribes to /autonomy/assist\_pose\_<arm>, which had zero subscribers anywhere -- the reason modes 4, 5 and 6 never drove the arm. It enters at request\_ik, the same door teleop uses, so collision-aware IK, the clearance floor, the guard, the flip reject and the e-stop are the same code with no per-mode exemption and no private route to the controller. autonomy\_executive gained command()/release() and a 20 Hz pump; its previous complete publisher list was speech, stage, decision and estop.} \\
\addlinespace[3pt]
{\scriptsize\ttfamily c6c31fe} & \textbf{docs: standing work brief as the session resume point}\newline{\small docs/WORK\_BRIEF.md carries the multi-part brief verbatim with a STATUS line per part, so a fresh session resumes from "Read docs/ENGINEERING_LOG.md, NEXT\_SESSION.md and docs/WORK\_BRIEF.md" with nothing re-pasted. Parts 1-4 are the user's text verbatim. Parts 5-8 (Tasks A/B/C and the dual-view GUI) are recorded as TEXT NOT SUPPLIED rather than reconstructed: only the one-line forward reference naming them exists, so any specification written here would be invented and then graded against itself.} \\
\addlinespace[3pt]
{\scriptsize\ttfamily 6736ee7} & \textbf{Part 1: autonomy moves the arm -- root cause was a missing import}\newline{\small ANSWER: YES. 1092 trajectories at a reachable target and 0.0700 m / 0.1204 m of real tf2 EE displacement over two legs, residual 0.0000 m on both. THE BUG. ik\_follower\_node.py had no \code{import time} while the autonomy path calls time.monotonic() at lines 713 and 733. A missing name is evaluated only when the path runs, so the node imported, launched, logged "Tracking enabled", and died on the first /autonomy/assist\_pose\_<arm> message with NameError. The followers launch without respawn, so it stayed dead.} \\
\addlinespace[3pt]
{\scriptsize\ttfamily cf9e4cd} & \textbf{docs: WORK\_BRIEF tracks the 7-part brief; Part 1 marked DONE}\newline{\small Supersedes the earlier 4-part version. Parts 1-7 verbatim with STATUS lines and commit hashes, START HERE pointing at Part 2. Part 1 recorded DONE at 6736ee7 with the Part 1 findings that feed directly into Part 2's diagnosis (bug class 5 instance found, SHM-at-startup confirmed necessary including the sem.fastrtps\_* glob miss, follower respawn left as a deliberate open decision).} \\
\addlinespace[3pt]
{\scriptsize\ttfamily 9e0dce4} & \textbf{Part 2: full diagnosis -- one bug wore three costumes, two instruments lied}\newline{\small DIAGNOSIS. Part 1's missing import killed the follower, and a dead node publishes nothing, so three separately-filed faults were one corpse seen from three angles: "autonomy does not move the arm", "/ik\_status\_<arm> publishes an empty data array" (Part 2D.3), and "probe\_dial\_safety recorded zero clearance samples". Measured after the fix, /ik\_status carries 11 fields with clearance at index 5 exactly as documented. D3 is closed by the Part 1 fix. A. Five bug classes. A1 one instance (fixed).} \\
\addlinespace[3pt]
{\scriptsize\ttfamily b6e5a8b} & \textbf{docs: Part 2 marked DONE, START HERE points at Part 3}\newline{\small Also records the live-stack state and the parked left arm, so the next session re-homes before any measurement that assumes the home pose.} \\
\addlinespace[3pt]
{\scriptsize\ttfamily c0ff406} & \textbf{Part 3.1: language sweep -- 7 MISUNDERSTOOD found, all seven closed}\newline{\small 40 phrasings, 13 categories, written against the parser's STRUCTURE rather than alongside it. The prior self-test (measure\_freeform\_language.py) reported 0 MISUNDERSTOOD over 30 phrases, but every one of those was written with the grammar it tests -- bug class 5, a test constructing the environment where the bug cannot occur.} \\
\addlinespace[3pt]
{\scriptsize\ttfamily 759dfe6} & \textbf{Part 3.2: the gripper no longer inherits a leftover as its open reference}\newline{\small The 2F-85 HOLDS POSITION when its commanding process goes away. So a node killed mid-grasp left the fingers part-closed and the next start found them there with nothing able to say whether that was a grip or a leftover. Two opposite mistakes follow and no default avoids both: open unconditionally -> drops whatever the hand was holding never open -> the travel range is wrong for the whole session THE MARKER IS WRITTEN WHEN THE GRIP IS TAKEN, NOT WHEN THE PROCESS EXITS. This is the decision the whole thing rests on.} \\
\addlinespace[3pt]
{\scriptsize\ttfamily c51ee0f} & \textbf{Part 3.3: the pot repair could not take effect, and the warning had never fired}\newline{\small TWO DEFECTS, AND NEITHER WAS "THE WARNING IS TOO QUIET". 1. THE RUNTIME READ A DIFFERENT FILE FROM THE ONE THE LAB WRITES. \code{check\_channels.sh} saves the new reference as channels\_<today>.json and diffs against \code{ls -t channels\_*.json | head -1}. \code{degraded\_mode.default\_baseline\_path()} returned the literal string "channels\_20260806.json".} \\
\addlinespace[3pt]
{\scriptsize\ttfamily d5bbec8} & \textbf{Part 3.4: five tasks -> three, verified N=10 over the full path, 116 min}\newline{\small TASKS A, B AND C, in src/srl\_experiments/experiments/abc/. A POSITIONING uncoupled; attention contrast; Fitts; warm-up B COORDINATED CARRY physical coupling, rigid AND compliant C DUAL PURSUIT simultaneity across the disjoint sets; b\_cross The binding constraint is the WEARER, not the science: >17 kg, 12 min continuous and 48 min total pack-on, and a two-hour session. Two modes x five tasks was costed at 167 minutes and does not fit.} \\
\addlinespace[3pt]
{\scriptsize\ttfamily 6b9d78a} & \textbf{docs: Part 3 marked DONE, START HERE points at Part 4}\newline{\small Findings write-up in NEXT\_SESSION.md. Records the three items that turned out worse than the brief described: the gripper's startup verdict was re-decided mid-motion, the degraded-mode warning had never fired at all, and the pot repair could not reach the runtime because the lab and the code disagreed on which baseline file is current.} \\
\addlinespace[3pt]
{\scriptsize\ttfamily 95adcf4} & \textbf{Part 4: commanded vs actual, with the divergence named -- and two panels proven not to work here}\newline{\small WHAT SHIPPED. scripts/srl\_gui.py rebuilt around the question the brief asks: what is the gap between what was commanded and what the arms actually did?} \\
\addlinespace[3pt]
{\scriptsize\ttfamily dd5ae82} & \textbf{docs: Part 4 marked DONE, START HERE points at Part 5}\newline{\small Records the two findings that changed the design: X11 reparenting embeds but does not clip without a window manager, so two side-by-side RViz panels are not available on this host; and /real/tf -- which the brief names -- does not exist, the real arm being prefixed into the shared /tf.} \\
\addlinespace[3pt]
{\scriptsize\ttfamily 93fcd5c} & \textbf{Part 5 preliminaries: capture frame time properly, and correct /real/tf}\newline{\small FRAME TIME IS NOW MEASURED, NOT READ OFF A SCREENSHOT. Part 4 could only report frame time by reading the status bar in a captured image, because the harness got an EMPTY string back from the GUI's stdout. Chasing that turned up three faults, two of them in the measuring tool and one of them serious: 1. \code{pgrep -f "Xvfb :99"} MATCHED ITS OWN SHELL COMMAND LINE, so ensure\_xvfb() believed the server was up and never started it. Every GUI run then died with "could not connect to display :99".} \\
\addlinespace[3pt]
{\scriptsize\ttfamily 3597e0a} & \textbf{Part 5: the clip verifier now refuses to report until it catches a broken clip}\newline{\small THE NEGATIVE CONTROL. Before any verdict is printed, the verifier is shown seven CONSTRUCTED clips whose answers are known and must get every one right. It refuses to report at all otherwise, because "0 clips failed" and "nothing was actually checked" print identically -- and this file has produced the second while reporting the first.} \\
\addlinespace[3pt]
{\scriptsize\ttfamily 93ce00c} & \textbf{Part 5 (partial): tasks A/B/C are runnable through a mode's own command path}\newline{\small WHAT NOW WORKS, AND IT IS THE THING THAT WAS BLOCKING THE WHOLE PART. Job A found that NO clip in this repository had ever been recorded through a control mode: every one came from record\_rviz.py, which calls /compute\_ik directly and never publishes /master\_arm\_pose\_*, so no follower, clutch, anchor, orientation lock or safety guard was in the path. The mode axis of the recordings tree was empty because nothing could fill it. \code{experiments/abc/run\_abc.py} fills it.} \\
\addlinespace[3pt]
{\scriptsize\ttfamily 39f8869} & \textbf{docs: Part 5 partly done, resume point is the capture sweep}\newline{\small Records what is proven (the A/B/C runner drives the arms through a mode's own command path; the verifier refuses to report until it catches a broken clip) and what is not (the 7-angle sweep, mode 01 suppression, mode 02 untested).} \\
\addlinespace[3pt]
{\scriptsize\ttfamily bc3b890} & \textbf{Part 5: both blockers resolved -- mode 01 diagnosed, VR driven end to end}\newline{\small BLOCKER 1 -- autonomy\_has\_control is REFUTED. It was never suppression. The decisive evidence was /ik\_status, whose layout is [success, fail, direct, slewed, rejected, clearance, blocks, redundancy]: success fail redundancy retries before 0 8 48 after a mode-01 run 0 15 (+7) 90 (+42) after a mode-04 run 8 (+8) 15 90 Mode 01 made SEVEN IK calls and FAILED ALL SEVEN, burning 42 redundancy retries (7 x redundancy\_samples). Every BlockMonitor blocker stayed clear throughout.} \\
\addlinespace[3pt]
{\scriptsize\ttfamily b071909} & \textbf{docs: all four modes drive the arms; the sweep needs mode isolation}\newline{\small Records the measured travel per mode and the precondition the sweep must honour: after a VR run, modes 01/04/06 reported 0.0000 m, and killing vr\_pose\_mapper restored them to 0.2000 m. The mapper stays engaged and keeps publishing on /master\_arm\_pose\_<arm> -- one source at a time, at the process level.} \\
\addlinespace[3pt]
{\scriptsize\ttfamily 131f2af} & \textbf{Part 5: the 7-angle sweep, recorded through the GUI, with isolation enforced}\newline{\small TWELVE CLIPS, 3 tasks x 4 modes, 8 files each (7 angles + the quad tile), every one started by pressing the GUI's own button (Gui.on\_launch on the real manifest). These are the FIRST clips in this repository recorded through a control mode: Job A found every earlier clip came from a recorder calling /compute\_ik directly, with no follower, clutch, anchor or safety guard in the path, which is why the mode axis of the tree was empty. THE ISOLATION RULE IS IN THE CODE AND IT FIRED ON THE FIRST RUN.} \\
\addlinespace[3pt]
{\scriptsize\ttfamily a843c67} & \textbf{docs: Part 5 DONE -- 12 clips through the GUI, isolation enforced}\newline{\small Records the two findings that shaped it: master:=false is required for a scripted-operator stack (master\_pose\_node is otherwise a second publisher and the isolation check refuses to record), and the A/B/C clips carry no task objects so the object-in-frame criterion is not claimed for them.} \\
\addlinespace[3pt]
{\scriptsize\ttfamily b4da854} & \textbf{chore: ignore RViz per-view log files written beside the clips}\newline{\small They are transient stderr from the capture displays, not evidence, and they were showing up as modified on every sweep.} \\
\addlinespace[3pt]
{\scriptsize\ttfamily e2dc845} & \textbf{Clips: 06\_full\_autonomy recorded -- 3 tasks, 7 angles + quad, captions burnt in}\newline{\small FIRST MODE IN THE ORDER, and pushed on its own so a dead session still leaves it on the remote. Full autonomy needs no operator input, which is why it goes first. A pick and place LEFT descends 0.10 m, closes, lifts 0.14 m, carries to the bin, releases; RIGHT holds still B hold and place LEFT holds the work at x=+0.25; RIGHT brings the part down 0.15 m and places it C multimeter LEFT presents and holds steady; RIGHT probes down 0.11 m, holds contact, retracts THE CLIP TASKS ARE BUILT FROM VERIFIED COORDINATES ONLY. [...]} \\
\addlinespace[3pt]
{\scriptsize\ttfamily 6799649} & \textbf{Clips: 01\_master\_teleop recorded -- 3 tasks, 7 angles + quad}\newline{\small Scripted master input, entering at /master\_arm\_pose\_<arm> -- the same topic master\_pose\_node publishes and the same follower callback live teleop uses, so the pose deadband, collision-aware IK, redundancy re-seed, clearance floor, step guard and every blocker are all in the path. The stack runs with master:=false, which is what makes the isolation check pass: master\_pose\_node is otherwise a second publisher on that topic and the sweep correctly refuses to record.} \\
\addlinespace[3pt]
{\scriptsize\ttfamily 515264b} & \textbf{Clips: 03\_shared\_autonomy recorded -- 3 tasks, 7 angles + quad}\newline{\small Autonomy commanding through /autonomy/assist\_pose\_<arm>, the arbiter's topic, entering request\_ik at the same door teleop uses. Isolation verified before every clip: 0 idle publishers, as predicted. 3/3, captions burnt in.} \\
\addlinespace[3pt]
\end{longtable}

\section*{2026-08-10 \textnormal{\small\textcolor{srlgrey}{--- 52 commits}}}
\addcontentsline{toc}{section}{2026-08-10}
\begin{longtable}{@{}p{0.055\textwidth}p{0.885\textwidth}@{}}
\endfirsthead\endhead
{\scriptsize\ttfamily b4913dc} & \textbf{Clips: 02\_vr\_teleop recorded -- 3 tasks, 7 angles + quad}\newline{\small Scripted Quest input on /vr/controller\_pose\_<side>, through vr\_pose\_mapper, which is the only publisher on /master\_arm\_pose\_* -- and the isolation check predicts exactly that (1 idle publisher, not 0 as for the other modes) and verified it before every clip. The VR path is genuinely relative: the mapper anchors on the clutch engage and commands motion since.} \\
\addlinespace[3pt]
{\scriptsize\ttfamily fa1ec28} & \textbf{Clips: 04\_vr\_shared recorded -- all five modes now complete}\newline{\small VR transport UP and autonomy owning the pose: vr\_pose\_mapper runs, but the command enters at /autonomy/assist\_pose\_<arm>, so the isolation check predicts 0 idle publishers on the follower's input (not 1 as for pure VR) and verified it before every clip. That the two VR modes predict DIFFERENT counts, and both were met, is the check doing real work rather than rubber-stamping. 3/3. Fifteen clips across five modes, 8 files each.} \\
\addlinespace[3pt]
{\scriptsize\ttfamily b7fda18} & \textbf{Clips: the language sweep on video -- 40 phrasings, every outcome captioned}\newline{\small 2m21s, one card per utterance, grouped by category, colour-coded by outcome. Every card is the SHIPPED grammar's real answer: voice\_intent.parse then resolve\_target against a constructed scene, the same path scripts/sweep\_language\_vision.py runs. The robot's spoken reply on each card is read off that call, not written by hand. CORRECT 10 ASKED 6 REFUSED 24 MISUNDERSTOOD 0 30 of 40 utterances end with the arm untouched, and the video says so: ASKED and REFUSED are SUCCESSES.} \\
\addlinespace[3pt]
{\scriptsize\ttfamily f1c6fdf} & \textbf{INDEX.md: generated from disk, the progress file and the verifier}\newline{\small 24 clips indexed, 24 verified. The five most important at the top, with how to open them from Windows (the \textbackslash \textbackslash wsl\$ path, or explorer.exe from a WSL shell). GENERATED, NEVER HAND-WRITTEN. This file was once overwritten by a script that dropped 38 clips and discarded a caveat with them, so every row is read from the clips themselves, the sweep's progress file and the verifier's own output -- it cannot claim a clip that is not there or a pass that was not measured.} \\
\addlinespace[3pt]
{\scriptsize\ttfamily 7de5beb} & \textbf{Graphs: language outcomes and per-mode EE travel, from the sweep's own data}\newline{\small Both read their numbers from measurements rather than restating them: language\_outcomes from recordings/baselines/language\_vision\_sweep.json, and mode\_travel from the tf2 path integrals the recording sweep produced. MISUNDERSTOOD is plotted even at ZERO and annotated as the dangerous category. A chart that omits an empty bar invites the reader to forget it was measured -- and zero is the entire result, so it has to read as a measured zero and not as an absence.} \\
\addlinespace[3pt]
{\scriptsize\ttfamily 68732d6} & \textbf{docs: clips recorded -- 15 across five modes, plus the language sweep}\newline{\small Records what is on the remote and names the one outstanding item (GUI tutorial clips) with the pieces that already exist for it.} \\
\addlinespace[3pt]
{\scriptsize\ttfamily daa3df5} & \textbf{Clips: objects in the scene, grasped for real -- 06\_full\_autonomy re-recorded}\newline{\small The first sweep drove each mode's command path with NO OBJECTS PUBLISHED, so fifteen clips evidenced that a mode commands the robot and nothing about anything being picked up -- and task C had no multimeter in it at all. The pixel check could only ever ask "did the arm move". NOW: table, teal container, orange block, green circuit box, tray, and a multimeter dummy body, at the verified task coordinates.} \\
\addlinespace[3pt]
{\scriptsize\ttfamily b2e282f} & \textbf{Clips: 01\_master\_teleop re-recorded with objects grasped}\newline{\small A block GRASPED->RELEASED carried 0.412 m; B part GRASPED->RELEASED carried 0.125 m; C multimeter GRASPED and held. Same grasp evidence as full autonomy, through the master command path.} \\
\addlinespace[3pt]
{\scriptsize\ttfamily f4bef7e} & \textbf{Clips: 03\_shared\_autonomy re-recorded with objects grasped}\newline{\small A 0.412 m, B 0.125 m, both GRASPED->RELEASED; C GRASPED and held.} \\
\addlinespace[3pt]
{\scriptsize\ttfamily b0a3d94} & \textbf{Clips: 02\_vr\_teleop re-recorded with objects grasped}\newline{\small A carried 0.206 m, B 0.062 m -- almost exactly HALF the other modes' 0.412 and 0.125. That is vr\_pose\_mapper's x0.5 scale showing up in the object's own travel, which is a stronger confirmation the VR path is real than the arm motion alone: the block is being carried by the gripper, so it inherits the scale. C GRASPED and held, carried 0.000 m -- the present-and-hold task.} \\
\addlinespace[3pt]
{\scriptsize\ttfamily 3feecfc} & \textbf{Clips: 04\_vr\_shared re-recorded -- all 15 now have objects grasped}\newline{\small A 0.412 m, B 0.125 m GRASPED->RELEASED; C GRASPED and held. Full distance, not VR's half, because autonomy owns the pose in this mode -- the scale difference tracks which path is commanding, exactly as it should.} \\
\addlinespace[3pt]
{\scriptsize\ttfamily 4f0ef88} & \textbf{GUI tutorials: six workflows, narrated, clicks are real}\newline{\small launch\_a\_mode run\_a\_task read\_the\_divergence precision\_speed\_dial estop\_trigger\_clear switching\_modes Every step calls the widget's own click() or the handler it is wired to, on the real GUI against a live ROS graph -- the same route verify\_gui\_buttons.py uses.} \\
\addlinespace[3pt]
{\scriptsize\ttfamily 31f1d6a} & \textbf{docs: 15 of 15 clips show the task completing, measured not judged}\newline{\small Records the carry distances per mode and the VR half-scale confirmation, the six GUI tutorials, and the one stale leftover directory from the superseded study-taskset sweep.} \\
\addlinespace[3pt]
{\scriptsize\ttfamily 90e7920} & \textbf{Three one-liners for things that were believed working}\newline{\small 1. recovery\_manager IS NOW LAUNCHED. It owns fault paths (b)-(e) -- Teensy reconnect, Kortex session loss, camera/zero-detection abort, network dropout -- each verified by fault injection and NEVER STARTED BY ANYTHING. Installed, documented, inert. A recovery layer that is not running is indistinguishable from one that is, right up until a fault. Verified in the graph: `recovery\_manager up.} \\
\addlinespace[3pt]
{\scriptsize\ttfamily b67739f} & \textbf{GUI resolved to one: console features ported, task set trimmed to A/B/C}\newline{\small srl\_gui.py is now the strict superset and the only launcher a lab session needs. Three things happened, and the third was a real defect. PORTED FROM srl\_console (the Dear PyGui tool it supersedes): Charts rolling strips for end-effector height, clearance to the wearer against its floor, and sim->real divergence. The divergence strip renders "no data -- not a flat line" when nothing is publishing, because a flat trace at zero is exactly what a missing real arm must never look like.} \\
\addlinespace[3pt]
{\scriptsize\ttfamily 7f48aff} & \textbf{Clip frame rate: measured, not raised -- 56\% of every clip was dead lead-in}\newline{\small "5-13 fps against 12 requested" turned out to be three separate things, and raising the capture rate would have fixed none of them. WHAT THE CONTAINER HEADER PROVES: nothing. x11grab always writes its requested rate -- when the source has not repainted it grabs the same pixels again -- so every clip on disk reported 12 fps whatever RViz managed. Counting frames that DIFFER from their predecessor is the only honest measure, and across nine clips it gave 0.6-1.3 fps.} \\
\addlinespace[3pt]
{\scriptsize\ttfamily 1c5a01c} & \textbf{Stale numbers re-measured, and three regressions the archive move caused}\newline{\small ONE PLACE WITH LIVE NUMBERS. docs/ENGINEERING_LOG.md now carries a CURRENT COUNTS block stating what was measured today, with the command that measures each. The historical log keeps its dated figures untouched -- "36 unit tests pass" in a 2026-08-05 section was true on 2026-08-05 and rewriting it would destroy the record. What was wrong was having no current figure to read instead.} \\
\addlinespace[3pt]
{\scriptsize\ttfamily 251a5ed} & \textbf{Consolidate: one owner for the grasp thresholds, three verifiers archived}\newline{\small GRASP CONSTANTS: six definitions -> one. srl\_teleop.gripper\_state now owns grip\_for(), holding(), and every threshold; record\_rviz, clip\_tasks, vr\_gripper\_node and verify\_gripper\_shutdown import them. TWO HAD ALREADY DRIFTED, and neither would have crashed: verify\_gripper\_shutdown FREE\_AIR\_RAD = 0.80 gripper\_state FREE\_AIR\_RAD = 0.74 the same name for two different quantities -- one is the knuckle's mechanical stop, the other the bound that separates "closed on nothing" from "closed on something".} \\
\addlinespace[3pt]
{\scriptsize\ttfamily d787d97} & \textbf{Seven gripper clips are entirely black, and the verifier could not see them}\newline{\small FOUND BY LOOKING AT THE FRAMES, which is the only reason it was found. verify\_rviz\_clips.py globbed \code{rviz\_front.mp4} and opened nothing else, so "24 of 24 clips pass" meant 24 FRONT VIEWS passed. Measuring all eight angles directly: 01\_master\_teleop/C/S1\_present\_probe gripper mean 1.0/255, 0 changing frames 02\_vr\_teleop/C/S1\_present\_probe gripper black 03\_shared\_autonomy/A/S1\_single\_arm gripper black 03\_shared\_autonomy/B/S1\_bimanual gripper black 04\_vr\_shared/B/S1\_bimanual gripper black 04\_vr\_shared/C/S1\_present\_pr [...]} \\
\addlinespace[3pt]
{\scriptsize\ttfamily 2ba2ce8} & \textbf{Black gripper views: root-caused to pkill leaving a live process that renders nothing}\newline{\small DIAGNOSED, NOT GUESSED. Reproduced by alternating the active arm four times and grabbing the display each time: step 0 arm=left brightness 99.00 left\_rviz=True right\_rviz=False step 1 arm=right brightness 98.92 left\_rviz=True right\_rviz=True step 2 arm=left brightness 0.04 left\_rviz=True right\_rviz=True step 3 arm=right brightness 0.04 left\_rviz=True right\_rviz=True The gripper camera is bolted to <arm>\_end\_effector\_link, so it needs a different config per arm, and ensure\_display kept ONE display for it: on an arm [...]} \\
\addlinespace[3pt]
{\scriptsize\ttfamily 6bec842} & \textbf{Task A is a carry now; framing on the workspace; the verifier can fail}\newline{\small TASK A GEOMETRY. A\_BIN 0.30 -> 0.62, so the transport is 0.300 m lateral against 0.130 m vertical. It was 0.132 m of which 0.130 was vertical: the block hovered directly above the bin and the "place" lowered it 130 mm. scripts/sweep\_task\_a\_bin.py walks the bin outward, verifying the WHOLE densified path at N=10 per candidate rather than checking one chosen value: x = 0.35 .. 0.68 PASS (lateral 0.030 .. 0.360 m) x = 0.71 fail <- the IK boundary, 3863 calls IK IS NOT THE BINDING CONSTRAINT -- the bench is.} \\
\addlinespace[3pt]
{\scriptsize\ttfamily 0b29ed3} & \textbf{Re-record mode 01: task A is a carry, framing tightened, grasps inside the clip}\newline{\small 01/A grasp at +5.6 s of 15.7 s 01/B grasp at +4.4 s of 12.4 s 01/C grasp at +5.0 s of 13.2 s WHAT I SEE, having looked at contact sheets of all three: A the orange block is gripped at the right of the bench, lifted, carried ACROSS to the teal bin and is visibly INSIDE it in the final frame. This is the change that mattered -- the previous version lowered the block 130 mm into a bin directly beneath it.} \\
\addlinespace[3pt]
{\scriptsize\ttfamily a609c04} & \textbf{Re-record all 15: every clip places its object at 0 mm from target}\newline{\small mode A B C 01\_master\_teleop 0 mm placed held 02\_vr\_teleop 0 mm placed held 03\_shared\_autonomy 0 mm placed held 04\_vr\_shared 0 mm placed held 06\_full\_autonomy 0 mm placed held every grasp lands inside its clip: +0.8 to +26.6 s of 12-34 s THREE DEFECTS FOUND BY RECORDING AND LOOKING, in the order they surfaced. 1. THE OBJECT WAS PINNED TO THE WRIST, NOT THE FINGERS.} \\
\addlinespace[3pt]
{\scriptsize\ttfamily d2b3918} & \textbf{FROZEN check used a statistic that dwell dilutes; use the max instead}\newline{\small Two clips that visibly move were failed as FROZEN at delta 0.011-0.012 against a 0.02 floor. The statistic was the mean absolute frame-to-frame difference over EVERY frame and EVERY pixel, so a small arm moving in a mostly-static frame scores near zero once the clip carries a dwell -- and task A now holds 4.5 s at the bin on purpose, so the check was penalising the fix for the previous defect. The question the check asks is "did anything ever happen".} \\
\addlinespace[3pt]
{\scriptsize\ttfamily 521f510} & \textbf{docs/ENGINEERING_LOG.md current counts: the 15 clips and their placement accuracy} \\
\addlinespace[3pt]
{\scriptsize\ttfamily 33cf23c} & \textbf{Clip verification green: 54 of 54, all eight angles per clip}\newline{\small 15 clips indexed, 15 verified 54 of 54 clips pass every automatic check (includes the older 00\_unclassified set) Every angle of every clip is now checked for a picture and for motion -- the pass that previously read "24 of 24" while seven gripper views were black for 34 s each, because the tool globbed rviz\_front.mp4 and opened nothing else.} \\
\addlinespace[3pt]
{\scriptsize\ttfamily 850a625} & \textbf{Part 1 (INCOMPLETE): the bench was decoration and the objects floated}\newline{\small REPORTING THIS AS UNFINISHED. The scene machinery is built and verified; the task geometry that has to sit on it does NOT yet verify, so nothing is re-recorded. WHAT WAS WRONG, both confirmed by measurement: 1. THE BENCH WAS DECORATION. clip\_scene published a MarkerArray to /task\_objects and NOTHING to /planning\_scene. Asked directly, move\_group returned "collision objects: NONE" while the bench looked solid on screen.} \\
\addlinespace[3pt]
{\scriptsize\ttfamily d9fc66c} & \textbf{Part 1 DONE: the sweep was self-cancelling; 40 failures -> 0 with the bench in}\newline{\small THE HARNESS BUG, FIRST, BECAUSE THE GEOMETRY WAS NEVER THE PROBLEM. 40 failures at every edge from 0.10 to 0.26 m. The cause: on\_bench() placed objects at BENCH\_NEAR\_Y + depth/2, so moving the edge moved every object with it and the wrist-to-edge geometry -- the only thing that decides whether the gripper hits the bench -- was INVARIANT. The parameter had been cancelled out of the quantity being scored. A flat score was the harness saying so.} \\
\addlinespace[3pt]
{\scriptsize\ttfamily 7b67a3f} & \textbf{Empty-world audit: the wearer was fine, the WORLD was not; guard added}\newline{\small THE WEARER IS LOADED, AND I WAS WRONG TO SAY OTHERWISE FIRST. My initial check looked for \code{human\_torso} and \code{human\_head}, found neither, and I raised the alarm. The links are named \code{torso}, \code{head}, \code{hips}, \code{left\_leg}, \code{right\_leg} plus the six human arm links -- 11 in the runtime URDF, all present in TF. A wrong name reads exactly like a missing link, which is the same class of error as everything else in this audit.} \\
\addlinespace[3pt]
{\scriptsize\ttfamily 083dc32} & \textbf{The right arm CAN support B and C: standoff was the lever, not the workspace}\newline{\small YOU ASKED WHETHER THE COLLAPSED BAND MEANS THE WORKSPACE CANNOT CARRY THE TASKS. It does not, and here is why the two figures are not in conflict. The band that collapsed 15/19 -> 1/19 was probed at a FIXED z = 1.15 and x = +/-0.32. The tasks do not live there: they live at the bench edge, y = 0.09..0.20, z = 1.06..1.20. Same rig, different place. "It passes" and "the envelope collapsed" were both true and neither answered the question that matters -- HOW CLOSE TO THE EDGE IS EVERY POSE.} \\
\addlinespace[3pt]
{\scriptsize\ttfamily 26adc92} & \textbf{Part 3: scene calibration, and the tolerance is +/-10 mm not +/-20 mm}\newline{\small WHAT scene\_fingerprint\_node DOES TODAY: consumes /perception/objects (it never opens a camera -- two processes on one camera is the serial-port fault in a new costume), accumulates during a sweep window, compares against config/scene\_fingerprint.json at 15 mm / 15 deg with a 250 mm gate, and re-registers ONLY what moved.} \\
\addlinespace[3pt]
{\scriptsize\ttfamily ccfb27b} & \textbf{Part 4 (font sub-item): Helvetica resolves, and it was never being requested}\newline{\small "Helvetica is not resolving" was not a fontconfig problem. Checked at both layers before changing anything: fc-match Helvetica -> LiberationSans-Regular.ttf QFontInfo(QFont("Helvetica")) -> "Liberation Sans" Both already resolved. The actual cause: NOTHING EVER ASKED FOR HELVETICA. srl\_hud.py requested "DejaVu Sans" and "DejaVu Sans Mono" throughout, so the alias was never exercised and the GUI simply rendered in DejaVu.} \\
\addlinespace[3pt]
{\scriptsize\ttfamily 147e704} & \textbf{CAD analysis: the alternation is confirmed and the interior chain agrees to 1 mm}\newline{\small /mnt/c IS HEALTHY. The I/O errors in docs/ENGINEERING_LOG.md are gone -- a WSL restart fixed the dead 9p mount at some point. No remount needed, nothing for you to run from Windows. READER: cadquery-ocp. pip wheels, no root, no conda, and it is the full OpenCASCADE kernel rather than a parser -- bounding boxes, volumes and surface classification all need the kernel. pythonocc-core is conda-first with unreliable pip wheels; headless FreeCAD needs apt and \textasciitilde 1 GB for the same kernel.} \\
\addlinespace[3pt]
{\scriptsize\ttfamily 820351c} & \textbf{pos\_tol 15 -> 10 mm, and the lab placement requirement recorded}\newline{\small (a) THE TOLERANCE WAS LOOSER THAN THE THING IT PROTECTS. The grasp's capture half-window is 17.5 mm on the widest object, so at 15 mm an object could be declared UNCHANGED by the fingerprint and then not be picked up.} \\
\addlinespace[3pt]
{\scriptsize\ttfamily f16b15f} & \textbf{Pure-CAD master arm URDF: builds, parses, and looks like the printed arm}\newline{\small PURE CAD, NOT HYBRID, as decided. Nothing in master\_arm\_cad.urdf.xacro comes from the measured lengths -- mixing them in would throw away the only thing this model is for, which is being an independent check. srl\_dual.urdf.xacro is UNTOUCHED.} \\
\addlinespace[3pt]
{\scriptsize\ttfamily 970d212} & \textbf{Perception without fiducials: the accuracy is in the FIT, not the detector}\newline{\small THE QUESTION WAS ASKED WRONG BEFORE. "Can vision hit sigma <= 2 mm" treats the detector as the thing that localises. It is not. The pipeline is detect -> crop the depth -> fit, and the box only has to point at the right REGION. So the two halves separate, and one of them is measurable here rigorously. MEASURABLE HERE, by the project's own rule -- synthetic is legitimate where the ground truth is CONSTRUCTED, not RENDERED.} \\
\addlinespace[3pt]
{\scriptsize\ttfamily 11e5f21} & \textbf{Lab procedure for the detector measurement, with what NOT to re-measure}\newline{\small The lab session is the only route that uses our objects in our lighting, so the procedure says plainly what is already settled and what is not. ALREADY KNOWN, do not re-measure: the pose accuracy is not the detector's job. The known-dimension fit gives 1.2 / 1.8 / 8.0 mm and tolerates a sloppy box; the naive centroid is wrong by 13-45 mm; range barely matters. WHAT THE LAB MUST MEASURE: the DETECTION RATE and the box quality.} \\
\addlinespace[3pt]
{\scriptsize\ttfamily 346f5b6} & \textbf{LM-O evaluation harness, ready to run when the download lands}\newline{\small measure\_detector\_on\_real\_rgbd.py runs the SAME estimator as the constructed measurement -- segment to the nearest surface, then fit known dimensions -- but against REAL RGB-D with real sensor noise and ground-truth 6D poses, and with the detector given nothing but an open-vocabulary prompt. It reports detection rate, median box IoU against ground truth, and the known-fit pose error with the fraction inside 10 mm.} \\
\addlinespace[3pt]
{\scriptsize\ttfamily 62eb61f} & \textbf{Multimeter respec'd 50x90x130 -> 30x70x45; LM-O measured; architecture written}\newline{\small 1. THE MULTIMETER, FIXED. It bound the object set three ways at once, and the three constraints pull on DIFFERENT dimensions -- so only changing both works: candidate pose err capture sigma<=2 50 x 90 x 130 (old) 8.0 mm 17.5 mm no 35 x 90 x 130 7.9 mm 25.0 mm no <- narrow only: pose barely moves 50 x 90 x 60 2.2 mm 17.5 mm no <- shallow only: capture unchanged 30 x 70 x 45 1.7 mm 27.5 mm YES <- all three WIDTH sets the capture window ((85-w)/2); DEPTH conditions the fit, because a deep object is seen from one face [...]} \\
\addlinespace[3pt]
{\scriptsize\ttfamily 3a950cc} & \textbf{PART 5: the empty RViz panel was a QoS durability mismatch, order-dependent}\newline{\small ROOT CAUSE, and it explains why it looked intermittent rather than broken. robot\_state\_publisher LATCHES /robot\_description -- one message at startup, TRANSIENT\_LOCAL, never republished. Confirmed from the live graph: publisher durability = 1 (TRANSIENT\_LOCAL). verification\_capture.rviz asked for it with Durability Policy: Volatile A volatile subscriber that arrives AFTER the latched message was sent never receives it, so the RobotModel has no description and the panel renders an empty scene.} \\
\addlinespace[3pt]
{\scriptsize\ttfamily c0d99e0} & \textbf{(a) CAD panel root-caused: a config key in the wrong place is SILENTLY IGNORED}\newline{\small SAME CLASS AS THE TRANSIENT-LOCAL BUG, exactly as you said it would be. ROOT CAUSE: the hand-written config put \code{Fixed Frame} directly under \code{Visualization Manager} instead of under \code{Global Options:}. RViz does not error on a misplaced key -- it silently uses the DEFAULT -- so the fixed frame fell back to one that does not exist, link transforms could not be resolved, and the panel rendered a single stub. \code{Marker Scale: 0.05} being ignored was the corroborating clue I had already seen and not followed.} \\
\addlinespace[3pt]
{\scriptsize\ttfamily 4df7bbf} & \textbf{PART 6: READY TO RUN, the participant session, and failure handling}\newline{\small NEVER GREEN ON UNKNOWN --- enforced STRUCTURALLY, not by remembering. srl\_experiments/readiness.py has no boolean anywhere, because a boolean cannot represent "nobody asked". Three states: OK / FAIL / UNKNOWN. Readiness is built from REQUIRED with every entry already UNKNOWN, and checks OVERWRITE. A check that crashes, is skipped, or was never wired writes nothing -- so it CANNOT write OK. Adding a name to REQUIRED and forgetting to implement it makes the system NOT READY, which is the correct direction to fail in.} \\
\addlinespace[3pt]
{\scriptsize\ttfamily 136674e} & \textbf{Tab-selection gap closed --- and it was TWO silent failures compounding}\newline{\small WHY THE PANEL COULD NOT BE CAPTURED FROM THE RUNNING GUI: 1. My \code{--tab} argument was never added. The patch matched \code{ap.add\_argument("--no-rviz", action="store\_true")} but the real line is multi-line with a help= on the next, so the replace silently did nothing. 2. \code{parse\_known\_args} then SWALLOWED \code{--tab ready} without a word. It is there so ROS can have its --ros-args, and it accepted an unknown flag silently.} \\
\addlinespace[3pt]
{\scriptsize\ttfamily 382b280} & \textbf{Part 4: the arm panels became diagrams, and one of them found the seam risk}\newline{\small Both master arms and both robot arms are now kinematic diagrams rather than tables. The master chain is drawn at TRUE bend angles with a uniform scale-to-fit, so the picture IS the posture; roll joints carry a rotating index mark, bend joints are hinges, and under each node sits the live angle and the TIME SINCE THE LAST DISTINCT VALUE -- the second number being the one this rig actually needs, because its characteristic failure is data that arrives on time and never changes.} \\
\addlinespace[3pt]
{\scriptsize\ttfamily 3ca6b4b} & \textbf{Archive the superseded clips: they are evidence, and the note says of what}\newline{\small recordings/verification/ -> extras/archive/recordings/verification\_20260810/, moved under the standing never-delete rule with GEOMETRY\_NOTE.md stating exactly what geometry they were recorded against and why they were replaced: planning scene EMPTY -- clip\_scene published a MarkerArray to /task\_objects and nothing at all to /planning\_scene the bench DECORATION -- avoid\_collisions could not see it, so the arm swept through a surface that looked solid on screen the objects FLOATING -- 0.107 to 0.357 m of clear air under eve [...]} \\
\addlinespace[3pt]
{\scriptsize\ttfamily c2d0b46} & \textbf{Re-record mode 06\_full\_autonomy --- and two grasp bugs found by the first take}\newline{\small 3 of 3 clips, 8 angles each, against the verified geometry. Against the stack's own record over the three runs: 54 of 54 IK windows at 100\%, mount guard PASS at 0.1610 m worst clearance (floor 0.15), zero clearance-floor blocks. The one block was a 0.38 s flip\_reject on the right arm that self-cleared -- the guard doing its documented job, not an IK failure. TWO REAL BUGS, BOTH FOUND BY THE FIRST TAKE FAILING RATHER THAN BY READING. 1. gripper\_state.holding(knuckle, width\_mm) HAD ONLY A LOWER BOUND.} \\
\addlinespace[3pt]
{\scriptsize\ttfamily cf14987} & \textbf{Re-record mode 01\_master\_teleop --- 3 of 3, the grasp fixes hold on a second path}\newline{\small A grasp at +5.2 s of 19.1 s, placed 0 mm from target B grasp at +4.4 s of 14.8 s C grasp at +3.8 s of 13.7 s The two grasp fixes from the previous commit were found on the autonomy path; this is the first evidence they hold on a different one. No spurious t=0 grasp, no grasp outside the window, on any of the three tasks.} \\
\addlinespace[3pt]
{\scriptsize\ttfamily dd4c13e} & \textbf{Re-record mode 03\_shared\_autonomy --- 3 of 3, and the clips are 13.4 fps not 5.4}\newline{\small A grasp at +4.6 s of 17.5 s, placed 0 mm from target B grasp at +20.7 s of 30.1 s C grasp at +4.3 s of 13.6 s B's 30.1 s against \textasciitilde 14 s on the other paths is the arbiter blending toward the target rather than a stall -- checked by looking, not assumed: both arms are engaged in every sampled frame, the left holds its pose throughout, and the right travels continuously to the part and on to the green box. DELIVERED FRAME RATE RE-MEASURED. 403 frames over a 30.1 s window = 13.4 fps against 15 requested.} \\
\addlinespace[3pt]
{\scriptsize\ttfamily 2fe0899} & \textbf{Mode 02\_vr\_teleop: A recorded, B and C did NOT --- reported, not papered over}\newline{\small A grasp at +15.8 s of 29.2 s, placed 0 mm from target OK B run exited 0, NO GRASP RECORDED at all FAIL C run exited 1, NO GRASP RECORDED at all FAIL NOT RE-RUN AND HOPED. What I can say from looking at B's frames: the right arm DOES travel to the part and ends up beside it; the part never moves; the burnt-in caption correctly turns red and reads "DID NOT COMPLETE -- run exited 0: NO GRASP RECORDED at all". So the transport moves the arm and the failure is at the grasp, not at the command path.} \\
\addlinespace[3pt]
{\scriptsize\ttfamily 0d774dd} & \textbf{Record the re-record's state and the one open VR question in NEXT\_SESSION}\newline{\small Three modes complete (9 clips), VR partial (1 of 3), 04\_vr\_shared not run. Names the measurement that settles the VR grasp question -- pad-to-object distance at closest approach -- and says explicitly not to tune the 30 mm gate or the mapper before taking it.} \\
\addlinespace[3pt]
{\scriptsize\ttfamily 48ec186} & \textbf{VR grasp failure SETTLED by measurement: neither hypothesis was right}\newline{\small The pair of numbers clip\_scene now logs answers it outright: VR A (works) min\_pad\_obj 0.000 m, knuckle 0.4136 at that instant VR B (failed) min\_pad\_obj 0.000 m, knuckle 0.0 at that instant min\_pad\_obj\_closed 0.339 m So the VR path does NOT land short -- the pads reach the part EXACTLY -- and the gate is NOT too tight. When the fingers closed they were 339 mm away, and when they were at the part they were fully open. The close was MISTIMED, and neither of the two hypotheses I offered was correct.} \\
\addlinespace[3pt]
{\scriptsize\ttfamily e642506} & \textbf{The physics gate FAILS, and it fails in both directions}\newline{\small The gate: close the shipped Robotiq 2F-85 on a 40 mm cube in Gazebo, hold it without jitter, drop it when released. Rig is the real Gen3 + 2F-85 hung upside down from a fixed post, every joint at zero, so the approach is top-down with no IK anywhere and the lift is a fold that preserves tool orientation exactly. No ros2\_control, no MoveIt, no ROS. Result at the SHIPPED fingertip friction (mu=100000): lifts 0.256 m, survives a 90 deg arm rotation at 0.24-0.48 mm RMS -- and will not let go.} \\
\addlinespace[3pt]
\end{longtable}

\section*{2026-08-11 \textnormal{\small\textcolor{srlgrey}{--- 60 commits}}}
\addcontentsline{toc}{section}{2026-08-11}
\begin{longtable}{@{}p{0.055\textwidth}p{0.885\textwidth}@{}}
\endfirsthead\endhead
{\scriptsize\ttfamily 264529a} & \textbf{Section 3: T3 replaces Task B; ALOHA written up; the randomisation region is EMPTY}\newline{\small THREE PARTS, and the third is the one with teeth. 1. T3 REPLACES TASK B. The slot keeps the key "B" -- renaming it would break the dispatcher, the GUI buttons and run\_experiment.sh, which is the exact "button labelled T3 runs something else" failure the archive README exists to prevent. The CONTENT is T3 and only T3: the compliant (sling) half is gone, 18 trial slots become 9. WHICH T3, because there are two and only one is usable.} \\
\addlinespace[3pt]
{\scriptsize\ttfamily c16737c} & \textbf{Archive README: the two T3s, and which one is the regression}\newline{\small Verified after the scope revert that the last commit touched ONLY the B slot: TASK\_A and TASK\_C are byte-identical to their pre-change text, the Task C header and OPTIONAL\_PICK\_PLACE are present, and the archived file contains no A or C content. Nothing needs restoring.} \\
\addlinespace[3pt]
{\scriptsize\ttfamily 6591a83} & \textbf{Gates 1-3 measured. Gate 1's premise is wrong, Gate 3 is impossible.}\newline{\small GATE 1 -- the timeline is not 2.5x over, and the cap does not need raising. The brief assumes a real-arm cap of max\_vel\_rad\_s 0.05. The shipped value is 0.15 (teleop.launch.py real\_max\_vel\_rad\_s); 0.05 is the HOMING velocity and real\_max\_step\_rad, which are different parameters.} \\
\addlinespace[3pt]
{\scriptsize\ttfamily 87e61fb} & \textbf{Gate report and experiment architecture (A-J), gates answered first}\newline{\small docs/research/12\_experiment\_gates\_and\_architecture.md. Gates are measured, the architecture is drawn around what they returned rather than around the brief's assumptions. Counterbalancing checked rather than assumed: williams\_square(5) returns 10 sequences (odd n needs 2n), and check\_balance says n=15 gives position imbalance 0 but SEQUENCE imbalance 2. n=10 and n=20 are perfect. Recruit 20 or state the residual -- do not use 15 and call it counterbalanced.} \\
\addlinespace[3pt]
{\scriptsize\ttfamily 69ee55b} & \textbf{Build 1/5: the task layer, five adapters, and the mode-independence PROOF}\newline{\small Build order item 1, done before any task code exists, because if the property does not hold then nothing built on top is comparable. task\_actions.py -- seven verbs (REACH\_TARGET, PICK\_OBJECT, PLACE\_OBJECT, GRASP, RELEASE, TRANSFER, RETURN) over an Adapter with exactly four methods. It names no mode. Phases are part of the trace signature so an adapter cannot quietly skip one -- teleporting to the object would produce a trace missing DESCEND and the proof would catch it.} \\
\addlinespace[3pt]
{\scriptsize\ttfamily 7a11b89} & \textbf{Build 2/5: T0 target reaching -- seeded, region-bounded, and it REFUSES}\newline{\small Extends experiments/abc/task0.py rather than adding a parallel spec. The fixed A/B/C/D set stays as the Fitts CALIBRATION; the randomised L1-L3 / R1-R3 set is the study instrument. Both share the same measured BAND so neither can drift from the geometry the other was verified against. sample\_targets(arm, seed) 3 spheres per arm, reproducible from the seed, drawn ONLY from the band verified with the bench in the planning scene, with an optional live-IK validator applied to every candidate.} \\
\addlinespace[3pt]
{\scriptsize\ttfamily dcea886} & \textbf{T1 and T2 re-measured against the corrected spec, before building either}\newline{\small T1 -- THE CUBE COSTS NOTHING AGAINST A CYLINDER, and the switch was unnecessary. A cube admits four grasp yaws, a cylinder all of them; measured on identical poses in each arm's own graspable strip: fixed cube(4 yaw) cylinder(8) top-down left 12/12 100\% 12/12 12/12 12/12 right 5/12 42\% 11/12 92\% 11/12 92\% 12/12 100\% 11/12 either way. Gate 2's 5/12-vs-11/12 was FIXED versus YAW-FREE, not cube versus cylinder.} \\
\addlinespace[3pt]
{\scriptsize\ttfamily 4f53151} & \textbf{T1 layout does NOT fit at the current bench edge -- measured, holding before changes}\newline{\small Option 3 (left arm only) is recorded, and three of its four reasons hold. The fourth -- "four cubes and both planes fit in the left arm's region" -- is FALSE at the current bench edge, and the brief said to come back rather than change anything, so nothing has been changed. scripts/verify\_t1\_layout.py searches rather than testing one layout: it enumerates every cell that is BOTH supported (near face at or behind the bench edge, i.e.} \\
\addlinespace[3pt]
{\scriptsize\ttfamily d44dbb8} & \textbf{The general constraint: support must come from BEHIND; T0 is the control}\newline{\small Two things to carry forward, recorded rather than remembered. 1. WHY A SHELF CANNOT HELP (docs/system/12\_randomisation\_region.md). The pinned tool axis is measured at (-0.153, +0.846, +0.511): 120.8 deg from straight down and 30.7 deg ABOVE horizontal. The hand enters from the near side and from BELOW, and the fingers close underneath the object and lift into it.} \\
\addlinespace[3pt]
{\scriptsize\ttfamily c7336dd} & \textbf{The rail does NOT restore MASTER\_TELEOP. Modelled, swept, and controlled.}\newline{\small A cantilevered lip flush with the bench top, thin in z so the space BENEATH it stays free for the fingers -- the only support geometry that can satisfy the 30.7-degree near-side-from-below approach. Swept at three edge positions, bench in scene, full densified pick path (standoff, descend, lift). rail edge FIXED cells TOPDOWN cells 6-item layout none (0.245) 0 18 no (18564 combos searched) 0.230 0 19 no (27132) 0.210 0 20 no (38760) 0.190 0 22 no (74613) FIXED IS ZERO AT EVERY RAIL POSITION.} \\
\addlinespace[3pt]
{\scriptsize\ttfamily 76631ff} & \textbf{Parts 1 and 2: T1 settled on option 4 with a verified layout; regression fixed}\newline{\small PART 1 -- T1 SETTLED. Options tested in the order the brief set. Option 2 (smaller standoff) cannot even be tested: at the current bench edge FIXED has zero supported cells on the GRASP POSE alone, so there is nothing for a smaller standoff to rescue, and with a rail at y=0.19 the grasp pose is still zero -- the rail occupies exactly the volume the fingers enter, as the general constraint predicts. --sweep-standoff REFUSED rather than reporting a ladder with no rungs.} \\
\addlinespace[3pt]
{\scriptsize\ttfamily 23c3a25} & \textbf{Blockers A and B settled; T2 re-spec'd; T3 built; all four MSc tasks verify 0}\newline{\small BLOCKER A -- the free-space check was the broken one, and it was hiding more than T2. Settled three independent ways, bench in scene: ARITHMETIC two declared T2 grip poses -- S2's start at (+/-0.25, 0.35, 1.10) -- lie INSIDE the bench slab (z 1.06-1.10, y 0.245-0.63). No solver involved, so it cannot be a solver artefact. PER WAYPOINT N=10, both arm assignments: S1 4/4 fail, S2 9/11, S3 8/10 = 21, exactly the number the regression block reported. CLEAR BAND feasible from z=1.28 at y=0.35 and z=1.30 at y=0.38.} \\
\addlinespace[3pt]
{\scriptsize\ttfamily e874da1} & \textbf{Resume point: parts 3-4 done, part 5 (recording) is the next boundary} \\
\addlinespace[3pt]
{\scriptsize\ttfamily b038d9c} & \textbf{Part 5 start: archive the superseded clips, record the two distinctions}\newline{\small ARCHIVED the 15 A/B/C clips to extras/archive/recordings/verification\_20260811/ with a GEOMETRY\_NOTE naming BOTH reasons they are superseded, not just the obvious one: 1. Every Task B clip was recorded against the declared band z 1.10-1.30, which is now known to be inside the bench slab (z 1.06-1.10). Two of its declared grip poses are geometrically INSIDE the slab -- arithmetic, no solver -- and per-waypoint at N=10 with the bench in scene the three paths fail 4/4, 9/11 and 8/10.} \\
\addlinespace[3pt]
{\scriptsize\ttfamily 1180cec} & \textbf{Part 5 wiring: clip paths for all four MSc tasks, BUILT but NOT YET VERIFIED}\newline{\small msc\_clip\_tasks.py -- the recording paths for T0/T1/T2/T3, in the same registry shape as clip\_tasks.TASKS so record\_abc\_sweep.py can drive them without a second sweep. Extends, does not duplicate. Every coordinate comes from an already-verified source and nothing new is invented: T0 from task0.sample\_trial at a FIXED seed (the randomised L1-L3 / R1-R3 set is the study instrument -- TARGETS\_LEFT is the A/B/C/D Fitts calibration set, and a clip must show what a participant sees), T1 from the option-4 layout, T2 from t [...]} \\
\addlinespace[3pt]
{\scriptsize\ttfamily 52cd314} & \textbf{Resume point: clip paths unverified, no clip recorded} \\
\addlinespace[3pt]
{\scriptsize\ttfamily ccefd1d} & \textbf{Clip paths VERIFIED through IK -- and the first run of this was vacuous}\newline{\small THE CLIP PATHS NOW PASS: 103 distinct waypoints at N=10, both arms, bench in scene, 0 failures. 5264 IK calls. t0 L 9 distinct 0 fail R 8 distinct 0 fail t1 L 61 distinct 0 fail R 1 distinct 0 fail t2 L 4 distinct 0 fail R 4 distinct 0 fail t3 L 8 distinct 0 fail R 8 distinct 0 fail Controls all correct, including the must-fail one: a pose inside the bench slab comes back unreachable. THE FIRST RUN OF THIS WAS A VACUOUS PASS AND I NEARLY REPORTED IT.} \\
\addlinespace[3pt]
{\scriptsize\ttfamily 5c5465a} & \textbf{clip\_scene: item tables for all four MSc tasks, and the guard T0 needs}\newline{\small Sweep wiring, continued. clip\_scene.py now has an item table for t0/t1/t2/t3: t0 NO OBJECTS -- and tick() now RETURNS EARLY on an empty item dict. Without that guard \code{next(iter(self.items.values()))} raises StopIteration; with a placeholder object instead, the one task whose whole point is that there is nothing to grasp would show something to grasp. T0 runs in every mode precisely because it has no object.} \\
\addlinespace[3pt]
{\scriptsize\ttfamily ef7a90d} & \textbf{Sweep wiring complete -- and the MSc task names COLLIDE with an archived refusal}\newline{\small Both wiring items done, plus one thing the resume point had not anticipated. THE COLLISION. run\_experiment.sh refuses t1|t2|...|t9 BY NAME -- they are the archived 300/310 mm bimanual set, and the refusal is the guard that stops a run producing a plausible number against the wrong tilt criterion. The MSc brief also calls its tasks T0-T3.} \\
\addlinespace[3pt]
{\scriptsize\ttfamily 8f71470} & \textbf{Mode 06: 32 clips recorded, ALL FOUR TASKS FAILED. Not a completed mode.}\newline{\small The sweep ran end to end for the first time on the MSc set and REFUSED all four tasks: FAIL run exited 1; no scene\_events.json -- the scene never ran 32 mp4s exist (4 tasks x 8 angles) and 0 scene\_events.json. The clips render; nothing has confirmed an object was ever in them. They are NOT a recorded mode and are not citable. Reported as a failure rather than as 32 files. Three bugs found by running it, each of which hid the next:} \\
\addlinespace[3pt]
{\scriptsize\ttfamily 7d2a5ce} & \textbf{Clean up: the sweep leaks one Xvfb per angle, and the locks are self-worsening}\newline{\small Found after the failed mode-06 run: 8 orphaned Xvfb servers -- one per angle -- and 8 stale /tmp/.X9*-lock files, left because the sweep failed before its teardown ran. Killed by explicit PID, locks removed by explicit path. Recorded in NEXT\_SESSION because it is not tidiness.} \\
\addlinespace[3pt]
{\scriptsize\ttfamily 5a6a597} & \textbf{Xvfb teardown on EVERY exit, proven by killing mid-angle; 32 bad clips deleted}\newline{\small ITEM 1 -- THE LEAK IS FIXED AND THE FIX IS EXERCISED ON THE FAILURE PATH. record\_rviz.teardown\_displays() kills every Xvfb and RViz BY PID -- pkill -f has killed the shell running it three times in this project -- and removes /tmp/.X<n>-lock and the socket for each of the 8 views. record\_rviz.stale\_displays() the precondition, returning what is live or still locked so the sweep can NAME it.} \\
\addlinespace[3pt]
{\scriptsize\ttfamily 7df7c81} & \textbf{Xvfb teardown proven on the failure path; three recording bugs found by running}\newline{\small THE LEAK IS FIXED AND THE FIX IS EXERCISED WHERE IT MATTERS. teardown on try/finally AND on SIGINT/SIGTERM/SIGHUP -- a \code{finally} does not cover the signal, which is how this sweep actually ends when a session runs out. Proven by SIGTERMing a run mid-angle: "killed 16 process(es), removed 8 lock(s)", then XVFB=0 LOCKS=0. A precondition now refuses on any live Xvfb or stale lock, naming each. THE VERIFIER'S BLIND SPOTS were already closed; this session VERIFIED that rather than assuming it.} \\
\addlinespace[3pt]
{\scriptsize\ttfamily 0e8e39c} & \textbf{Recording blocked: no joint\_state\_broadcaster. Chain diagnosed to one contradiction.}\newline{\small VERIFIER (item 1) -- all three named blind spots green, re-run BEFORE any recording as asked: 7/7 controls correct including black\_capture and hud\_text\_only, plus the per-angle control in both directions -- "black gripper is CAUGHT" and "all-good clip is not flagged". It refuses to report until those pass. RECORDING (item 2) -- still blocked, but the chain is now established rather than suspected: probe: TIMEOUT joint\_state\_msgs=0 ik=True followers=1 1.} \\
\addlinespace[3pt]
{\scriptsize\ttfamily 2be27d6} & \textbf{The respawn storm fixed at SOURCE; my master:=false diagnosis was wrong}\newline{\small ITEM 1 -- WHICH LINK WAS LYING: NONE OF THEM. Run directly, \code{run\_teleop.sh gate:=false master:=false} yields 0 master\_pose\_node lines, so "\$@" forwarding and the launch file's IfCondition are both honest, and sim\_session passes the flag correctly. MY PREVIOUS REPORT WAS WRONG.} \\
\addlinespace[3pt]
{\scriptsize\ttfamily 0a80b7f} & \textbf{Mode 06 recorded: 2 of 4 pass the sweep, and LOOKING shows none is correct yet}\newline{\small FIRST CLIPS THAT ARE NOT WORTHLESS. The stack comes up, the arms move, the scene runs, captions burn in, and scene\_events.json is written. T0 and T2 passed every automatic check; T1 and T3 failed with "NO GRASP RECORDED at all". THEN I EXTRACTED FRAMES AND LOOKED, WHICH IS THE POINT OF ITEM 4, AND BOTH "PASSING" CLIPS ARE ALSO WRONG: T0 Both arms reach out, caption right, arms clear of the wearer -- and there are NO SPHERES AND NO BENCH in the frame.} \\
\addlinespace[3pt]
{\scriptsize\ttfamily 24e013a} & \textbf{Four mode-06 defects fixed; T1/T3 grasp diagnosed from pad position}\newline{\small T1 and T3 failed for DIFFERENT reasons, and the existing pad diagnostics said so without any guessing at schedule indices: T1 cube\_0 min\_pad\_obj\_m 0.0957 knuckle 0.4235 closed\_m 0.0957 T3 box min\_pad\_obj\_m 0.0978 knuckle 0.0 closed\_m None T1: the fingers CLOSED (knuckle 0.42) 95.7 mm from the cube, against a 30 mm gate. 0.0957 is |PAD\_OFFSET| to the millimetre.} \\
\addlinespace[3pt]
{\scriptsize\ttfamily dd89e1d} & \textbf{Mode 06 re-record: 4/4 FAIL, and the failure MOVED -- resume point written}\newline{\small The four fixes are in and the scene is right: scene\_events.json records fixtures, T0 lists its six spheres, T2 lists its ball, and the subject check's control fails as required. The re-record still failed, and reporting it as progress would be false. WHAT ACTUALLY HAPPENED: all four gated the capture on TIMEOUT, which means graph.moved() never fired -- the arm did not move -- and the runner exited 1 for its own reason, 'no arm moved at least 0.010 m'.} \\
\addlinespace[3pt]
{\scriptsize\ttfamily f92297e} & \textbf{kill\_stack() was blind: ps comm truncates at 15 chars, so it killed almost nothing}\newline{\small THIS IS WHY THE 06 RE-RECORD FAILED 4/4, AND IT IS NOT THE SCENE. stack\_pids() compared \code{ps -o comm} against a set of full node names. comm truncates at 15 characters, so three of the five entries could never match: robot\_state\_publisher (21) -> comm reports robot\_state\_pub ros2\_control\_node (17) -> ros2\_control\_no ik\_follower\_node (16) -> ik\_follower\_nod Only move\_group, rviz2 and spawner were ever killed.} \\
\addlinespace[3pt]
{\scriptsize\ttfamily 2d1b5be} & \textbf{Resume point: the 06 blocker is DISCOVERY, not the scene}\newline{\small Replaces yesterday's resume point, which drew the wrong conclusion from the right data. It blamed the scene changes for four failed recordings. They were not the cause and the pad distances it quoted were measured on an arm with no TF; that paragraph is withdrawn. The sweep was recording a graph that had 66 leaked stack processes, 10 leaked launch parents, no move\_group, no /tf publisher and 153 stale shm segments -- because kill\_stack() had never been able to see any of it (ps comm truncates at 15 chars).} \\
\addlinespace[3pt]
{\scriptsize\ttfamily 6a77a0d} & \textbf{Narrow the 06 blocker: the stack is fine, the SEPARATELY-STARTED process is isolated}\newline{\small Decisive evidence, from the stack's own log: fsr\_gripper\_node reports 'OPEN REFERENCE CONFIRMED from /joint\_states at 0.0000 rad' for both arms. So /joint\_states is published and received. The broadcaster is not dead and the topic is not dead. What fails is any process started SEPARATELY from the launch: the readiness probe (joint\_state\_msgs=0 while ik=True, three times), a fresh 'ros2 node list --no-daemon', and 'ros2 topic hz /joint\_states'. Nodes launched together see each other; a new process joins nothing.} \\
\addlinespace[3pt]
{\scriptsize\ttfamily bbbb9cb} & \textbf{The environment theory is WRONG -- tested, reverted, recorded}\newline{\small Sourcing scripts/env.sh in wait\_ready() was the one-line fix my own resume point predicted. It was tried and it made the probe WORSE: without env.sh joint\_state\_msgs=0 ik=True with env.sh joint\_state\_msgs=0 ik=False With the variables pinned the probe stopped seeing even /compute\_ik. So the missing ROS\_DOMAIN\_ID / ROS\_AUTOMATIC\_DISCOVERY\_RANGE / RMW\_IMPLEMENTATION are not what isolates a separately-started process, however well the theory fitted docs/ENGINEERING_LOG.md's own account of env.sh.} \\
\addlinespace[3pt]
{\scriptsize\ttfamily 3518ed2} & \textbf{Exclude three more candidates by measurement; the signature points elsewhere}\newline{\small (a) The machine's DDS is FINE. Sanity floor with NO stack: ros2 topic pub in one bare shell, ros2 topic hz in another -> 5.002 Hz. Separately-started processes CAN discover each other here, so 'wsl --shutdown' is not indicated and the fault is not the box. (b) Discovery RANGE is not it. teleop.launch.py pins SUBNET via setdefault and this machine has eth0/eth1 DOWN with only eth2 up, so SUBNET multicast with no healthy interface was a good theory.} \\
\addlinespace[3pt]
{\scriptsize\ttfamily 0608c22} & \textbf{ROOT CAUSE: Fast DDS shared memory. Discovery was fine; DATA never flowed}\newline{\small The probe was instrumented to print LISTS instead of counts, per the resume point, and settled it in a single run: probe sees 1 nodes: ['sim\_ready\_probe'] probe sees 81 topics /joint\_states publishers=1 subscribers=7 PRESENT in topic list Discovery was never broken. The probe found all 81 topics and counted the publishers and subscribers correctly -- it knew exactly who was there. What never arrived was DATA: zero samples on a topic with one publisher and seven subscribers already matched.} \\
\addlinespace[3pt]
{\scriptsize\ttfamily 81dcceb} & \textbf{The 06 blocker was NOT the transport: it was clearing /dev/shm under the launch}\newline{\small Five sessions of transport and environment theories all measured the same three zeros, and the reason they could not tell each other apart is that both things they were varying were fine. WHAT IS ACTUALLY WRONG. sim\_session cleared every /dev/shm/fastrtps\_* segment and launched the stack in the same instant.} \\
\addlinespace[3pt]
{\scriptsize\ttfamily 2808bbe} & \textbf{Per-task scenes: T0 loses the bench, T1 gains its planes, and "objects rest on a surface" is measured to be impossible for this wrist}\newline{\small FOUR SCENE DEFECTS, and the third one turned into a platform finding. (a) T0 HAD FURNITURE IT DOES NOT USE. It is reaching only -- no objects, no grasp, nothing to rest on -- and it was being handed the bench, the bin and a circuit box. With them came the bench's limits: T0's target band was derived as the largest rectangle BOTH arms can work with the bench in scene, 200 x 80 mm, which is why its three targets landed 76 mm apart.} \\
\addlinespace[3pt]
{\scriptsize\ttfamily e775862} & \textbf{T0's three targets become three DIRECTIONS, re-derived in T0's own empty scene}\newline{\small The targets were 76 mm apart and the whole task was 0.219 m of commanded travel. They came from \code{task0.BAND}, which is the largest rectangle BOTH ARMS can work WITH THE BENCH IN THE SCENE -- 200 x 80 mm -- and T0 has no bench. The band was correct for the question it was derived for and wrong for this one, and nothing said so because the same scene was applied to every task.} \\
\addlinespace[3pt]
{\scriptsize\ttfamily 825408f} & \textbf{Resume point: the geometry is clean, mode 06 is recorded by the NEXT session}\newline{\small What changed about the state of the project, in one line: mode 06 is no longer BLOCKED, it is simply not RECORDED. Every earlier resume point named a blocker -- the scene, then discovery, then the transport -- and all three are settled and measured.} \\
\addlinespace[3pt]
{\scriptsize\ttfamily cca29fc} & \textbf{Archive the pre-fix mode-06 clips with a geometry note before re-recording}\newline{\small Four clips recorded against a scene that has since changed in four ways, from a run that gated every capture on timeout with the arm at home. They are evidence about the failure, not about the tasks, and the note says which.} \\
\addlinespace[3pt]
{\scriptsize\ttfamily 73d8f46} & \textbf{Mode 06 recorded, and SIX defects found by looking at the frames}\newline{\small All four tasks pass the sweep AND inspection. Mode 06 is citable. task angles EE travel L/R (m) grasp/release sweep T0 8 1.716 / 1.477 - / - OK, motion-gated T1 8 2.001 / 0.239 4 / 4 OK, placed 3 mm T2 8 0.711 / 0.604 - / - OK, motion-gated T3 8 0.249 / 0.361 2 / 2 OK, placed 0 mm Every clip is 7 angles plus the quad, captions burnt in (mode, task, scenario, expected, caveat, result), Xvfb :91-:98, capture gated on MOTION in all four -- not one gated on timeout.} \\
\addlinespace[3pt]
{\scriptsize\ttfamily 8d2e050} & \textbf{Mode 01 recorded, and T1's LAST cube exposed a lag defect that had to be fixed in the TASK, not the mode}\newline{\small Mode 01 master teleop: four tasks, eight angles each, all motion-gated. task EE travel L/R (m) grasp/release sweep T0 1.569 / 1.283 - / - OK T1 2.099 / 0.000 4 / 4 OK, placed 0 mm T2 0.711 / 0.604 - / - OK T3 0.909 / 1.016 2 / 2 OK, placed 0 mm THE DEFECT, and it is the reason to record one mode at a time. T1's first mode-01 pass grasped THREE of four cubes and still passed the sweep, because the placement check reads items[0] and cube\_0 was fine.} \\
\addlinespace[3pt]
{\scriptsize\ttfamily 5b2cb67} & \textbf{Mode 03 recorded: four tasks, all four grasps closing at 0.0 mm}\newline{\small task EE travel L/R (m) grasp/release sweep T0 1.541 / 1.294 - / - OK T1 2.165 / 0.515 4 / 4 OK, placed 16 mm T2 0.711 / 0.605 - / - OK T3 0.911 / 1.023 2 / 2 OK, placed 0 mm Every capture gated on MOTION, eight angles each, captions burnt in. THE LENGTHENED PICK HOLD HOLDS UP ACROSS MODES, which is what it was for.} \\
\addlinespace[3pt]
{\scriptsize\ttfamily addf3b1} & \textbf{Mode 02 VR teleop recorded, and the VR lag shows up exactly where it should}\newline{\small task EE travel L/R (m) grasp/release sweep T0 1.707 / 1.453 - / - OK T1 2.118 / 0.480 4 / 4 OK, placed 22 mm T2 0.718 / 0.575 - / - OK T3 0.904 / 0.998 2 / 2 OK, placed 0 mm THE ISOLATION RULE WAS ASSERTED, NOT ASSUMED, and it named the right thing: isolated: no unowned publisher on the follower input (the MAPPER is the only publisher; the runner drives it upstream on /vr/controller\_pose\_*) VR IS MEASURABLY THE LAGGIEST PATH, and it falls out of the numbers rather than being asserted: mode grasp at clip T1 placed c [...]} \\
\addlinespace[3pt]
{\scriptsize\ttfamily 1dcac33} & \textbf{Mode 04 recorded -- ALL FIVE MODES DONE, one identical task definition}\newline{\small 04\_vr\_shared task EE travel L/R (m) grasp/rel sweep T0 1.760 / 1.478 - / - OK T1 2.160 / 0.446 4 / 4 OK, placed 9 mm T2 0.710 / 0.604 - / - OK T3 0.909 / 1.024 2 / 2 OK, placed 0 mm THE WHOLE SET, 20 clips x 8 angles = 160 files, every capture gated on MOTION and not one on timeout: mode T0 travel L/R T1 g/r T1 placed worst closed 06\_full\_autonomy 1.716 / 1.477 4/4 14 mm 0.0000 m 01\_master\_teleop 1.569 / 1.283 4/4 0 mm 0.0000 m 03\_shared\_autonomy 1.541 / 1.294 4/4 16 mm 0.0000 m 02\_vr\_teleop 1.707 / 1.453 4/4 22 mm [...]} \\
\addlinespace[3pt]
{\scriptsize\ttfamily aea7379} & \textbf{Resume point: all five modes recorded, and the seven defects only LOOKING found}\newline{\small The recording is done. What the file now leads with is the part that matters for whoever picks this up: four of the seven defects passed the sweep, and every one of them was an arm moving with objects moving alongside it -- the shape that looks correct in every metric.} \\
\addlinespace[3pt]
{\scriptsize\ttfamily 29146f1} & \textbf{Freeze the verified mode-06 clips: archived copy, README, and a tag}\newline{\small The four mode-06 clips are the first in this project to pass BOTH gates -- the automatic sweep and frame-by-frame inspection -- and four of the seven defects found on the way there had passed the sweep, so a clip set that survives both is not cheap to reproduce. Copied to extras/archive/recordings/mode06\_GOOD\_20260811/, 32 mp4 plus every scene\_events.json and the sweep's progress file, 19 MB.} \\
\addlinespace[3pt]
{\scriptsize\ttfamily 061f77b} & \textbf{Real-robot readiness: a paper audit, and the safety path is weaker than the arm path}\newline{\small docs/system/14\_real\_robot\_readiness.md. No hardware was contacted. Every verdict is read off the code and the machine's current state, and where the code and docs/ENGINEERING_LOG.md disagree the code wins and the disagreement is named. THE HEADLINE: the ARM path is in better shape than the SAFETY path.} \\
\addlinespace[3pt]
{\scriptsize\ttfamily a81797d} & \textbf{Task 1: a real two-level table, per-cube slots, measured workspace markings -- and the front centre is NOT REACHABLE, which changes 2(b)}\newline{\small VERIFIED N=10 OVER THE FULL PATH AFTER EVERY MOVE: 6304 IK calls, TOTAL FAILURES 0, all four controls correct, every task against its own scene. (b) THE FRONT CENTRE IS NOT BUILDABLE, and this is the measurement rather than an opinion.} \\
\addlinespace[3pt]
{\scriptsize\ttfamily ac0de7d} & \textbf{Note in the frozen archive that the scene changed after those clips}\newline{\small The clips are untouched and still valid for what they show. The README now says what moved afterwards -- the two-level table, the planes 15 mm outboard, the per-cube slots and the workspace markings -- so a reader does not assume the current scene matches, and marks T1 in that set as superseded geometry.} \\
\addlinespace[3pt]
{\scriptsize\ttfamily ea0c961} & \textbf{Resume point: items 0-2 done, 3-5 not started, and the region item 3 needs}\newline{\small The file now opens with the six-part brief as a state table and the measured reachable region item 3 samples from, with the two limits that travel with it -- the survey box truncates it and cells were tested at the grasp pose only. It states 2(b) as a finding rather than an omission: zero reachable cells anywhere with |x| <= 0.10, nearest reachable x is 0.30, so planes centred in the front centre cannot exist on this rig.} \\
\addlinespace[3pt]
{\scriptsize\ttfamily be28e16} & \textbf{Fix the two silent real-robot failures, and pin both against regression}\newline{\small Both had the same shape: a safety mechanism that runs, reports success, and is about something other than the real arm. Neither showed up in any check, because every check was satisfied -- the clearance number was a number and the halt logged a line. (a) THE CLEARANCE MEASURED THE WRONG ARM AND SAID NOTHING. measure\_clearance() looked up \code{f"\{arm\}\_\{link\}"} -- unprefixed, the SIM tree -- while the real arm publishes under \code{real\_*} (both real launches run robot\_state\_publisher with frame\_prefix="real\_").} \\
\addlinespace[3pt]
{\scriptsize\ttfamily f3e8a5e} & \textbf{The unprefixed-name class, swept: five more services and one live bug}\newline{\small This class has now appeared six times and every instance was found separately -- reactivate\_gripper as a GPIO, the Robotiq COM port, tcp/twist.* and reset\_fault/* as C++ literals, reactivate\_gripper again in the driver, and last commit's emergency halt. Finding them one at a time IS the defect, so this is the sweep rather than another instance. THE RULE APPLIED, because "prefix everything" is not it.} \\
\addlinespace[3pt]
{\scriptsize\ttfamily f1d3122} & \textbf{Make the test suite unavoidable: the session runner refuses to RECORD while it fails}\newline{\small Two tests broke last session and went unnoticed for a whole session because the suite was not run. Of the three ways to make that impossible, I chose the third, and the reasoning is the point:} \\
\addlinespace[3pt]
{\scriptsize\ttfamily 0b843dd} & \textbf{T1 Stage 2: both arms at once, positions from the surveyed region -- and the region had to be re-surveyed to be usable}\newline{\small A THIRD KIND OF BIMANUAL, and the write-up must not merge it with the other two. The project already separates: T2 PHYSICAL COUPLING -- one rigid body at two grips 500 mm apart. Neither arm's pose is free given the other's: a height difference at one grip IS a tilt at the other. Remove an arm and the task is IMPOSSIBLE. Tilt and separation are JOINT metrics no single arm can produce. T3 BIMANUAL BY ROLE -- two objects, two places, two jobs. Neither arm constrains the other.} \\
\addlinespace[3pt]
{\scriptsize\ttfamily e4bce0c} & \textbf{Item 4: the sequence, run end to end -- and THE PIPELINE PRODUCES NO DATASET}\newline{\small The question was "so I know the pipeline produces a complete dataset before a participant sits down". Run end to end in sim, the answer is NO, and that is the deliverable. CAPTURED: 0 sample rows across 0 files \code{run\_abc.py} is the ONLY entry point the MSc tasks have -- run\_experiment.sh --taskset msc dispatches straight to it -- and it contains ZERO references to TrialLogger.} \\
\addlinespace[3pt]
{\scriptsize\ttfamily 8cc305b} & \textbf{Resume point: the data path is missing, and item 5 is not started}\newline{\small The file now opens with the two things that decide what happens next: the MSc set has a CLIP path and no DATA path -- measured, 0 sample rows across 0 files -- and the demonstration mode is not started, with a note that its perception half rests on a detection rate that is still unmeasured. Everything done this session is a one-line summary underneath, including the two silent real-robot failures and the five further unprefixed services the sweep found.} \\
\addlinespace[3pt]
{\scriptsize\ttfamily abc00ae} & \textbf{The MSc tasks now HAVE a data path, and every mode moves}\newline{\small Item 4's finding was that the whole study's output was missing: run\_abc.py was the only entry point the MSc tasks had and it never constructed a TrialLogger. Measured then: 0 sample rows across 0 files. Measured now, all 19 cells of the sequence, every task in every mode it runs in: 19 of 19 cells OK, 0 FAIL 1520 sample rows across 19 files 42 of 66 sample columns populated THE LOGGER IS OPENED BY run\_abc ITSELF, not by a wrapper.} \\
\addlinespace[3pt]
{\scriptsize\ttfamily 9fa4352} & \textbf{Choreography verified; unmask the failure the teardown was eating}\newline{\small THE TEARDOWN WAS DESTROYING THE EVIDENCE. run\_one returns grabs=None on every failure path, and stop\_grabs(None) raised AttributeError -- before the sweep logged WHY the run failed. So a real, nameable failure surfaced as a stack trace in the cleanup code, pointing at the wrong file. Fixed at both ends: stop\_grabs tolerates an empty map, and the sweep prints the reason at the moment it decides, not after teardown.} \\
\addlinespace[3pt]
{\scriptsize\ttfamily 6d2c3cd} & \textbf{Mock RGB-D at the camera boundary; T1 stage 2 recorded}\newline{\small MOCK THE BOUNDARY, THEN SAY WHERE IT IS. Same pattern as mock\_real.launch.py and the desktop Quest mock: stand a plausible producer at the edge so everything inside it runs, and be explicit about what that cannot show. Publishes on the topics the real driver owns -- /<arm>\_camera/color/image\_raw (rgb8), /color/camera\_info, /depth/image\_raw (16UC1 MILLIMETRES). Depth is mm on purpose: vlm\_object\_locator carries a mm-or-metres branch and the mm side has never run.} \\
\addlinespace[3pt]
{\scriptsize\ttfamily af5a297} & \textbf{Status table, read from disk -- and two false positives in it}\newline{\small Every MSc task x every mode, with four columns counted SEPARATELY because they fail separately: built, verified N=10 over the full path, clip on disk, trial rows on disk. "It works" has meant each of these at different times in this repo. task B master vr\_tel shared vr\_shr full\_auto T0 target reaching 9 CD CD CD CD CD T1 pick and place 10 CD CD CD CD CD T1 s2 both arms 10 .D .D .. .D CD T2 coordinated carry 9 CD C. C. C. CD T3 circuit + meter 10 CD C. CD C.} \\
\addlinespace[3pt]
\end{longtable}

\section*{2026-08-12 \textnormal{\small\textcolor{srlgrey}{--- 30 commits}}}
\addcontentsline{toc}{section}{2026-08-12}
\begin{longtable}{@{}p{0.055\textwidth}p{0.885\textwidth}@{}}
\endfirsthead\endhead
{\scriptsize\ttfamily 8212e58} & \textbf{Fix the stale spec count the suite caught; resume point}\newline{\small test\_only\_the\_two\_CURRENT\_task\_sets\_are\_offered asserted 20 MSc specs (4 tasks x 5 modes) against 25. The CODE was right: T1 has two STAGES and they are separate dispatcher keys with separate session cells, because stage 2 is the simultaneity condition and a participant must do stage 1 first or they learn the task and the simultaneity at once. A stale constant in the test written to catch drift is precisely the drift it exists to catch -- which is an argument for the test, not against it.} \\
\addlinespace[3pt]
{\scriptsize\ttfamily 159637f} & \textbf{T2/T3 VR rows: the path is fine, the PLAN does not ask. Table now says so}\newline{\small CAUSE, MEASURED BEFORE CHANGING ANYTHING. The missing T2/T3 VR data cells are not the VR data path failing again. msc\_session.PLAN runs T2 in the two ANCHORS ONLY and T3 in the anchors plus shared autonomy -- the deliberately half-crossed design the module already documents and costs out. The logged cells match the plan exactly: planned cells 19 cells with rows 19 PLANNED BUT MISSING DATA : NONE logged though NOT planned: NONE And the path itself was re-proven rather than argued.} \\
\addlinespace[3pt]
{\scriptsize\ttfamily cd3fac5} & \textbf{T1 stage 2 filmed in all five modes -- no real gaps remain}\newline{\small clip dirs 25 logged cells 19 of 19 planned REAL GAPS -- missing clip : NONE REAL GAPS -- planned, no data : NONE Inspected, not just counted: the 02\_vr\_teleop clip shows the right caption, two yellow cubes on the left arm's side and two cyan on the right, both arms moving at once and neither entering the other's cells.} \\
\addlinespace[3pt]
{\scriptsize\ttfamily 13eb160} & \textbf{SCAN POSE: the wrist cameras cannot see the work surface from home}\newline{\small THIS IS THE FIRST THING THAT WOULD GO WRONG ON CAMERA DAY, and it is now answered before anyone stands in a lab. Measured in sim: at home the bench projects to pixel (954, 539) and the table to (1115, 765) for the right camera -- IN FRONT of it and outside a 640-wide image. The left camera is parked looking up and back.} \\
\addlinespace[3pt]
{\scriptsize\ttfamily 812cafe} & \textbf{Choreography FILMED -- three routines, eight angles, caveat in every caption}\newline{\small Built and verified last session (2420 IK calls, 0 failures); now recorded. 06\_full\_autonomy/d1/D1\_wave OK 8 files, quad 06\_full\_autonomy/d2/D2\_mirror OK 8 files, quad 06\_full\_autonomy/d3/D3\_sweep OK 8 files, quad -> recordings/demo\_20260812/ THE CAVEAT IS IN THE PICTURE, not only in the module. Every frame carries: "A DEMONSTRATION IS NOT EVIDENCE. This routine produces no trial data, no condition, no metric and no comparison, and nothing in it may be cited as a result.} \\
\addlinespace[3pt]
{\scriptsize\ttfamily 4d97734} & \textbf{CONFIRMED: the band is at the body, and the BENCH is why. Measured, not argued}\newline{\small CAUSE CONFIRMED AND EXTENDED. The reported cause is right and understates it. The surveyed band is y 0.05..0.20; the bench surface spans y 0.245..0.630, so it overlaps the band by EXACTLY 0.000 m. The bench is entirely unreachable. And all four T1 cubes are over NO surface at all -- not over the bench (they sit 15-75 mm in front of its near edge) and 0.170 m above the table. Nothing ever checked object-against-surface, so this was never reported.} \\
\addlinespace[3pt]
{\scriptsize\ttfamily b8d5b89} & \textbf{Mount measured (it is NOT the fix); dance rebuilt with beat and phrasing}\newline{\small === 1. THE NUMBER I OWED YOU. (a) IS NOT THE FIX === Furniture transformed per arm along with the target -- the honest version of the mount experiment -- probing 20 mm above the table top: baseline L 0.050 R 0.050 tilt -30 deg L 0.125 R 0.150 tilt -45 deg L 0.200 R 0.225 fwd 0.15 + tilt -30 L 0.125 R 0.150 fwd 0.30 L 0.050 R 0.050 The best candidate buys 0.150 m and STILL leaves the work at the waist.} \\
\addlinespace[3pt]
{\scriptsize\ttfamily 0f790af} & \textbf{ONE TABLE: bench deleted, work moved out front -- NOT YET RECORDABLE}\newline{\small APPLIED * the bench is gone from every MSc scene. It cost 0.250 m of forward reach (no furniture 0.425, table only 0.425, bench only 0.175) and its surface overlapped the reachable band by EXACTLY ZERO, so nothing could rest on it. The A/B/C set keeps it -- those clips were verified against it.} \\
\addlinespace[3pt]
{\scriptsize\ttfamily dce8a7d} & \textbf{The dance timing DOES survive to the arm -- measured, three stages}\newline{\small CLAIM CHECKED BEFORE ACTING ON IT, and it does not hold. The recorder does not teleport through waypoints: routine recorded clip flow 18.1 s 49.8 s pulse 17.4 s 50.9 s play 25.0 s 68.1 s The clips are LONGER than the routines, not shorter -- real time plus the approach from home and the settle. run\_abc publishes one waypoint per spin(0.05), which is wall-clock 20 Hz, so duration is honoured by construction; there is no per-waypoint duration for the recorder to ignore because the runner IS the clock.} \\
\addlinespace[3pt]
{\scriptsize\ttfamily f84c22b} & \textbf{LIFT MEASURED: it is not the limitation. The LEFT ARM cannot grasp at all}\newline{\small Asked at every real T1 grasp point, N=10 per 5 mm step, lift = the last height that passed N/N (a MARGINAL height does not count): PINNED, scene left --, --, --, -- NONE REACHABLE PINNED, scene right 300, 300, 300, 300 >= 300 mm PINNED, no furniture left --, 300, 300, 300 cube\_0 unreachable PINNED, no furniture right 300, 300, 300, 300 >= 300 mm FREE wrist, scene left --, --, --, -- NONE REACHABLE FREE wrist, scene right 300, 300, 300, 300 >= 300 mm 300 mm is the SWEEP CAP, not a limit -- the right arm never ran ou [...]} \\
\addlinespace[3pt]
{\scriptsize\ttfamily ae7c140} & \textbf{Placement answered from scene\_events, not frames: 20/20 on their OWN plane}\newline{\small TOTAL own 20 WRONG 0 off 0 never released 0 Four cubes x five modes, every one on the plane of its own colour. The question no longer depends on frame sampling: a cube is 40 mm, a plane is 140 x 100 mm and the capture runs at 5-11 fps, so the moment of release can fall between frames entirely -- colour-pixel counting is the wrong instrument for a question about POSITION. THREE VERDICTS, NOT TWO.} \\
\addlinespace[3pt]
{\scriptsize\ttfamily 27a97ed} & \textbf{Could NOT reproduce the Section 1 premise -- reporting before building on it}\newline{\small Everything in items 1b, 1c and 3 rests on "TRAC-IK honours OrientationConstraint and top-down is 0/128 with the wearer". Item 3 would write that into the thesis as a design finding, so it was probed first. It does not reproduce here.} \\
\addlinespace[3pt]
{\scriptsize\ttfamily b861660} & \textbf{FK CHECK: the constraint IS honoured. Tolerance artefact confirmed}\newline{\small pinned and top-down are 169.7 deg apart (tool axis 120.8 deg) request achieved exact pinned, no constraint 0.00 deg from pinned pose PINNED + constraint TOP-DOWN 0.05 REFUSED 4/4 exact top-down, no constraint 0.00 deg from top-down pose TOP-DOWN + constraint TOP-DOWN 0.05 0.00 deg from top-down ROW 2 IS THE WHOLE ANSWER, and it is why the test asks for a pose and a constraint that DISAGREE.} \\
\addlinespace[3pt]
{\scriptsize\ttfamily 91e6dd5} & \textbf{TOP-DOWN TESTED AS A REAL APPROACH: it is complementary, not better}\newline{\small N=10, full path, wearer AND furniture, constrained yaw-free IK at 0.05 rad -- the mechanism modes 3/4/6 would actually use: task pinned near-side top-down constrained T1 0/4 reachable 0/4 reachable T2 2/2, lift >= 300 mm 0/2 reachable T3 0/2 reachable 1/2, lift >= 300 mm THE TWO APPROACHES ARE COMPLEMENTARY, NOT ORDERED. Each rescues a task the other cannot do.} \\
\addlinespace[3pt]
{\scriptsize\ttfamily fe6db97} & \textbf{T1 layout: the PEDESTALS are fatal in every option. Removed}\newline{\small Three options, N=10, wearer and furniture, before choosing anything: option cubes planes left arm, table only 2/4 1/2 left arm, stands ON 0/4 0/2 RIGHT arm, table only 4/4 2/2 <- the only clean pass right arm, stands ON 0/4 0/2 split, table only 3/4 2/2 THE PEDESTALS ARE FATAL EVERYWHERE -- 0/4 in every layout that has them.} \\
\addlinespace[3pt]
{\scriptsize\ttfamily f2da765} & \textbf{The workspace marking is STALE -- the mismatch is 170 mm, not 15 mm}\newline{\small Measured against WORKSPACE['left'] = x 0.300..0.700, y 0.050..0.200: object x outside by y outside by cube\_0 60 mm 10 mm cube\_1,2,3 0 10 mm plane\_0 70 mm 170 mm plane\_1 0 170 mm WIDENING x TO 0.365 WOULD NOT FIX IT, so it is not applied. The dominant error is in y and it is 170 mm on both planes. WORKSPACE was surveyed from the work plane WITH THE BENCH IN SCENE, where the reachable band really was y 0.05..0.20.} \\
\addlinespace[3pt]
{\scriptsize\ttfamily e18717d} & \textbf{Two design properties written up: mode-independence and the 40 mm floor}\newline{\small MODE-INDEPENDENCE IS STRUCTURAL, NOT EMPIRICAL, and that is the stronger claim. In run\_abc.py the commanded trajectory is \code{wp = spec["build"]()} -- built from the TASK spec with NO mode argument. \code{mode} selects only the topic the poses go out on, the isolation check and the logging labels. Every mode therefore commands identical waypoints to the same follower under the same limits, so 100\% / 100\% / zero variance is the NECESSARY consequence rather than a surprising measurement.} \\
\addlinespace[3pt]
{\scriptsize\ttfamily f82b814} & \textbf{MEASURED: the wrist DOES point up at home. +30.8 deg left, +22.1 deg right}\newline{\small arm tool axis elevation angle to table normal left (-0.153, +0.846, +0.511) +30.8 deg 120.8 deg right (+0.259, +0.890, +0.376) +22.1 deg 112.1 deg left fingertip sits +40 mm ABOVE the EE origin right fingertip sits +49 mm ABOVE the EE origin SO THE OBSERVATION IS CORRECT AND IT IS NOT A FRAME BUG. The gripper approach axis points UP AND FORWARD at the home pose on both arms, 112-121 degrees away from the table normal, and the fingertips sit ABOVE the wrist rather than below it.} \\
\addlinespace[3pt]
{\scriptsize\ttfamily a3af107} & \textbf{Pad measurement is INVALID -- FK returned the home pose. Not reporting it}\newline{\small The per-grasp pad measurement ran to completion and its output must not be used. Every row reports the SAME EE and pad height regardless of the grasp requested: t1 left (0.320, 0.190, 1.120) EE z 1.146 PAD z 1.186 t2 left (0.250, 0.350, 1.320) EE z 1.146 PAD z 1.186 t3 left (0.357, 0.095, 1.065) EE z 1.146 PAD z 1.186 Three grasps 255 mm apart in z returning one number to the millimetre.} \\
\addlinespace[3pt]
{\scriptsize\ttfamily 8e360bb} & \textbf{PADS MEASURED PROPERLY: not in the table -- 40 mm PAST the object, upward}\newline{\small Instrument first, because the last one lied. /compute\_fk returned the home pose whatever solution it was given, so it is not used. The EE->pad offset is CONSTANT in the end-effector frame, so pad\_world = p\_grasp + R(q) @ offset is exact and needs no solver. Validated by recomputing the home pad position from the offset: error 0.00e+00 m (left) and 5.55e-17 m (right). Negative control: a grasp 300 mm below the table is flagged at -260.5 mm, so the check CAN fail.} \\
\addlinespace[3pt]
{\scriptsize\ttfamily bf9cdc0} & \textbf{Part 2b: the work surface height has ONE owner, and 20 mm is now NAMED}\newline{\small FIXED AT SOURCE, not with a check bolted on top. srl\_experiments/work\_surface.py holds the height; clip\_scene READS it rather than defining it. The fallback literal survives only inside an except branch that WARNS, because a silent fallback to a hardcoded height is precisely how a measured surface can exist and reach nothing. THE SPLIT THAT MATTERS is DECLARED versus MEASURED. Declared is what the sim builds and what every task coordinate was verified against; it is a constant and should be.} \\
\addlinespace[3pt]
{\scriptsize\ttfamily 04c8f30} & \textbf{PART 3: twelve new faults injected. THREE FAIL SILENTLY}\newline{\small fault\_injector.py handles 14 and predates the camera work, the work-surface owner and the object-pose gap. These twelve had never been injected. fault what named recovers object\_rotated\_30deg SILENT NO no object\_missing SILENT NO no wearer\_moves SILENT NO no camera\_garbage PARTIAL yes unknown gui\_killed PARTIAL yes relaunch, session resumable table\_20mm\_high DETECTED yes yes table\_20mm\_low DETECTED yes yes table\_never\_measured DETECTED yes n/a object\_moved\_after\_scan DETECTED yes re-scan (srv\_resweep) camera\_absent [...]} \\
\addlinespace[3pt]
{\scriptsize\ttfamily 0e78162} & \textbf{A1 first half: attachment tests ORIENTATION, so a rotated object now FAILS}\newline{\small WHAT WAS SILENT. Attachment tested DISTANCE alone -- \code{near = d <= GRASP\_NEAR\_M} -- so an object at an angle was grasped as though square and the sim reported a clean grasp. On a 40 mm cube a 30 deg error is about 20 mm of misalignment across the face; in reality the fingers catch a corner or shove the cube away. Nothing anywhere compared the gripper's closing axis to the object's, so simulation could not show it and it would have appeared on real hardware as an inexplicable miss. NOW MEASURED.} \\
\addlinespace[3pt]
{\scriptsize\ttfamily 0a2c93c} & \textbf{Part F: the accuracy table is now a REPRODUCIBLE SCRIPT, not a quoted figure}\newline{\small Flagged three times that the 100\% / 0.000 mm / zero-variance numbers lived only in session reports and could not be re-derived. scripts/accuracy\_table.py derives them from scene\_events.json, which every clip already writes, so the table is a consequence of committed data rather than a recollection. No new trials are needed to reproduce it. AND THE NUMBERS DO NOT MATCH WHAT WAS QUOTED.} \\
\addlinespace[3pt]
{\scriptsize\ttfamily 15af670} & \textbf{Split docs/ENGINEERING_LOG.md: 98k tokens -> 11 kB. Nothing deleted, verified line by line}\newline{\small docs/ENGINEERING_LOG.md was 4789 lines and 98.2k tokens, a TENTH of a session's context window consumed before any work started, every session, whether or not any of it was relevant. It is now 209 lines and about 3k tokens. docs/system/findings.md 2935 lines + the old header as an appendix docs/system/hardware.md 689 lines docs/system/wsl.md 375 lines docs/system/architecture.md 632 lines NOTHING WAS DELETED AND THAT IS CHECKED, NOT ASSERTED.} \\
\addlinespace[3pt]
{\scriptsize\ttfamily 5f50355} & \textbf{PART C: three dance buttons dispatched to the WRONG taskset and exited 1}\newline{\small FOUND BY RUNNING, which is the only way this class is ever found. [1] dispatch d1 FAIL exit 1 unknown msc task key 'D1' [1] dispatch d2 FAIL exit 1 [1] dispatch d3 FAIL exit 1 I added d1|d2|d3 to the MSc case arm, which hardcodes \code{--taskset msc}, so \code{run\_experiment.sh d1} exited 1. THE GUI HAPPENED TO SURVIVE IT: its spec appends \code{--taskset demo} and argparse takes the last one, so the button worked while the command it claims to run did not.} \\
\addlinespace[3pt]
{\scriptsize\ttfamily 4977a43} & \textbf{Part D blocked: the two gates the brief records as closed are OPEN}\newline{\small Checked in the repo rather than assumed: T1\_CUBES still at POSITIVE x -- T1 still runs on the left arm at 0/4 WORKSPACE still y (0.05, 0.20) -- the 170 mm stale value pad offset no 0.098 / tool\_axis anywhere in the task or action code Recording against that films the left arm failing to reach four cubes, inside a marking that excludes both targets, with the pads landing 39.5 mm past each object. That is a fourth superseded set and it costs an hour to produce. NOT RECORDING.} \\
\addlinespace[3pt]
{\scriptsize\ttfamily 271eb34} & \textbf{Sweep names what it is about to film, and refuses the wrong taskset}\newline{\small CORRECTION TO THE PREMISE, checked rather than accepted: the switch already exists and already defaults to the MSc set --} \\
\addlinespace[3pt]
{\scriptsize\ttfamily 396f418} & \textbf{T1 mirrored to the right arm: 49 failures -> 1. Not yet zero, so NOT recording}\newline{\small MEASURED AT EVERY STEP, N=10 full path, wearer and furniture in scene: starting point (left arm, pedestals) 49 failures mirrored to the RIGHT arm 11 cube row moved into the surveyed band 2 planes pulled 20 mm nearer 1 T1 now drives the RIGHT arm: 4/4 cubes and 2/2 planes at the grasp points, 146 waypoints, 8 grip changes. T1\_ARM is honoured by the path and the grip schedule instead of "left" being hardcoded in three places, so the side is one edit rather than four.} \\
\addlinespace[3pt]
{\scriptsize\ttfamily c4c4040} & \textbf{THE LAYOUT GATE IS CLOSED: 0 failures across all five tasks, N=10 full path}\newline{\small t0 50 distinct waypoints, 0 failures t1 80 distinct waypoints, 0 failures t1s2 50 distinct waypoints, 0 failures t2 22 distinct waypoints, 0 failures t3 16 distinct waypoints, 0 failures Wearer and furniture in the scene, N=10 per waypoint over the densified path.} \\
\addlinespace[3pt]
\end{longtable}

\section*{2026-08-13 \textnormal{\small\textcolor{srlgrey}{--- 37 commits}}}
\addcontentsline{toc}{section}{2026-08-13}
\begin{longtable}{@{}p{0.055\textwidth}p{0.885\textwidth}@{}}
\endfirsthead\endhead
{\scriptsize\ttfamily 9db677f} & \textbf{Archive note: everything before the T1 mirror is superseded}\newline{\small Records what changed and why the old clips should not be used: T1 moved to the right arm from 0/4 on the left, the bench and the pedestals are gone, and the cube row moved into the surveyed band. The layout now verifies at zero IK failures, N=10 full path, wearer and furniture in scene. mode06\_GOOD is untouched.} \\
\addlinespace[3pt]
{\scriptsize\ttfamily 88e5fe6} & \textbf{Information card plays FIRST, instead of an overlay across the footage}\newline{\small The caption used to sit on top of the arms for the whole clip, competing with the thing it describes. A viewer reads it while trying to watch the motion. A card that plays first is read once and then gets out of the way. Six seconds of full-frame card, then the footage, on EVERY angle rather than only the front one. Card states the task, how it is being driven, what happens and what to watch for.} \\
\addlinespace[3pt]
{\scriptsize\ttfamily df7d3d5} & \textbf{T1's scene still watched the LEFT arm after the task moved right}\newline{\small 'NO GRASP RECORDED at all' on a run where the arm plainly picked cubes up and placed one 112 mm from target. The task moved to the right arm; clip\_scene still had arm='left' hardcoded on every T1 cube, so the right gripper closed on the object while the scene watched the left one. Now follows MCT.T1\_ARM, so the side is one edit rather than two files that can disagree. Scenario renamed S1\_left\_arm -> S1\_right\_arm, which was the visible tell and which I left stale when I mirrored the task.} \\
\addlinespace[3pt]
{\scriptsize\ttfamily fb7e425} & \textbf{Orientation check REPORTS instead of gating; it was refusing every grasp}\newline{\small I added it last session to make a rotated object visible in sim. It made attachment conditional on the gripper being within 20 deg of square, and the pinned wrist sits well off square, so a live sweep produced four clips reading 'NO GRASP RECORDED at all' on runs where the arm visibly picked cubes up. The check now measures and warns without deciding whether the object attaches. max\_yaw\_err\_deg still lands in scene\_events.json and the warning still fires, which is what made the failure visible.} \\
\addlinespace[3pt]
{\scriptsize\ttfamily d2a0ece} & \textbf{Camera was framing where the work USED to be}\newline{\small T1 moved to the right arm, so its cubes sit at x -0.32..-0.50. The capture camera still looked at x=0, which put the whole task at the right edge of frame and partly outside it, with the empty left-arm workspace marking occupying the middle of the picture. A viewer sees a marked box with nothing in it and the work happening off the edge. Focal point x 0 -> -0.10, the true centre of every task's span (-0.54..0.34), and distance 2.25 -> 2.55 so the whole span fits. Found by extracting a frame and looking at it.} \\
\addlinespace[3pt]
{\scriptsize\ttfamily 81eae18} & \textbf{Part E: recording is clean, the CAMERA is aimed at the wrong side}\newline{\small Mode 01 recorded 5 of 5 with 0 failures, twice. T1 shows 4 grasps with the pads on the cube at 0.000 m, and the card plays first and reads well. Both bugs from the previous sweep stayed fixed. What is wrong is the framing, and I found it by extracting frames and looking rather than by any check: the work sits at the right edge of frame and partly outside it, while the EMPTY left-arm workspace marking holds the middle of the picture. A viewer sees a marked box with nothing in it.} \\
\addlinespace[3pt]
{\scriptsize\ttfamily 8deb5ca} & \textbf{Framing: eliminated two wrong controls by experiment, found the third}\newline{\small The front view's aim comes from NEITHER of the things I edited. Established by running, not by reading: verification\_capture.rviz aims the standalone RViz only. Edited twice, the frame never moved. \_FOCUS in record\_rviz.py aims the other angles. Set to 0.15 and to -0.10; the front view was IDENTICAL both times. the FRONT view FOLLOWS THE TASK record\_rviz.py line 121 says so, and the config cache holds front\_t1.rviz next to front.rviz. That is the one to fix.} \\
\addlinespace[3pt]
{\scriptsize\ttfamily 71a39bc} & \textbf{Framing fixed: TASK\_FOCUS aims the front view, and T1's entry was stale}\newline{\small TASK\_FOCUS in record\_rviz.py aims the FRONT view per task. T1's entry was (0.34, 0.22, 1.19) which was the OLD left-arm cube row. T1 moved to the right arm at x -0.32..-0.50, so the task sat at the right edge of frame while the empty left-arm marking held the middle of the picture. Now (-0.30, 0.22, 1.19) at distance 1.45. Checked by looking: the wearer is fully in shot, the right arm reaches out over the work, the blue and green cubes sit on their matching mats, and the cyan right-arm marking encloses the work.} \\
\addlinespace[3pt]
{\scriptsize\ttfamily 0777214} & \textbf{Mode 01 recorded against the fixed layout and framing}\newline{\small 5 of 5, 0 failures. Cards play first, the work is in frame with its marking around it, and T1 shows 4 grasps with the pads on the cube.} \\
\addlinespace[3pt]
{\scriptsize\ttfamily 4f3386c} & \textbf{Sweep: 19 of 25 cells recorded, and two patterns worth more than the count}\newline{\small 01\_master\_teleop 5/5 T1 4/4 T1S2 4/4 T3 2/2 02\_vr\_teleop 2/5 T1 0/4 T1S2 0/4 T3 0/2 03\_shared\_autonomy 4/5 T1 0/4 T1S2 4/4 T3 2/2 04\_vr\_shared 5/5 T1 4/4 T1S2 4/4 T3 2/2 06\_full\_autonomy 3/5 T1 0/4 T1S2 4/4 T3 2/2 dance 1/3 02\_vr\_teleop loses EVERY grasp while its arms travel 1.0 to 2.6 m, so the pose command lands and the gripper command does not. It is not that VR is broken: 04\_vr\_shared drives the same controllers and gets everything. The fault is specific to the 02 path.} \\
\addlinespace[3pt]
{\scriptsize\ttfamily dfdc1e4} & \textbf{TASK\_SPEC is the source of truth, and the pre-recording pass found nine gaps}\newline{\small docs/TASK\_SPEC.md holds the brief verbatim as acceptance criteria, and scripts/audit\_task\_spec.py asks every MUST BE TRUE item of the CODE before a single clip is filmed. 12 self-test probes, each breaking one input on purpose, because a check that cannot fail is not a check. 40 PRESENT, 0 MISSING, 1 BLOCKED. Nine items were missing when it started:} \\
\addlinespace[3pt]
{\scriptsize\ttfamily 5d4d133} & \textbf{Archive the three superseded clip sets with a note naming all nine changes}\newline{\small recordings/verification, verification\_20260813 and verification\_20260811\_t1rework move to extras/archive/recordings/superseded\_20260813\_prespec/. mode06\_GOOD is untouched. GEOMETRY\_NOTE.md names what changed under them, item by item: the re-surveyed marking, cells instead of a bounding box, one marking for a one-arm task, the presentation pose, T2's carry series, T3's card, the wrong arm in T1's expect text, and the leaked loop variable that measured left-arm yaw against the right gripper.} \\
\addlinespace[3pt]
{\scriptsize\ttfamily 9a465f9} & \textbf{Found what makes T1 flaky: ee\_for() uses the LEFT arm's pad offset}\newline{\small A single T1 clip under 01\_master\_teleop, the mode that used to score 4/4, records NO GRASP AT ALL with the pads stopping 31.4 mm from three of four cubes against a 30 mm gate, while the arm travels 2.72 m. THREE CUBES REPORTING THE IDENTICAL 0.0314 IS THE CLUE. They sit in a row 60 mm apart, so an identical closest approach to all three is the perpendicular distance from a path running parallel to the row: a CONSTANT offset, not the accumulating lag the earlier note described.} \\
\addlinespace[3pt]
{\scriptsize\ttfamily 1db80dc} & \textbf{One pad offset per arm, shared by the task layer and the scene}\newline{\small PAD\_OFFSET was the LEFT arm's wrist-to-pad vector and was used for both arms. Measured off TF at home, which IS the pinned anchor and the only orientation a task path ever holds: left (-0.0171, +0.0945, +0.0572) identical to the old constant right (+0.0289, +0.0995, +0.0421) 48.3 mm away from it ee\_for(obj, arm) now takes the arm and defaults to left, so the A/B/C set does not move. Every MSc task passes its own. AND clip\_scene.\_pad\_offset() reads the SAME table instead of sampling live TF.} \\
\addlinespace[3pt]
{\scriptsize\ttfamily d387b23} & \textbf{MOVE\_TO reaches the resolver, and T3's four points go on the box}\newline{\small TWO GAPS, and both were "the piece exists and nothing connects to it". 1. THE MOVE\_TO VERB. voice\_intent has parsed "move to the front centre" since it was written and NAMED\_PLACES said which places are reachable -- and the executive fell through to "I don't know how to 'goto'", so the system refused a command it was designed to accept.} \\
\addlinespace[3pt]
{\scriptsize\ttfamily cbc742e} & \textbf{Part 1: the pose in the screenshots is the PRESENTATION pose, and it is a T}\newline{\small IDENTIFIED FIRST, with numbers, because getting it wrong in either direction is expensive. scripts/describe\_pose.py reports both poses through /compute\_fk: HOME wrist axis +30.8 / +22.1 deg (UP) hand span 1.459 m PRESENTATION wrist axis +6.5 / +7.5 deg (LEVEL) hand span 1.633 m The discriminator is the wrist.} \\
\addlinespace[3pt]
{\scriptsize\ttfamily 184633c} & \textbf{The clearance floor, not the search, is what forbids hands at torso width}\newline{\small CORRECTING MY OWN MEASUREMENT. measure\_ready\_pose\_envelope.py reported both arms reaching |x| = 0.14 with a level wrist and I reported the brief as achievable. That sweep asked MoveIt's collision-aware IK, and the WEARER PAIRS ARE SRDF-EXCLUDED -- the trap mount\_guard\_node exists for. Measured geometrically instead, with no exclusions: hand |x| = 0.18 clearance -0.016 m the arm is INSIDE the torso 0.20 -0.004 0.26 +0.002 .. +0.041 0.30 +0.019 .. +0.045 0.45 +0.151 ..} \\
\addlinespace[3pt]
{\scriptsize\ttfamily cdcb99f} & \textbf{Looked at the render: the pose is still splayed, and the cost said it was fine}\newline{\small RENDERED IT AND LOOKED, which is the whole point of showing a pose before recording with it. scripts/render\_pose.py stages a pose and grabs a still per view, reusing record\_rviz's displays so the framing matches the clip. WHAT THE PIXELS SHOW, against a cost that scored it as satisfied:} \\
\addlinespace[3pt]
{\scriptsize\ttfamily 586dcba} & \textbf{The staging move now pauses the follower, so the arm actually goes where told}\newline{\small THE CONTROL CAUGHT IT. A pose search reported "home tool axis +22.1 deg (TF -0.0)" and refused to save: FK computed home while TF showed the arm still standing in the previous staged pose. The arm had never been homed -- \code{stage\_presentation\_pose.py --home} reported success and moved nothing. WHY. ik\_follower\_node streams position commands to the same arm controller this script publishes a trajectory to, and once the follower has a target it holds the arm there.} \\
\addlinespace[3pt]
{\scriptsize\ttfamily 8e220e6} & \textbf{A symmetric presentation pose, rendered and looked at. It is still not right.}\newline{\small All controls pass now that the arm is genuinely homed, and the pose is mirrored by construction: both hands at (+/-0.550, 0.360, 1.100), wrists level to 0.0 deg, wearer clearance 0.1610 m against a 0.150 floor, hand travel 0.19 and 0.24 m from home. WHAT THE RENDER SHOWS, and this is the part the numbers cannot say. It is no longer a T-pose and the two arms finally mirror each other, which is the fault the last render exposed. But:} \\
\addlinespace[3pt]
{\scriptsize\ttfamily d34118f} & \textbf{The dance had no scene node at all, and its envelope was a tenth of the room}\newline{\small TWO FAULTS, and the first explains why the dance clips were thin. 1. clip\_scene REFUSED the dance by its own argument list. "d1" was not in --task's \code{choices}, so every dance clip was filmed with the scene node failing to start: no furniture, no markers, and no scene\_events.json to say how far the arms travelled. The sweep launched it, argparse rejected it, and nothing downstream asked.} \\
\addlinespace[3pt]
{\scriptsize\ttfamily 2f99716} & \textbf{Dance envelope verified clean, framed per routine, and renders tidy up}\newline{\small verify\_dance\_paths.py on the widened envelope: 1918 IK calls, 0 failures across all three routines, both arms, collision-aware with the wearer in scene, and all three controls correct. |x| 0.26..0.50, y 0.22..0.44, z 1.02..1.58 -- so the inherited envelope holds on the paths that will actually run, which is what it needed to be checked for. TASK\_FOCUS gains d1/d2/d3, centred on the wearer because both arms work both sides and pulled back to 2.10 m.} \\
\addlinespace[3pt]
{\scriptsize\ttfamily aeb6279} & \textbf{Watched the dance: the canon was invisible and the subject was too small}\newline{\small WATCHED d1 rather than checked it, which is the only way this question can be answered. Frames at 8, 18 and 28 s: the arms tuck in, open out and unwind, so the motion is real and the widened envelope shows. Two faults in the picture: 1. THE CANON READ AS UNISON. Measured, mirror-corrected, the two arms differed by 68 mm on average -- real in the numbers and a few pixels on screen at a 1-beat lag.} \\
\addlinespace[3pt]
{\scriptsize\ttfamily b7f91f2} & \textbf{Archive the dance and demo sets before re-recording, with what changed}\newline{\small recordings/dance\_20260812 and demo\_20260812 move to extras/archive/recordings/superseded\_20260813\_dance\_prespec/. mode06\_GOOD untouched. GEOMETRY\_NOTE.md names five changes, and the first is the reason the rest were invisible: the dance clips were filmed with NO SCENE NODE, because clip\_scene rejected "d1" in its own argument list. Then the envelope (0.18 x 0.14 x 0.18 m), the canon (68 mm, a few pixels), the framing (no TASK\_FOCUS entry) and the presentation pose (lost to the follower).} \\
\addlinespace[3pt]
{\scriptsize\ttfamily 2cdd912} & \textbf{The THIRD consumer of PAD\_OFFSET: the gate that decides when the hand closes}\newline{\small T1 under 01\_master\_teleop recorded "NO GRASP RECORDED at all" with the pads 0.1 mm from the cube. The geometry was finally exact and the knuckle stayed at 0.0 for the whole clip. run\_abc.py held \code{PAD\_OFF = CT.PAD\_OFFSET} -- the LEFT arm's wrist-to-pad vector -- and uses it to decide when the arm has ARRIVED and the pending grip may be issued. The task path and the clip scene were both moved to the per-arm table; this one was missed.} \\
\addlinespace[3pt]
{\scriptsize\ttfamily c9d4e3c} & \textbf{Looked at the T1 framing: in frame, but not centred and nearly edge-on}\newline{\small The framing was verified numerically and this is the first time it has been looked at. The card plays first and reads well, the work IS in frame and not cut off, and the wearer is fully in shot -- so the TASK\_FOCUS fix held. Two things the numbers could not say:} \\
\addlinespace[3pt]
{\scriptsize\ttfamily 3f26c35} & \textbf{Mode 01 records 5 of 5, and the grasp chain is fixed end to end}\newline{\small With the third PAD\_OFFSET consumer corrected, T1 records FOUR grasps and FOUR releases with all four cubes carried 0.16-0.30 m and the fingers closing 0.0-4.4 mm from the object, against a 30 mm window. T3 records two and two. Both were "NO GRASP RECORDED at all" one run earlier, with the pads 0.1 mm from the cube.} \\
\addlinespace[3pt]
{\scriptsize\ttfamily cbede31} & \textbf{Reap ORPHANED DDS segments so the presentation pose stops losing the port race}\newline{\small Five retries did not help, because retrying does not free a port. Measured mid-sweep: 156 fastrtps segments in /dev/shm, 76 held by a live process and 80 ORPHANED. Fast DDS leaks one per participant on this host and every clip runs the stager as a fresh process, so the range fills and the clip silently opens on home.} \\
\addlinespace[3pt]
{\scriptsize\ttfamily 4ff4e58} & \textbf{A failed restore disarms the arm, so restoring is now verified by read-back}\newline{\small MY OWN BUG, and the worst kind: one clip logged "paused (motion\_enabled true -> false), RESTORE FAILED" and every later clip in that mode failed, because the follower it had silenced never spoke again. The single set\_parameters call was believed on its return value alone. resume\_followers() now retries and READS THE VALUE BACK, retrying again if it is still false, and the script exits non-zero if it cannot restore -- even when the pose itself arrived.} \\
\addlinespace[3pt]
{\scriptsize\ttfamily 48401ad} & \textbf{The SHM reaper was deleting the DISCOVERY PORTS. Exclude them.}\newline{\small /dev/shm holds two kinds of Fast DDS segment and they must be treated differently: fastrtps\_port7400 a DISCOVERY PORT, the rendezvous every participant uses. Not mapped continuously by anybody. fastrtps\_<hex> one per participant. These are what leak. The "orphaned" test -- nothing in /proc/*/maps has it open -- calls EVERY port segment orphaned, so the reaper deleted the graph's rendezvous points.} \\
\addlinespace[3pt]
{\scriptsize\ttfamily 340a48d} & \textbf{Remove the SHM reaper. "Not currently mapped" is not "orphaned".}\newline{\small I added it to stop the presentation pose losing a DDS port race, and it broke the graph twice. Both failures measured, both worse than the problem: 1. IT DELETED THE DISCOVERY PORTS. fastrtps\_port7400 is the rendezvous every participant uses and nothing maps it continuously, so the "orphaned" test called every port orphaned. After a clean restart that had just reported READY with 30 joint-state messages, no new process could join. 2. EXCLUDING THE PORTS WAS NOT ENOUGH.} \\
\addlinespace[3pt]
{\scriptsize\ttfamily ec2da99} & \textbf{Mode 01 recorded 5 of 5, every clip opening on the presentation pose}\newline{\small With the reaper gone and the graph healthy, 01\_master\_teleop completes all five MSc tasks at the corrected framing. Every clip logs 'presentation pose: OK' and the progress file records opened\_on=presentation for each, so the claim is per clip rather than per run.} \\
\addlinespace[3pt]
{\scriptsize\ttfamily 0d300f6} & \textbf{Modes 01 and 04 recorded: 10 of 10 cells OK}\newline{\small 01\_master\_teleop and 04\_vr\_shared both record all five MSc tasks at the corrected framing. 9 of 10 clips open on the presentation pose; the one that did not is recorded as opened\_on=home rather than assumed.} \\
\addlinespace[3pt]
{\scriptsize\ttfamily a1b6332} & \textbf{Mode 06 records 5 of 5, and mode 03's failures are CAPTURE TIMING not grasps}\newline{\small 06\_full\_autonomy previously scored 0/4 on T1 and now records all five tasks. With 01 and 04 also at 5/5, three of the four modes that showed the "T1 is flaky across modes" pattern are clean, which is what the left-arm PAD\_OFFSET in run\_abc's arrival gate predicted.} \\
\addlinespace[3pt]
{\scriptsize\ttfamily fa7c397} & \textbf{Sweep complete: 22 of 28 cells, and the two failure modes are DIFFERENT}\newline{\small 01, 04 and 06 record 5 of 5 each; the dance records 3 of 3. 06 previously scored 0/4 on T1, so three of the four modes that showed last session's "T1 is flaky across modes" pattern are now clean -- which is what the left-arm PAD\_OFFSET in run\_abc's arrival gate predicted. The six failures are two unrelated faults and must not be lumped together: 03\_shared\_autonomy -- CAPTURE TIMING, not grasping.} \\
\addlinespace[3pt]
{\scriptsize\ttfamily 6c934d9} & \textbf{25 of 28 cells, every one opening on the presentation pose}\newline{\small Mode 03 re-recorded 3 of 3 on a quiet stack, so its "GRASP OUTSIDE THE CLIP" failures were CONTENTION, not a synchronisation bug -- the same three cells that failed at -16.4 s of a 4.4 s window now record cleanly with 4 grasps and 4 releases. Worth knowing before anybody rewrites the motion gate: the fault did not reproduce. Standing at 01, 03, 04, 06 at 5/5, the dance at 3/3, and all 25 successful clips recording opened\_on=presentation.} \\
\addlinespace[3pt]
{\scriptsize\ttfamily 341d0ce} & \textbf{Part 5: the T2 carry series found what an end-state check cannot}\newline{\small The tilt/separation series was added because T2 could only be scored pass/fail at the end. On its first full sweep it earned that: mode tilt\_rms tilt\_max time above 6.8 deg 01\_master\_teleop 1.085 5.994 0.0 s 02\_vr\_teleop 1.643 3.522 0.0 s 03\_shared\_autonomy 1.146 5.723 0.0 s 04\_vr\_shared 10.482 29.475 8.28 s 06\_full\_autonomy 10.478 32.047 11.19 s T2's whole claim is that tilt past 6.8 degrees drops the ball.} \\
\addlinespace[3pt]
\end{longtable}

\section*{2026-08-15 \textnormal{\small\textcolor{srlgrey}{--- 20 commits}}}
\addcontentsline{toc}{section}{2026-08-15}
\begin{longtable}{@{}p{0.055\textwidth}p{0.885\textwidth}@{}}
\endfirsthead\endhead
{\scriptsize\ttfamily d782743} & \textbf{T1 to the left arm, and the workspace measured against the wearer for the first time}\newline{\small The brief was four changes to T1 and one question about the workspace. The question turned out to be the session: every workspace number this project has ever quoted comes from /compute\_ik with avoid\_collisions, and the SRDF excludes torso/harness/backpack against each arm's base, shoulder and half\_arm\_1 -- the pairs a shoulder-mounted arm actually threatens. So MoveIt returns \code{valid} for poses with the tube inside the person, and nothing had ever measured the distance.} \\
\addlinespace[3pt]
{\scriptsize\ttfamily b133a65} & \textbf{Two witnesses for a moving arm, and the wearer reaches the spoken commands}\newline{\small THE GATE. record\_abc\_sweep trusted run\_abc's exit code as its only evidence that an arm moved. The two measure the same quantity in different processes off different tf2 listeners, so the runner can see motion the scene node never received -- and that is what got filed on 2026-08-15: clips reporting TRAVEL L 0.00 R 0.00, four cubes carried 0.000 m, recorded OK. scene\_travel\_verdict() is a second, independent gate reading the scene node's own ee\_travel\_m.} \\
\addlinespace[3pt]
{\scriptsize\ttfamily e2e85b3} & \textbf{Mode 01 recorded: 5 of 5, every clip on the presentation pose}\newline{\small T0, T1, T1 stage 2, T2, T3 under 01\_master\_teleop, eight angles each plus the card. Both travel witnesses agree on every clip: t0 left 1.4071 right 1.4167 both arms reach t1 left 3.1056 right 0.5956 one-arm task, left t1s2 left 2.0120 right 1.9336 both arms, as stage 2 requires t2 left 1.0747 right 1.0071 bimanual carry t3 left 1.2575 right 1.2239 bimanual by role TWO BUGS FOUND BY RUNNING IT, both of which quietly damage a set. 1. THE LEDGER WAS OVERWRITTEN BY A SUBSET RUN.} \\
\addlinespace[3pt]
{\scriptsize\ttfamily 35d032c} & \textbf{Modes 06 and 03 recorded, and the looking pass finds three things no check can}\newline{\small 10 clips, eight angles each, every one opening on the presentation pose. Both travel witnesses agree on all ten. THREE FINDINGS FROM THE PIXELS AND THE GEOMETRY, none of which any automatic check in this repository can produce, because each is a thing that is true BY CONSTRUCTION and therefore cannot fail. 1. T2's TRAY IS ELASTIC.} \\
\addlinespace[3pt]
{\scriptsize\ttfamily 457c7d5} & \textbf{02\_vr\_teleop: 5 of 5, and the cause was the mapper's LIFETIME}\newline{\small 02\_vr\_teleop failed three of five cells on the first pass -- T1, T1 stage 2 and T3, which are exactly its three tasks that record a grasp. It never touched a cube: the pads' closest approach was 88.1, 115.4, 156.2 and 205.9 mm against a 30 mm gate, while 01, 03 and 06 closed on all four at 0.0000 m. The final frame shows all four cubes still in their starting row. TWO CAUSES, BOTH ISOLATED, AND ONE OF THEM I GOT WRONG FIRST. THE MAPPER HOLDS THE ARMS.} \\
\addlinespace[3pt]
{\scriptsize\ttfamily a3fd7b3} & \textbf{All 25 cells recorded, and three verifiers were checking nothing}\newline{\small 25 of 25 cells OK, every one opening on the presentation pose. 28 clip dirs, 19 of 19 planned data cells, no real gaps. verify\_rviz\_clips 28 of 28 with all seven of its own negative controls correct. 510 tests pass, 2 allowed. 04\_vr\_shared's two failures were the DEGRADED LONG-RUNNING STACK, not the mode.} \\
\addlinespace[3pt]
{\scriptsize\ttfamily 6309856} & \textbf{Resume note: the set is complete, and what to fix first}\newline{\small 25 of 25 cells, every one on the presentation pose. 28 clip dirs, 19 of 19 planned data cells, no real gaps. 510 tests pass. The note leads with the clearance ledger, because the thing most likely to go wrong next is somebody citing a pre-fix workspace number in good faith.} \\
\addlinespace[3pt]
{\scriptsize\ttfamily b762cad} & \textbf{Part 1: the centre is shut by the TORSO, not by the wearer's arms}\newline{\small The project's answer to "why can the work not be in the centre" has been the wearer's own forearm, measured against a mannequin whose arms hang rigidly at its sides. That is a modelling choice, so the posture is now a variable with one source and the sweep was re-run in four postures plus the limit case of no arms at all. The centre does not open in any of them.} \\
\addlinespace[3pt]
{\scriptsize\ttfamily e48f5bf} & \textbf{Part 2: the render was right, and the numbers were about something else}\newline{\small Reported: both elbows 0.14 m below the shoulders, wrists level, symmetric to 1e-16. Rendered: elbows at shoulder height and wrists angled down. Measured from /tf, which is what RViz draws, with three stills from the same boot: elbow below its own shoulder 0.0793 m left, 0.2715 m right left elbow outboard of its hand +0.188 m hands 1.10 m apart, against a 0.36 m torso tool axis elevation -0.01 deg, so the wrists ARE level forearm mirror residual 0.2826 m, against 0.0000 m at the hands It is not a cached config.} \\
\addlinespace[3pt]
{\scriptsize\ttfamily eea4cd0} & \textbf{Part 3: the null space is real now, and it still cannot help}\newline{\small An earlier pass built a seed-fan sampler, measured it worth 4.4 mm, and said plainly that seed sampling is not redundancy resolution and the real answer is an explicit parameterisation. That was built. The conclusion did not change and it is now a much stronger statement.} \\
\addlinespace[3pt]
{\scriptsize\ttfamily ca59be0} & \textbf{Part 4: the tray was elastic and T1's two stages were two tasks}\newline{\small T2. clip\_scene drew the tray between the grippers with length \code{span + 0.06}, where span is the LIVE distance between them, so the prop resized itself every frame to whatever the arms were doing. Over the shipped clips it ran 0.487 to 1.211 m against a 0.560 m spec, from 73 mm short to 651 mm long, held at both ends throughout. That is why T2-1 ("held by BOTH grippers") and T2-2 ("ball visible ON the tray") could not fail.} \\
\addlinespace[3pt]
{\scriptsize\ttfamily 7673334} & \textbf{Part 5: finding 3 is real, and its size is still not known}\newline{\small The claim under test was that scripted clips with no operator show different achieved motion by mode, made on a single run per cell against a 3\% same-mode spread taken from the one cell recorded twice. A single repeat is not a measurement, so the two ways it could have been an artefact were tested directly against the recorded set. Not recording length.} \\
\addlinespace[3pt]
{\scriptsize\ttfamily 165b238} & \textbf{Part 6 (pre-record): T1 stage 2 seed 0 was a stale pool, not a bad draw}\newline{\small The N=10 re-verification before recording reproduced the known regression: stage 2, seed 0, RIGHT arm, 26 waypoints inside the 150 mm wearer floor, worst 0.1135 m, with zero IK failures. Seeds 1 and 2 were clean, which is what made it look like one unlucky draw. It was the pool. work\_surface\_region.json was surveyed before the 2026-08-15 home change and its right-arm clear\_cells reach |x| = 0.400.} \\
\addlinespace[3pt]
{\scriptsize\ttfamily 9efd6e1} & \textbf{Part 6: mode 01 recorded, and the missing producer three verifiers wanted}\newline{\small Mode 01 recorded all five MSc cells on the teleop stack: every clip opens on the presentation pose, both arms travel (T1S2 1.635 / 1.881 m), the grasps land at +5.2 s and +6.5 s, T3 places 1 mm from target, and verify\_rviz\_clips passes 28 of 28 clips on every automatic check. Three other verifiers failed, and they were right to.} \\
\addlinespace[3pt]
{\scriptsize\ttfamily a4cb334} & \textbf{Part 6: mode 02 recorded, and T2 carries the tray with its hands open}\newline{\small Mode 02 recorded all five cells. With the grip trace now being produced, the three verifiers that had been refusing to report could finally run, and two more faults had to be cleared before they could: their glob knew only the legacy four-level path and CRASHED unpacking four components from the MSc's three, and GRASP\_TASKS held lower-case retired names while the path spells the task in capitals, so no MSc clip was ever classified as a grasp clip.} \\
\addlinespace[3pt]
{\scriptsize\ttfamily 678984b} & \textbf{Part 6: mode 03 recorded, and "0 of 10 genuine grasps" was two instrument faults}\newline{\small Mode 03 recorded all five cells. verify\_grasp\_quality ran on the MSc set for the first time and said 0 of 10 clips show a genuine grasp. The scene's own record says otherwise and is not ambiguous: T1 logs GRASPED and RELEASED for all four cubes, knuckle 0.383 to 0.395 against a needed 0.3812, min\_pad\_obj\_m 0.0 on every cube, carried 0.23 to 0.40 m, all four finishing on their pads.} \\
\addlinespace[3pt]
{\scriptsize\ttfamily f845a84} & \textbf{Part 6: mode 04 recorded}\newline{\small All five cells OK on 04\_vr\_shared, with the grip trace now produced so the three gripper verifiers have real input.} \\
\addlinespace[3pt]
{\scriptsize\ttfamily d1dbe77} & \textbf{Part 6: mode 06 recorded, all five modes now on the current geometry}\newline{\small All five cells OK on 06\_full\_autonomy. Mode 01 is being re-recorded because its clips predate the grip-trace producer and the set has to be uniform.} \\
\addlinespace[3pt]
{\scriptsize\ttfamily 08ba80b} & \textbf{Part 6: mode 01 re-recorded, so all 25 cells carry the grip trace}\newline{\small All five cells OK and five grip\_trace.json on disk. The set is now uniform: every mode recorded on the current geometry with the producer in place.} \\
\addlinespace[3pt]
{\scriptsize\ttfamily 25e1ee4} & \textbf{Part 6 complete: 25 cells in five modes, plus the three routines}\newline{\small The dance could not be filmed before today and the check that knew it crashed: verify\_dance\_paths printed an example failure with "\%s" \% worst[0] against a tuple, so it raised TypeError exactly when it had a failure to report. Fixed, it then found d1 failing 45 of 382 waypoints and d3 49 of 554, at x = -0.28 on the right arm and +0.272 on the left. The cause is the same constraint as parts 1 and 2.} \\
\addlinespace[3pt]
\end{longtable}

\section*{2026-08-16 \textnormal{\small\textcolor{srlgrey}{--- 7 commits}}}
\addcontentsline{toc}{section}{2026-08-16}
\begin{longtable}{@{}p{0.055\textwidth}p{0.885\textwidth}@{}}
\endfirsthead\endhead
{\scriptsize\ttfamily 7752e5c} & \textbf{Part 7: looked at the clips, and four of the six things asked for cannot be shown}\newline{\small Frames pulled from the re-recorded set and read against the MUST BE TRUE lists, pixels first with the measurement beside them. Stills in docs/img/home\_render. symmetric home, elbows down NO 0.0793 (L) and 0.2715 (R) below their shoulders, forearm mirror residual 0.2826 objects ON the table NO they float; visible in T1, T1S2 and T3 pads and cubes in the CENTRE NO min |x| 0.325 / 0.450, and 0.300 / 0.375 even with the wearer's arms deleted arms in FRONT of the person NO T1S2 shows both out at the table's ends the tra [...]} \\
\addlinespace[3pt]
{\scriptsize\ttfamily 9d246dc} & \textbf{Home solved from constraints, T1 layout re-measured, vision drives the grasp}\newline{\small THE HOME POSE WAS SOLVED, NOT CHOSEN. scripts/solve\_home\_pose.py states it as nine geometric constraints and searches the 14-DOF joint space for them, with FK validated against live TF to 0.06 mm. Achieved, measured from the render: elbows 0.2776/0.2794 m below their shoulders (was 0.0793/0.2715), camera 12.05 deg off vertical (was 56.4 mm BELOW the fingertips -- the hand was upside down and nothing had ever constrained the roll), approach level to 1.5 deg and facing forward, whole-chain mirror residual 0.0272 m ag [...]} \\
\addlinespace[3pt]
{\scriptsize\ttfamily ffc50ba} & \textbf{Record T0 and T3 in mode 03 against the new home and layout}\newline{\small Both cells clean, run exited 0. Mode 02 is deferred: it records a STATIONARY ARM on both cells (ee\_travel 0.0000 either side) while its isolation reports correctly -- vr\_pose\_mapper started, stopped for the staging move and restarted. Its clips are left at the committed versions rather than overwritten with failed runs.} \\
\addlinespace[3pt]
{\scriptsize\ttfamily ac3919b} & \textbf{Record T0 and T3 in mode 06 against the new home and layout}\newline{\small Both cells clean. Mode 04 deferred with mode 02: both VR modes record a STATIONARY ARM (ee\_travel 0.0000 either side) while every non-VR mode (01, 03, 06) records cleanly. The split is exactly the vr\_pose\_mapper path. Their clips are left at the committed versions.} \\
\addlinespace[3pt]
{\scriptsize\ttfamily 639b2fc} & \textbf{Record T0 and T3 in mode 02, and fix the staging pose that was failing runs}\newline{\small THE PRESENTATION POSE WAS STILL THE 2026-08-15 ONE. The sweep stages it before every clip, and after home was re-solved it sat 2.9465 rad (left) / 2.7723 (right) away -- a 7.6-9.3 s move that regularly did not arrive. Cells then recorded a STATIONARY ARM (ee\_travel 0.0000 both sides) and one clip logged "NOT STAGED -- this clip opens on HOME", which breaks R-7 silently. The pose is now loaded from config/home\_positions\_\{arm\}.txt, so it cannot drift from home again.} \\
\addlinespace[3pt]
{\scriptsize\ttfamily 1d001d2} & \textbf{Record T0 and T3 in mode 04: all five modes now on the new home and layout}\newline{\small T0 and T3 are complete across 01, 02, 03, 04 and 06 against the solved home, the re-validated region and the per-arm pad layout.} \\
\addlinespace[3pt]
{\scriptsize\ttfamily 89c8dad} & \textbf{Record d1 and d2 in all five modes; d3 is REFUSED, not missing}\newline{\small The restart-per-mode protocol is what makes these land: after any VR-mode cell the DDS graph is degraded and the NEXT run stages nothing and films a stationary arm, whatever mode it is. Clearing /dev/shm and relaunching before each mode gives 2 of 2 on 02, 03, 04 and 06. D3 IS NOT RECORDED AND MUST NOT BE. verify\_dance\_paths, N=3 collision-aware from the new home: d1 0 of 382 fail, d2 0 of 356 fail, d3 18 of 560 FAIL -- e.g. right waypoint 124 at (-0.740, 0.435, 1.567).} \\
\addlinespace[3pt]
\end{longtable}

\section*{2026-08-17 \textnormal{\small\textcolor{srlgrey}{--- 29 commits}}}
\addcontentsline{toc}{section}{2026-08-17}
\begin{longtable}{@{}p{0.055\textwidth}p{0.885\textwidth}@{}}
\endfirsthead\endhead
{\scriptsize\ttfamily e26389a} & \textbf{T1 is BLOCKED in the vision path, not the layout: 28-31 mm detections against a 30 mm gate}\newline{\small THE GEOMETRY IS CLEAN AND THAT IS WHY THIS IS NOT A LAYOUT PROBLEM. verify\_t1\_paths at the solved home, N=10, wearer and furniture, clearance geometric: 0 IK failures and 0 waypoints inside the 150 mm floor over 7522 IK calls, stage 1 and all three stage-2 seeds, both controls correct. t1\_paths.json had been stale since 2026-08-15 -- it predates the home solve -- and is refreshed here. Stage 2 seed 0's right arm, 25 of 98 waypoints inside the floor at the old home, is CLEAN; that regression is closed.} \\
\addlinespace[3pt]
{\scriptsize\ttfamily 3f1a509} & \textbf{A run starts from HOME, and the second one did not}\newline{\small The reported symptom was "RViz boots into the new pose but a task run starts from the old one". Measured on a fresh teleop stack, reading /joint\_states against config/home\_positions\_\{arm\}.txt at the instant the FIRST waypoint is published on the mode's entry topic: at boot left 0.0000 right 0.0000 rad run 1, first commanded waypoint left 0.0000 right 0.0000 between runs (nothing else ran) left 1.2338 right 0.6248 run 2, first commanded waypoint left 1.2338 right 0.6248 So the pose is NOT lost between launch and the [...]} \\
\addlinespace[3pt]
{\scriptsize\ttfamily e68dfe7} & \textbf{T1 from a typed sentence, grounded on what the camera saw}\newline{\small WHAT THE SHIPPED GRAMMAR DID WITH T1'S OWN INSTRUCTION, measured before any of this was written: "pick up the blue cube and put it on the blue pad" -> grab / target "blue cube" / destination None. The destination was SILENTLY DROPPED, which is the dangerous category "put the green ones on the green mat" -> refused, "no known verb" "move that blue block to its colour" -> refused, "heard 'move to' but no place I know" One misunderstood, two refused, and the task the whole clip set is about could not be said.} \\
\addlinespace[3pt]
{\scriptsize\ttfamily b4d9813} & \textbf{The T1 vision fault, reproduced exactly and fixed at both its causes}\newline{\small REPRODUCED, and worse than the 2026-08-17 note recorded. Inside the recording sweep the four T1 cubes deprojected 31.7 / 32.7 / 34.0 / 35.7 mm from truth against a 30 mm capture gate, so NONE of the four was grasped -- not just cube\_3. The pad miss at closure in the same clip was 31.6 / 32.6 / 33.9 / 35.6 mm the same numbers to 0.1 mm. The grasp failure IS the detection error, one to one, which is the strongest evidence in the file that nothing else about T1 is wrong.} \\
\addlinespace[3pt]
{\scriptsize\ttfamily b1b664d} & \textbf{THE FOLLOWER PAUSE WAS DECORATIVE. That is the root of the T1 vision fault.}\newline{\small \code{ik\_follower\_node} read \code{motion\_enabled} ONCE at construction into \code{self.motion\_enabled} and registered no parameter callback. \code{set\_parameters} therefore changed the parameter server's copy and the follower carried on publishing -- so every caller that "paused the follower" by lowering it paused NOTHING, and the node's own blocker recovery text (\code{ros2 param set motion\_enabled true}) could not work either.} \\
\addlinespace[3pt]
{\scriptsize\ttfamily 554924f} & \textbf{Restore the committed T1/01 clip: the last commit swept in the FAILED one}\newline{\small The mode-01 recording that ran while the follower pause was still decorative overwrote the clip files on disk, and \code{git add -A} took them. That clip is the one whose four cubes closed 31.6-35.6 mm from the pads and which the sweep itself reported as NO GRASP RECORDED at all -- it must not sit in the tree as evidence for anything. Restored from b4d9813, which is the 2026-08-16 set.} \\
\addlinespace[3pt]
{\scriptsize\ttfamily 03a863b} & \textbf{The stillness gate WAITS for a still frame instead of refusing the first}\newline{\small It refused the first bad frame, and inside the recording sweep that turned a transient into a lost clip -- twice, both times reading EXACTLY 0.0170 m. A value that repeats to 0.1 mm across two differently loaded runs is a fixed lag, not an arm still travelling, and the refusal message reported only the distance so the two could not be told apart. It now prints both positions and the delta.} \\
\addlinespace[3pt]
{\scriptsize\ttfamily 07df66b} & \textbf{T1 re-recorded under 01\_master\_teleop, and it grasps all four again}\newline{\small The first T1 clip since the vision fault was found. Same code path, same sweep, same load as the clip that failed this morning. in-sweep detection error 1.30 / 0.50 / 1.30 / 2.30 mm was 31.7 / 32.7 / 34.0 / 35.7 mm min\_pad\_obj\_closed\_m 0.0000 on all four cubes was 0.0316 / 0.0326 / 0.0339 / 0.0356 against a 30 mm gate -- 0 of 4 grasped carried 0.2074 / 0.3267 / 0.3005 / 0.2569 m was 0.0000 on all four placed 30 mm from target, was 101 mm left EE path 5.1561 m Every cube went to the pad of its own colour: cube\_0 and [...]} \\
\addlinespace[3pt]
{\scriptsize\ttfamily 9d99007} & \textbf{T1 under 03\_shared\_autonomy: 4 of 4, every closure at 0.0000 m}\newline{\small detection error 2.30 / 1.22 / 1.52 / 2.19 mm min\_pad\_obj\_closed\_m 0.0000 on all four carried 0.2118 / 0.2963 / 0.3020 / 0.2541 m placed 31 mm from target left EE path 5.1722 m Blue cubes released at x = 0.5952 and 0.5952 against a blue pad at 0.595; green at 0.8149 and 0.8223 against a green pad at 0.825.} \\
\addlinespace[3pt]
{\scriptsize\ttfamily 63b5b37} & \textbf{T1 under 04\_vr\_shared: 4 of 4, every closure at 0.0000 m}\newline{\small detection error 1.48 / 1.55 / 2.08 / 3.09 mm min\_pad\_obj\_closed\_m 0.0000 on all four carried 0.2168 / 0.3282 / 0.2984 / 0.2452 m placed 30 mm from target left EE path 5.1551 m Blue released at x = 0.5951 and 0.5950 against a blue pad at 0.595; green at 0.8226 and 0.8211 against a green pad at 0.825. This mode also exercises the VR mapper being stopped for staging and restarted before the run, which is where 02 and 04 have failed before.} \\
\addlinespace[3pt]
{\scriptsize\ttfamily 484b336} & \textbf{T1 under 06\_full\_autonomy: 4 of 4, every closure at 0.0000 m}\newline{\small detection error 1.94 / 1.36 / 0.92 / 1.94 mm min\_pad\_obj\_closed\_m 0.0000 on all four carried 0.2102 / 0.3240 / 0.3067 / 0.2568 m placed 30 mm from target left EE path 5.1843 m Blue released at x = 0.5951 and 0.5950 against a blue pad at 0.595; green at 0.8208 and 0.8227 against a green pad at 0.825.} \\
\addlinespace[3pt]
{\scriptsize\ttfamily fa51d46} & \textbf{T1 under 02\_vr\_teleop: 4 of 4, and that CONTRADICTS a documented state}\newline{\small docs/ENGINEERING_LOG.md records 02\_vr\_teleop as 'cannot grasp -- all three of its grasping tasks fail; the pads miss by 88 to 206 mm against a 30 mm gate, and the miss ACCUMULATES'. On T1, measured now: detection error 1.53 / 1.55 / 2.08 / 3.09 mm min\_pad\_obj\_closed\_m 0.0000 on all four carried 0.2130 / 0.3066 / 0.2662 / 0.2440 m placed 30 mm from target left EE path 5.1453 m Blue released at x = 0.5951 and 0.5958 against a blue pad at 0.595; green at 0.8176 and 0.8226 against a green pad at 0.825.} \\
\addlinespace[3pt]
{\scriptsize\ttfamily 027d265} & \textbf{The objects float because Rig was measuring at the wrong wrist, and T1-1 is BLOCKED}\newline{\small THE INSTRUMENT FIRST. \code{measure\_what\_binds.Rig} -- which verify\_t1\_paths, search\_centre\_on\_surface, measure\_objects\_on\_the\_table and every other reachability sweep in this repo solves through -- set \code{self.quat} to the LIVE end-effector orientation read off TF at construction, i.e. the HOME wrist. The tasks do not command that: \code{run\_abc.send()} writes \code{WORKSPACE\_ORIENT} into every waypoint of every mode.} \\
\addlinespace[3pt]
{\scriptsize\ttfamily ca4bab1} & \textbf{T1S2 seed 2 was always 16 failures, and the docs now say which wrist each number came from}\newline{\small The refreshed \code{t1\_paths.json} baseline, at the anchor on the raised surface: stage 1 clean everywhere -- the full 171-waypoint path, every per-object pick and both pad slots on both pads, 0 IK failures and 0 below the clearance floor at N=10 -- and stage 2 seeds 0 and 1 clean on both arms. STAGE 2 SEED 2's LEFT ARM IS 16 OF 91, AND IT IS NOT THE SURFACE. Measured at the raised 0.980 and again at the old 0.950: identically 16.} \\
\addlinespace[3pt]
{\scriptsize\ttfamily 658d861} & \textbf{The arms are NOT at home at the first commanded waypoint, and now a clip fails for it}\newline{\small The operator looked at the last frame of the newest T1 clip and said the arms were at the old pose. They are, the clip is not stale, and nothing in the pipeline could see it. WHAT EVERY EXISTING CHECK SAID, recording T1 under 01\_master\_teleop: stage\_presentation\_pose.py succeeded, so the clip records \code{opened\_on: "presentation"}; the staged look reported \code{home to 0.0000 rad}; \code{require\_home()} returned homed and the run exited 0; verify\_rviz\_clips passed 37 of 37.} \\
\addlinespace[3pt]
{\scriptsize\ttfamily f413e21} & \textbf{Two checks were pointing at the wrong thing: the home gate and the yellow cube}\newline{\small THE HOME GATE I ADDED LAST SESSION WAS MEASURING THE WRONG INSTANT, and it failed every clip including good ones. \code{run\_abc} deliberately drives the arms to waypoint 0 and settles them there before it starts measuring -- travel taken from wherever the previous run left the arm reported 0.0118 m instead of 0.4761 m, which is why that block exists. My probe sat at the first iteration of the send loop, i.e. AFTER that approach, called it "the first commanded waypoint" and read 0.7991 rad.} \\
\addlinespace[3pt]
{\scriptsize\ttfamily e54049b} & \textbf{Mode 01: T1 and T1S2 re-recorded on the corrected geometry}\newline{\small Both cells, all eight angles, card first. Both open on the current home -- AT HOME 0.0000 rad on both arms at the last instant before the approach, which is the reading that answers the question. T1 4 of 4 cubes closed at 0.0000 m, each placed on the pad of its own colour: blue slots at x 0.595, green slots at x 0.825. T1S2 4 of 4 closed at 0.0000 m, both arms, cubes blue/green by pad.} \\
\addlinespace[3pt]
{\scriptsize\ttfamily 47ca64a} & \textbf{Mode 02: T1 and T1S2 recorded, and the detector stopped counting pads as cubes}\newline{\small T1 REFUSED TWICE BEFORE THIS, with "saw 8 cubes" and then "saw 9 cubes, the task expects 4", while the same code standalone on the same stack saw 4. The refusals were correct -- they skipped rather than filming from a stale detection file -- but the message was only a count, so the sweep's own dump was extended to list every accepted and rejected blob with its pixel size, the size a cube would be at that range, and its depth. The nine then split at a glance: 4 detections z 1.1118 ..} \\
\addlinespace[3pt]
{\scriptsize\ttfamily 64e89f8} & \textbf{Mode 03: T1 and T1S2 recorded on the corrected geometry}\newline{\small Both cells AT HOME 0.0000 rad on both arms at the pre-approach instant, 4 of 4 cubes closed at 0.0000 m, each on the pad of its own colour. Detection clean at four cubes with the work-plane and duplicate rules in place.} \\
\addlinespace[3pt]
{\scriptsize\ttfamily 9cee332} & \textbf{Mode 04: T1 and T1S2 recorded on the corrected geometry}\newline{\small Both cells AT HOME 0.0000 rad on both arms at the pre-approach instant, 4 of 4 cubes closed at 0.0000 m, each on the pad of its own colour.} \\
\addlinespace[3pt]
{\scriptsize\ttfamily dcf421d} & \textbf{Mode 06: T1 and T1S2 recorded -- all ten cells now on the corrected geometry}\newline{\small Both cells AT HOME 0.0000 rad on both arms at the pre-approach instant, 4 of 4 cubes closed at 0.0000 m, each on the pad of its own colour. That completes T1 and T1S2 across all five modes, every cell re-recorded against the raised surface, the outline marking and the corrected detector. Every one reports AT HOME 0.0000 and 4 of 4 closures at 0.0000 m.} \\
\addlinespace[3pt]
{\scriptsize\ttfamily 481fda2} & \textbf{A clip counts for the task it CONTAINS, not the folder it sits in}\newline{\small WHAT I WAS ASKED TO FIX, AND WHAT I FOUND. The report was that T1S2 showed as covered because T1's clips were being counted. It is not what is happening, and the numbers say so plainly: T1S2 has its own clips in all five modes, every one containing genuine stage-2 content (cube\_left\_0/1, cube\_right\_0/1), and \code{status\_table} already keys on the exact directory name -- the prefix bug ("t1s2".startswith("t1")) was fixed in that file earlier and its scar is in the comments.} \\
\addlinespace[3pt]
{\scriptsize\ttfamily 34f9944} & \textbf{The marking goes back to the plane it bounds, and the check can now fail on the real scene}\newline{\small TWO THINGS WERE ASKED FOR, AND ONE OF THEM I CANNOT DO. Taking them in order. THE MARKING WAS 116.5 mm BELOW THE PADS, and that was my regression from the previous session. I moved it from the work plane down to the table on the argument that "paint goes on the thing it is painted on". That is wrong for the reason the frame shows immediately: the marking is the boundary a participant is told to keep the WORK inside, and drawing it under the work put the whole task outside its own boundary.} \\
\addlinespace[3pt]
{\scriptsize\ttfamily e05e384} & \textbf{T1S2 mode 01: re-recorded with the marking at the plane it bounds}\newline{\small AT HOME 0.0000 rad on both arms, 4 of 4 closures at 0.0000 m. The markings now enclose the work instead of sitting 116.5 mm under it.} \\
\addlinespace[3pt]
{\scriptsize\ttfamily 29f06fb} & \textbf{T1S2 mode 02: re-recorded with the marking at the plane it bounds}\newline{\small AT HOME 0.0000 rad on both arms, 4 of 4 closures at 0.0000 m. The first attempt aborted: the virtual display would not render after three restarts and record\_rviz REFUSED rather than filing a black view. Two orphaned rviz2 processes from an earlier run were holding the display; the per-mode cleanup now kills viewers as well as displays.} \\
\addlinespace[3pt]
{\scriptsize\ttfamily 5baa428} & \textbf{T1S2 mode 03: re-recorded with the marking at the plane it bounds}\newline{\small AT HOME 0.0000 rad on both arms, 4 of 4 closures at 0.0000 m.} \\
\addlinespace[3pt]
{\scriptsize\ttfamily 7914958} & \textbf{T1S2 mode 04: re-recorded with the marking at the plane it bounds}\newline{\small AT HOME 0.0000 rad on both arms, 4 of 4 closures at 0.0000 m.} \\
\addlinespace[3pt]
{\scriptsize\ttfamily efa878d} & \textbf{T1S2 mode 06: all five cells re-recorded with the marking at the plane it bounds}\newline{\small AT HOME 0.0000 rad on both arms, 4 of 4 closures at 0.0000 m. That completes T1S2 across all five modes on the corrected marking height.} \\
\addlinespace[3pt]
{\scriptsize\ttfamily b782a9b} & \textbf{Record the marking regression, the anchor re-price of supports, and T1S2 coverage}\newline{\small findings.md carries the full account: the marking sat 116.5 mm below the pads because I moved it to the table last session on an argument about paint that is the wrong argument for a boundary; the on-table check was asking about T1S2 and measuring T1; and the support family, re-priced at the anchor, costs 109 of 171 waypoints where the home-wrist note claimed 22. Also recorded: the two ordering bugs in my own negative control, caught by running it rather than by reading it.} \\
\addlinespace[3pt]
\end{longtable}

\section*{2026-08-18 \textnormal{\small\textcolor{srlgrey}{--- 19 commits}}}
\addcontentsline{toc}{section}{2026-08-18}
\begin{longtable}{@{}p{0.055\textwidth}p{0.885\textwidth}@{}}
\endfirsthead\endhead
{\scriptsize\ttfamily 4fd9a76} & \textbf{T1 verifies at the measured pad offset: symmetric pads, and the grasp lands on the cube}\newline{\small The T1 rebuild and the mount change were in the working tree, uncommitted, with \code{test\_t1\_layout\_is\_verified} failing on all three of its assertions -- the shipped constants and the walked record described different layouts. Walking the composed path found two real faults on the way. grasp\_frames.PAD\_MID\_EE was 13.47 mm too long. It had been DERIVED, as the magnitude of clip\_tasks.PAD\_OFFSET\_BY\_ARM -- a world-frame vector recorded at the anchor -- rather than measured.} \\
\addlinespace[3pt]
{\scriptsize\ttfamily 5f746df} & \textbf{T1 universally commandable: selectors, compounds, an answerable question, and stage 2 on stage 1's geometry}\newline{\small THE ONE MISUNDERSTANDING, AND IT WAS LIVE. "put the blue ones on the blue pad and the green ones on the green pad" -- two complete instructions sharing a verb -- planned the FIRST clause and dropped the second in silence: two cubes moved of four, reported as a success. Three of the 75 phrasings did it, all compounds. \code{split\_destination} cut at the first destination cue, so the head was "put the blue ones" and the two-targets check, which counts targets in the head, saw one and passed it.} \\
\addlinespace[3pt]
{\scriptsize\ttfamily 3948edd} & \textbf{The pre-recording audit is alive again, and T1 passes all twelve of its criteria}\newline{\small TASK\_SPEC section 8 requires a code-side pass over every MUST BE TRUE item before any recording, and \code{scripts/audit\_task\_spec.py} is its mechanical form. It had been dead since the 2026-08-17 rebuild: it calls \code{msc\_clip\_tasks.t1()}, the builder moved to \code{t1\_task} and the alias went with it, so the check that exists to catch "reported done and later found absent" died with AttributeError the first time it was run afterwards.} \\
\addlinespace[3pt]
{\scriptsize\ttfamily 2851f71} & \textbf{Drop three imports my own edits left unused} \\
\addlinespace[3pt]
{\scriptsize\ttfamily 7e5beb2} & \textbf{The GUI button audit passes for the first time, and getting there found four defects}\newline{\small \code{verify\_gui\_buttons} level 2 presses every button in the operations window and requires each press to produce an observable outcome. It had never completed: it hung on button 05 of 69 and was killed, every button after it unpressed, the run reporting nothing at all. What it was hung on is the modal CONSENT dialog on START SESSION, which is exactly what should block a machine.} \\
\addlinespace[3pt]
{\scriptsize\ttfamily 3eddbea} & \textbf{The observe pose, re-solved for the geometry that exists, and a staging check that could not fail}\newline{\small THE SOLVER REFUSED ON ITS OWN FK CONTROL, and it was right to. \code{srl\_fk} compares itself against live TF recorded in \code{home\_render.json}, and that recording is of the PRE-MOUNT-CHANGE robot: shoulder at 0.3473 against the 0.4977 the arm is at now, 0.1553 m apart, which is the 150 mm the mounts moved. The FK was correct and the reference was stale. Refreshed from live TF; FK now agrees to 0.000061 m.} \\
\addlinespace[3pt]
{\scriptsize\ttfamily 6f77517} & \textbf{Bring the state table and the spec up to what has been measured}\newline{\small docs/ENGINEERING_LOG.md: the unit count, the marking (T1 and T1S2 draw their OWN cells, not the global survey, because not one surveyed cell is over this task's work), the observe pose, the GUI prompt box and the button audit. TASK\_SPEC: R-5 now requires the typed instruction on the card verbatim, because under 06 the instruction IS the trial's input and a viewer cannot tell from the footage which sentence produced it. And a new P-D: every detection must be AT an object.} \\
\addlinespace[3pt]
{\scriptsize\ttfamily 6bcc031} & \textbf{sim\_session lost every argument containing a space}\newline{\small \code{" ".join(cmd)} handed the shell a string, so \code{--instruct "put every cube on the pad of its own colour"} reached the sweep as \code{--instruct put} plus six unrecognised extras. The recording died with an argparse usage message AFTER the stack was up and the unit suite had run -- a five minute setup spent to be told the command line was wrong. The shell is still needed (SRC sources the ROS setup and that has to happen in the same shell), so the arguments are quoted rather than the shell dropped.} \\
\addlinespace[3pt]
{\scriptsize\ttfamily eabfb25} & \textbf{The run started the schedule before the arms reached waypoint 0, and it cost the first grasp of each arm}\newline{\small MEASURED, on the first instruction-driven T1 recording. The clip ran, exited 0, travelled 4.69 m on the left and 1.43 m on the right -- and TWO OF FOUR CUBES WERE NEVER PICKED UP: cube\_0 left carried 0.0000 m hand closed 47.8 mm away cube\_1 left carried 0.3125 m hand closed 0.0 mm away cube\_2 right carried 0.0000 m hand closed 49.6 mm away cube\_3 right carried 0.3204 m hand closed 0.0 mm away against a 30 mm capture gate. The pattern is the diagnosis: it is the FIRST cube of each arm, and only the first.} \\
\addlinespace[3pt]
{\scriptsize\ttfamily 321217d} & \textbf{Stage 2 was commanded at the pinned anchor while its coordinates were solved at T1's approach}\newline{\small MEASURED, on the first stage 2 recording, and it read as a perfect approach and a dead gripper: every cube pads 1.4 to 4.2 mm away knuckle 0.00 throughout every cube carried 0.0000 m NO GRASP RECORDED AT ALL The pad distances look like the best numbers this task has ever produced. They are two descriptions sharing one error.} \\
\addlinespace[3pt]
{\scriptsize\ttfamily fce9730} & \textbf{One look per arm, because no single arm can see this layout's row at a usable range}\newline{\small MEASURED, and it is a property of the geometry rather than a bug. The rebuilt T1 spans 1.08 m of table with a pad pair in the middle. Asked for a pose from which ONE arm sees all twelve columns within 1.05 m, the solver returns NOTHING -- and the poses that saw everything spanned 0.62 to 1.34 m. WHAT THAT SPAN COSTS. A 40 mm cube at 1.34 m is 18 px across and the same-coloured pad 60 mm behind it is 200.} \\
\addlinespace[3pt]
{\scriptsize\ttfamily b422005} & \textbf{The mat is excluded by WHERE IT IS, not by a six-millimetre threshold}\newline{\small The pads are the cubes' own colours -- deliberately, so "each cube ended on the pad of its own colour" is judgeable from a frame -- and on the rebuilt layout they are mats lying ON the table the cubes stand on. So a fragment of pad is the right colour, at very nearly the right height, and at 0.5 m it is the right SIZE too: the left camera returned five blue blobs where there are two cubes, the extra three being pad fragments 37 to 51 px against an expected cube of 38 to 45.} \\
\addlinespace[3pt]
{\scriptsize\ttfamily 6951888} & \textbf{Keep what a launched job said, and record where T1 stands}\newline{\small \code{Gui.on\_launch} sent a launched job's stdout AND stderr to DEVNULL, so a job that REFUSED -- naming its reason, as every refusal in this repository is written to do -- produced a bare exit code and nothing else. Measured the hard way: \code{run\_abc} exits 2 inside the recording sweep and runs perfectly outside it, and the one sentence saying which of its two refusals had fired was being discarded at that line. Each launch now writes to its own file in the scratch directory and the event log says where.} \\
\addlinespace[3pt]
{\scriptsize\ttfamily d1cd8bb} & \textbf{Archive two clips their own verifiers rejected, so no failed footage sits in the verification tree}\newline{\small Both were recorded and both were REFUSED by the checks that ran immediately after them, and \code{abc\_sweep\_progress.json} records ok=False for each. Left in \code{recordings/verification/} they would read as the T1 and T1S2 evidence, which is precisely the confusion this project has been bitten by before. T1: \code{run exited 2; NO HOME RECORD; NO GRASP RECORDED at all}.} \\
\addlinespace[3pt]
{\scriptsize\ttfamily 32751e3} & \textbf{Wait for the fingers, not for a fixed number of ticks}\newline{\small The first mode-06 T1 re-record picked two cubes and left two, and reported exit 0. clip\_scene's closest-approach numbers separate two different failures: cube\_0's pads reached 11.0 mm of it with the knuckle at 0.0007 -- hand OPEN -- and did not read closed until 48.8 mm away; cube\_2 never got closer than 42.3 mm at all, 58\% of the way down a 100 mm standoff. The same run launched by hand, with nothing recording, picked all four.} \\
\addlinespace[3pt]
{\scriptsize\ttfamily 962e264} & \textbf{Pull frames out of a clip at the moments its own event log names}\newline{\small The brief ends "the frame is the evidence", and every automatic check here has at some point passed a clip that was visibly wrong. This judges nothing: it extracts frames and prints where they are. Not evenly spaced -- evenly spaced frames of a 30 s clip are mostly the arm traversing, which is the part nothing goes wrong in. The instants come from scene\_events.json's own wall clock, offset by the clip's t0\_wall, so the frames land on the card, each grasp, each release and the final state.} \\
\addlinespace[3pt]
{\scriptsize\ttfamily 84457b1} & \textbf{Measure stage 2 at the approach it actually commands}\newline{\small The scene decided where a graspable object is from \code{Scene.\_pad\_offset()}, guarded by \code{if self.task == "t1"}. Stage 2 was rebuilt on stage 1's geometry on 2026-08-18 and commands stage 1's approach -- msc\_clip\_tasks gives t1 and t1s2 the same \code{orient} -- but an exact match on "t1" is not a match on "t1s2", so every stage-2 object was measured at the ANCHOR instead.} \\
\addlinespace[3pt]
{\scriptsize\ttfamily 448b57f} & \textbf{Record T1 both stages under 06, from a typed sentence}\newline{\small Two clips, mode 06 only, eight angles each, every one opening on a card that carries the operator's sentence verbatim and then the LOOK before the task: 06\_full\_autonomy/T1/S1\_both\_arms\_centre two blue cubes left, two green right, 4 grasps, 4 releases 06\_full\_autonomy/T1S2/S2\_both\_arms\_random seed 3, a 3/1 draw, 4 grasps, 4 releases into three distinct blue slots and one green Both ran from the instruction "put every cube on the pad of its own colour", both started AT home (0.0000 rad, both arms), and the destinati [...]} \\
\addlinespace[3pt]
{\scriptsize\ttfamily 3abb809} & \textbf{Write down what the two clips cost, and what still reads wrong}\newline{\small docs/ENGINEERING_LOG.md gains the four rows this session earned: T1 recorded under 06 both stages, the stage-2 anchor defect, the frame extractor that never extracted a grasp, and the mislabel control's re-measurement. NEXT\_SESSION gets a resume point whose point is that every one of the four failures was the INSTRUMENT, not the robot -- including a "new" failing test that was two leftover stacks.} \\
\addlinespace[3pt]
\end{longtable}

\section*{2026-08-19 \textnormal{\small\textcolor{srlgrey}{--- 4 commits}}}
\addcontentsline{toc}{section}{2026-08-19}
\begin{longtable}{@{}p{0.055\textwidth}p{0.885\textwidth}@{}}
\endfirsthead\endhead
{\scriptsize\ttfamily 5bb604f} & \textbf{Record two things that are true today and would be assumed wrong tomorrow}\newline{\small Neither is a defect. Both are assumptions a reader would make from the numbers as they stand, so TASK\_SPEC 4B states them and the section-4 gap table points at them. U-4. A VOICE SCORE IS A PARSER SCORE. The 75-phrasing sweep runs \code{voice\_listener}'s TEXT source -- no audio in that path at all.} \\
\addlinespace[3pt]
{\scriptsize\ttfamily 438c91a} & \textbf{Re-record T1 both stages under 06, 2026-08-19}\newline{\small A fresh recording of both cells on the current tree, not a re-use of yesterday's. Same instruction, same seed, and the same result, which is the point: the two clips reproduce. T1/S1\_both\_arms\_centre 2 blue left, 2 green right; grasps at 0.42, 0.48 (left) and -0.48, -0.42 (right); 4 releases into the two pads' slots T1S2/S2\_both\_arms\_random seed 3, the 3/1 draw; grasps at 0.42, 0.48, 0.54 (left) and -0.48 (right); three distinct blue slots at 0.2198 / 0.2798 / 0.3398 and one green at -0.2798 Both started AT home (0 [...]} \\
\addlinespace[3pt]
{\scriptsize\ttfamily 800194a} & \textbf{VR over wifi works end to end with a real Quest, no adb}\newline{\small The WebXR route was in the tree as a documented fallback and had never been run. Running it found that it could not start at all, and that its axis map was for a different frame. Both fixed; the route is now verified with a physical headset on the lab wifi.} \\
\addlinespace[3pt]
{\scriptsize\ttfamily 20c8acc} & \textbf{Redesign for desk operation: motion capture, not VR}\newline{\small The operator sits across the room FACING the wearer, holds the controllers as a 6-DOF input, and watches the real robot. Nobody wears the headset -- it stands on a shelf as the tracking reference. Every head-worn assumption was wrong about who is where, and two of them were load-bearing. ROUTE DECISION, ON FACTS SteamVR/OpenVR over Quest Link was evaluated and rejected. Checked on this machine: Meta Quest Link PC app NOT installed, SteamVR NOT installed (Steam is), no OpenXR runtime registered, Quest not on USB.} \\
\addlinespace[3pt]
\end{longtable}

\section*{2026-08-20 \textnormal{\small\textcolor{srlgrey}{--- 4 commits}}}
\addcontentsline{toc}{section}{2026-08-20}
\begin{longtable}{@{}p{0.055\textwidth}p{0.885\textwidth}@{}}
\endfirsthead\endhead
{\scriptsize\ttfamily 3b8f051} & \textbf{One window for the whole session, and the real arms are in it}\newline{\small The GUI starts from one command and everything is reachable from it. Nothing was rewritten; the pieces that already worked are where they were. TWO VIEWS, SIDE BY SIDE. RViz is the COMMANDED view and is embedded where a window manager exists. That is a correction, not a preference: the "embedding does not work here" conclusion was measured on a host with no window manager and written down as if it were a property of the technique.} \\
\addlinespace[3pt]
{\scriptsize\ttfamily 41cdf41} & \textbf{The wearer is a person now, and the camera may only make them bigger}\newline{\small The scene camera path is built end to end: a USB camera node, both calibration tools, a markerless body tracker, per-segment gating, the fused body in move\_group's planning scene, and a GUI tab whose top line is which body the robot is planning against. NO CAMERA HAS EVER BEEN ATTACHED TO THIS MACHINE. Everything optical is still inferred and is listed as such in doc 19 section 6 rather than estimated into a table.} \\
\addlinespace[3pt]
{\scriptsize\ttfamily 6e2bf68} & \textbf{The observer can be somewhere, and the system can tell when they leave}\newline{\small Finishing the VR phase for a lab day. The deliverable is docs/system/VR\_REAL\_ARM\_RUN.md: a numbered procedure with a pass condition and a failure branch at every step, covering every connection problem this project has hit, and a section 14 listing what is still unverified ranked by how likely it is to bite. FOUR DEFECTS, ALL IN THE SAFETY PATH. The observer e-stop had NEVER BEEN SATISFIABLE.} \\
\addlinespace[3pt]
{\scriptsize\ttfamily 534a033} & \textbf{One button starts VR, and pressing it repeatedly is how the bugs fell out}\newline{\small START VR TELEOP, top of the RUN tab. Twelve steps in an order that is load-bearing -- the simulation before the mapper, because the mapper refuses to engage without robot TF and the symptom is a clutch that silently does nothing; the certificate and the port before the bridge, because a bridge that dies on a bound port leaves four healthy nodes and a URL nobody is serving. Each step fails in plain words with the fix as a button. It cannot reach a real arm, by construction and by test.} \\
\addlinespace[3pt]
\end{longtable}

\section*{2026-08-21 \textnormal{\small\textcolor{srlgrey}{--- 1 commit}}}
\addcontentsline{toc}{section}{2026-08-21}
\begin{longtable}{@{}p{0.055\textwidth}p{0.885\textwidth}@{}}
\endfirsthead\endhead
{\scriptsize\ttfamily d1390f0} & \textbf{The instrument was wrong six times out of eight, and the arms were fine}\newline{\small Two real-arm sessions. Perception found the cube every time and nothing was picked up, because the failures were never detection: they were calibration, kinematics, and eight defects in the code that COMMANDS the arms. The one that was already in the tree: real\_homing\_node looked the arm link up as \code{real\_<arm>\_<link>} and the wearer part up as a bare \code{torso}. The real stack publishes the whole URDF under \code{real\_}, so the wearer clearance guard resolved 0 of 9 parts in every run that has ever been made.} \\
\addlinespace[3pt]
\end{longtable}

\section*{2026-08-22 \textnormal{\small\textcolor{srlgrey}{--- 15 commits}}}
\addcontentsline{toc}{section}{2026-08-22}
\begin{longtable}{@{}p{0.055\textwidth}p{0.885\textwidth}@{}}
\endfirsthead\endhead
{\scriptsize\ttfamily 3c1eb46} & \textbf{The pinned wrist cost a factor of seven, and freeing the roll was worth nothing}\newline{\small "The arms barely move" has had four answers in this repository and none of them carried a size. Measured offline, wearer floor 0.15 m enforced under every policy, walking outward until the first refusal: from the arm's own HOME EE -- far outboard, roomy pinned 0.460 m mean reach over 28 walks, 6 wearer-bound free 0.547 m 6 from the WORK POINT (0.400, 0.175, 1.120) -- where the tasks run pinned 0.065 m mean reach over 12 walks, 11 of 12 wearer-bound cone15 0.319 m 7 cone45 0.440 m 5 free 0.450 m 4 The mechanism is n [...]} \\
\addlinespace[3pt]
{\scriptsize\ttfamily 25a28e2} & \textbf{Every joint parks a third of a degree short, and I called that speed-independent}\newline{\small recordings/baselines/arm\_directional\_calibration.json has held 36 real runs since 2026-08-21 -- both arms, six directions, 100 mm each, photographs in recordings/calibration\_frames/ -- and no code in this repository read it. The end effector lands 7.2 mm RMS from where it was sent. In joint space 211 of 252 per-joint terminal errors sit at +/-0.25..0.30 deg with the sign a coin flip: a BOUND, not a spread. A joint parking just inside a band, on the side it arrived from.} \\
\addlinespace[3pt]
{\scriptsize\ttfamily 6b10bf5} & \textbf{It can say what is on the table, and the grasp pipeline's floor could never fire}\newline{\small "It is unable to check what's on the table, it will just move here and there." That was accurate, and it was a MISSING piece rather than a broken one. There was a detector that finds a NAMED thing, a planner that grasps ONE given cloud, and a task layer commanding coordinates written down in advance. Nothing looked at a surface and enumerated what was standing on it, so a wrong name produced a move to somewhere plausible and empty.} \\
\addlinespace[3pt]
{\scriptsize\ttfamily bc99719} & \textbf{The window said VR only, and RViz was embedding an invisible helper}\newline{\small The button audit read 203/203 the whole time these were live, because it presses buttons and none of these are about buttons. RVIZ WAS NEVER ACTUALLY EMBEDDED. \_rviz\_windows() matched any X window with "rviz" in its line and the caller took the LOWEST id -- which is "Qt Selection Owner for rviz2", invisible. The geometry dump measured the CONTAINER and read 1005..1554 exactly as designed, every check passed, and the real RViz sat in its own window on top of the GUI.} \\
\addlinespace[3pt]
{\scriptsize\ttfamily 70ef8ae} & \textbf{LeRobot is worth its dataset format and not its policies, and the docs to run all this}\newline{\small Installed lerobot 0.6.1, exported 23 episodes / 20 440 frames of the real teleop recording as a LeRobotDataset, and wrote down the honest reading. THE CHANNEL GATE IS THE INTERESTING PART. The newest channel baseline says the master arm was 5 ALIVE, 1 INTERMITTENT, 6 INCOHERENT, 2 DEAD. A dead channel is a column of one repeated value; exported without comment it becomes a zero-variance feature indistinguishable from a joint held still on purpose.} \\
\addlinespace[3pt]
{\scriptsize\ttfamily 0c069e4} & \textbf{The hand lands 33 mm from where it was told, and the wearer moves faster than the guard can hear}\newline{\small Asked "how good is the system", this repository could produce a hundred numbers and no budget. A budget is what tells you which number to work on. scripts/measure\_control\_budget.py composes the recorded baselines into one; it measures no new hardware and prints UNMEASURED where nothing has been measured, because a budget with a guessed term hides the thing you should go and measure. WHERE THE HAND ENDS UP: 33.4 mm worst case against a 30 mm grasp gate.} \\
\addlinespace[3pt]
{\scriptsize\ttfamily 1fceffc} & \textbf{There was no motion planner, and we were reading seven numbers off the arm}\newline{\small Two findings, and they are most of "everything works in simulation and then real hardware messes up". THIS PROJECT HAD NO MOTION PLANNER. ik\_follower\_node solves ONE IK per commanded pose and clamps the joint step per cycle -- pose streaming with a rate limiter. There is no OMPL in the command path, no move\_group.plan(), no Cartesian path service; the only OMPL reference in the tree is a vendored Kinova launch file this project does not run. MoveIt is used for /compute\_ik and the planning scene, and nothing else.} \\
\addlinespace[3pt]
{\scriptsize\ttfamily 384db95} & \textbf{Each joint was walked to the IK solution on its own, and the hand went 51 mm somewhere nobody asked for}\newline{\small \code{ik\_follower\_node.clamp\_towards()} was the whole of motion generation for teleoperation. It clamps EACH JOINT INDEPENDENTLY to \code{max\_step\_rad} per control cycle, so a joint needing 0.10 rad arrives in one cycle while one needing 0.50 takes two: the joints are at different fractions of their own travel at every instant, the path through joint space is not the straight line the IK solution implies, and the end effector traces a curve nobody asked for. On a rig whose grasp capture gate is 30 mm.} \\
\addlinespace[3pt]
{\scriptsize\ttfamily dc726df} & \textbf{The diagnosis reproduced to the decimal place, and then four things wrong with the instrument}\newline{\small \code{scripts/measure\_teleop\_motion.py}. It reproduces the recorded diagnosis exactly before it is allowed to report anything else -- max 135.3 mm, mean 45.1 mm on the left arm's shipped home plus [0.50 -0.30 0.10 0.25 -0.05 0.15 0.02] at 50 Hz -- because a baseline you cannot reproduce is not a baseline.} \\
\addlinespace[3pt]
{\scriptsize\ttfamily e3e4956} & \textbf{The follower generates motion through the generator now, in BOTH branches -- the fast path was the unguarded one}\newline{\small \code{clamp\_towards} is replaced in \code{ik\_follower\_node}, and the replacement is in both branches on purpose. It used to be two: a solution within \code{max\_step} was published RAW -- a 0.35 rad jump in one 20 ms cycle, 17.5 rad/s against a 1.3963 rad/s joint limit -- and only a solution outside it went through \code{clamp\_towards}. So the path with "no guard needed" was the faster one. \code{motion\_generator:=} is a parameter and a launch argument on both stacks. \code{ruckig} is the default in every mode.} \\
\addlinespace[3pt]
{\scriptsize\ttfamily d22d832} & \textbf{The audit's launch stub reimplemented the front of on\_launch, so nothing added there could ever be checked}\newline{\small THE GUI RULE. \code{How the arms move (next launch)} on SET UP: which generator the next launch uses, what each costs in plain words and millimetres, and a button that runs the whole comparison offline with no stack up. It applies to the NEXT launch and says so -- every parameter in \code{ik\_follower\_node} except \code{motion\_enabled} is read once at startup, so a control that appeared to change a running follower would be a control that changed nothing.} \\
\addlinespace[3pt]
{\scriptsize\ttfamily 9490a41} & \textbf{The excursion settles at 68 mm rather than going away, so it was never a step size}\newline{\small The write-ups. \code{docs/system/24\_motion\_planning.md} moves Ruckig from "not yet wired; next" to what it measured, with the three decisions worth carrying forward and what is still unmeasured. \code{docs/system/findings.md} carries the full account, including the five things that were wrong in my own work on the way -- four of which would have passed a test suite -- and the six definitions of CONTINUOUS\_IDX. docs/ENGINEERING_LOG.md gains one row.} \\
\addlinespace[3pt]
{\scriptsize\ttfamily 6753293} & \textbf{The pads stopped 13.45 mm short of every object T0, T2 and T3 declared, and the picture agreed with the task the whole time}\newline{\small The largest single term in \code{scripts/measure\_control\_budget.py}'s error budget, against a 30 mm grasp capture gate, and it survived because it cancelled everywhere it was looked at. \code{clip\_tasks.PAD\_OFFSET\_BY\_ARM} was two hand-recorded world vectors of magnitude 0.1118 m. \code{grasp\_frames.PAD\_MID\_EE} is the same quantity in the end effector frame, measured by /compute\_fk on the two finger-tip links on 2026-08-18: 0.09833 m, both arms agreeing to 0.000 mm.} \\
\addlinespace[3pt]
{\scriptsize\ttfamily 29ccfe0} & \textbf{The pad midpoint is a curve, not a constant, and staging went to a sixth copy of the home pose}\newline{\small Asked to take the accuracy as far as it goes and then check that everything works. The error budget was the map. RSS of the measured terms 18.64 mm -> 12.91 mm WORST CASE (they add) 33.39 mm -> 19.99 mm against a 30 mm gate \#\# THE PAD MIDPOINT DEPENDS ON HOW OPEN THE HAND IS The Robotiq 85 is a FOUR-BAR LINKAGE: its fingers swing, they do not translate.} \\
\addlinespace[3pt]
{\scriptsize\ttfamily 7f13d57} & \textbf{The recording sweep could not record: how to run it, and what it was refusing about}\newline{\small \code{scripts/record\_everything.sh} drives the whole verification set one mode-task cell per stack, resumable, in \code{mode\_upstreams.MODE\_ORDER}. One cell per \code{sim\_session} invocation and not one long sweep, for the reason \code{record\_t1\_both\_stages.sh} is written the same way: churning processes against a live stack degrades the DDS graph until a fresh process receives nothing while every node is still up. \code{--fresh} only stops it SKIPPING.} \\
\addlinespace[3pt]
\end{longtable}

\section*{2026-08-23 \textnormal{\small\textcolor{srlgrey}{--- 16 commits}}}
\addcontentsline{toc}{section}{2026-08-23}
\begin{longtable}{@{}p{0.055\textwidth}p{0.885\textwidth}@{}}
\endfirsthead\endhead
{\scriptsize\ttfamily 83e467a} & \textbf{The scene drew T1's cubes 11.43 mm from where the hand was told to close, and filed the clip "1 OBJECT NEVER MOVED"}\newline{\small MY OWN REGRESSION, CAUGHT BY THE RECORDING IT BROKE, and it is the third time this one function has produced that sentence. \code{t1\_task.ee\_for} moved onto the gripper opening a 40 mm cube needs earlier today -- the Robotiq's fingers swing on a four-bar, so the pad sits 0.09833 m from the wrist wide open and 0.10976 m closed on a cube. \code{clip\_scene .\_pad\_offset} stayed on the WIDE-OPEN value. So the task commanded a wrist that puts the pads on the cube and the scene drew the cube 11.43 mm from there.} \\
\addlinespace[3pt]
{\scriptsize\ttfamily 12db4b7} & \textbf{T2's grippers had never closed, and three instruments were reporting defects that were not there}\newline{\small THE WHOLE VERIFICATION SET IS RE-RECORDED against the fixed geometry: 17 cells, eight camera angles each, every mode. The old set is archived, not deleted. grasp success, all five modes 100\% positioning error at closure 0.000 mm T1 placement 319.7 -> 56.3 mm (var 264 -> 0.59) verify\_rviz\_clips 33 of 33 verify\_gripper\_motion 8 of 12 -> 12 of 12 status\_table REAL GAPS NONE \#\# T2'S GRIPPERS HAVE NEVER CLOSED, AND THE TRAY WAS NOT WHY docs/ENGINEERING_LOG.md records "T2-1 'held by BOTH grippers' is false and has been for the life of [...]} \\
\addlinespace[3pt]
{\scriptsize\ttfamily 8e6b7ee} & \textbf{The one section whose job is "which file to open first" named four files that had not existed for weeks}\newline{\small \code{make\_clip\_index.py} generates \code{recordings/verification/INDEX.md} and its own header says it is generated so that it "cannot claim a clip that is not there". That was true of the table of every clip and false of the shortlist above it: the five recommended paths were written out by hand against the A/B/C task set, and four of them -- \code{06\_full\_autonomy/B/S1\_bimanual}, \code{01\_master\_teleop/A/S1\_single\_arm}, \code{02\_vr\_teleop/B/S1\_bimanual} and \code{06\_full\_autonomy/A/S1\_single\_arm} -- were retired when the MSc tasks replaced th [...]} \\
\addlinespace[3pt]
{\scriptsize\ttfamily 6d20f42} & \textbf{The planner has never known the table exists, and everything runs on coordinates a file declares}\newline{\small The operator watched the re-recorded clips: the arm "does not scan the area properly in a straight line, it moves in random directions", and asked for what a 3-D printer does before it prints -- probe each point in order. Then the part that matters more: the real cell has a different table and different objects, teleoperation is unusable because there is no proper planning, and the robot should calibrate the environment from every source it has, build a 3-D map, and let the path and grasp layers adapt to THAT.} \\
\addlinespace[3pt]
{\scriptsize\ttfamily c1e39a3} & \textbf{Objects come out of the picture now, not out of the point pile}\newline{\small The calibration stage found things on a surface the way this repository always has: fit a plane to the fused cloud, cluster the points above it by Euclidean distance. Run against the T1 scene it reported object 0 at (-0.356, 0.471, 1.272), 140 mm across the jaws, NOT graspable where two 40 mm cubes are. They sit on a 60 mm pitch, so the gap between two faces is exactly 20 mm, which is the cluster distance. The map refused an object that does not exist, with a confident reason naming the gripper.} \\
\addlinespace[3pt]
{\scriptsize\ttfamily 5621520} & \textbf{The camera pose has to come off the frame, and the viewing angle is a measurement}\newline{\small Two defects in the live half of the calibration stage, and the second one is the reason the first three rounds of this looked like a perception problem. THE POSE. env\_probe.camera\_pose looked the transform up with Time() -- the LATEST -- which is where the camera is now, not where it was when the frame was taken. The map came back with the cubes REPLICATED, one copy per view, each displaced by about one cell of the sweep: 13 objects for a 6-object scene, and two runs of the same command gave different maps.} \\
\addlinespace[3pt]
{\scriptsize\ttfamily 4592bd1} & \textbf{A pick and place planned from the map the robot measured}\newline{\small Everything else in this repository that picks something up was told where it is. run\_abc builds T1's path from T1\_CUBES, four coordinates in a source file. That is verified N=10 over densified paths and it survives exactly one cell. pick\_from\_map reads what the arm SAW during the calibration sweep and picks an object out of it. Move the cube and the plan moves. Put a different object down and it appears in the list. Nothing in it knows what T1 is. And it is the first caller joint\_planner.VoxelWorld has ever had.} \\
\addlinespace[3pt]
{\scriptsize\ttfamily 74ce92c} & \textbf{The window says what the robot is doing, and the clip shows what it found}\newline{\small THE GUI RULE: a capability that only exists in a terminal does not exist for the person running the session. MAP THE ENVIRONMENT is the top panel of the RUN tab -- above the fold, checked by screenshot and not assumed, because this window has already had one defect where every mode button was below it. It carries a live phase banner off /robot\_say, which is the robot's own sentence and the same one the log gets, so the window and the terminal cannot disagree.} \\
\addlinespace[3pt]
{\scriptsize\ttfamily 306d3d8} & \textbf{Write down what was measured, and what is still not true}\newline{\small docs/system/25\_calibrate\_then\_plan.md is the new stage end to end: the sweep, the segmenter, the map, the seam to the planner, the pick, the voice, the accuracy scored against the renderer's own geometry, and the viewing-angle table. 22\_grasping.md gets the finding itself, with sources -- including that LeRobot has no perception stage at all, so there is nothing there to copy for "where is the cube in metres". Its loop is teleoperate, record, train, deploy, and the policies eat raw video.} \\
\addlinespace[3pt]
{\scriptsize\ttfamily 2cb7d60} & \textbf{The recordings of the calibration sweep and the pick it planned}\newline{\small recordings/calibration/20260823\_125722/ holds the run as it happened: calibrate.mp4 the arm sweeping the workspace in straight serpentine rows pick.mp4 the pick and place, planned from the map that sweep made narration.txt all 24 sentences the robot published, timed world\_map.json the map the pick was planned from summary.txt what it is and what it is not The map is drawn in both clips: the measured surface, the two cubes green at 40 and 41 mm, the two pads red and labelled with the reason the jaws cannot close on [...]} \\
\addlinespace[3pt]
{\scriptsize\ttfamily c0ef577} & \textbf{Ignore the regenerated occupancy cloud}\newline{\small 1.1 MB of points written by every calibration run. pick\_from\_map names its absence rather than silently planning without the table.} \\
\addlinespace[3pt]
{\scriptsize\ttfamily 0dce39b} & \textbf{The occupancy clouds really are out now, and the run logs are in}\newline{\small The previous commit's message said world\_map\_occupancy.npy was deliberately not committed. It was committed -- 2.3 MB of it, twice. The ignore rule I wrote covered recordings/baselines/ and the calibration run directories are somewhere else, so the message described an intention rather than the tree. Found by reading \code{git ls-files} instead of trusting what I had just written. Removed from tracking, ignored properly, and left on disk.} \\
\addlinespace[3pt]
{\scriptsize\ttfamily d60ed9a} & \textbf{The wrist holds still now, the sweep is a volume, and both arms build one map}\newline{\small Four things the operator asked for, and the first is the one that mattered. THE GRIPPER NEVER HELD STILL. \code{look\_at\_quat} aimed the wrist AT each cell, so the orientation changed at every step of the grid: the arm reoriented between neighbouring cells 90 mm apart, swung continuously, and no two views were taken from the same attitude -- which means a residual wrist error was a DIFFERENT error in every frame and could never cancel or be measured. A 3-D printer's probe does not re-aim.} \\
\addlinespace[3pt]
{\scriptsize\ttfamily 2de054b} & \textbf{The two-arm volume sweep, recorded}\newline{\small recordings/calibration/20260823\_202409/ is the run described in the previous commit, filmed: calibrate.mp4 both arms sweeping the volume, three fixed-attitude passes each over two layers -- 144 cells planned, 70 reachable, 67 photographed pick.mp4 the pick and place, planned from the map that sweep made narration.txt 89 sentences, timed, including which pass and which attitude run.log every cell, every refusal, every region count world\_map.json the map itself summary.txt coverage, stillness refusals and the measure [...]} \\
\addlinespace[3pt]
{\scriptsize\ttfamily fe4dc5a} & \textbf{Write down the printer-probe sweep, and why the scene camera cannot help}\newline{\small The measured figures for the two-arm volume sweep, the coverage split into the two honest numbers, and the evidence rule. And one thing checked rather than assumed: THE SCENE CAMERA CANNOT CONTRIBUTE TO THE MAP AT ALL, and it is not a wiring gap I could close. \code{scene\_camera\_node} publishes image\_raw and camera\_info and NO DEPTH -- it is a monocular colour camera. One colour image gives a direction to an object and never a distance, so it cannot place a point in metres.} \\
\addlinespace[3pt]
{\scriptsize\ttfamily 1f66b98} & \textbf{The swept volume is measured now, and the sweep clip shows the map arriving}\newline{\small Two things left over from the last pass, both of which I had written down as not done. THE VOLUME WAS A GUESS. x 0.30-0.58, y 0.30-0.52, and 74 of the 144 cells planned inside it came back "outside the arm's envelope". The obvious reading is that the arms are small. The real one is that a guessed box sat badly inside a much larger one. \code{measure\_reachable\_volume.py} solves IK at the shipped attitude over a deliberately over-wide grid, 924 poses per arm, commanding nothing: reachable at SOME facing+layer x 0.10 ..} \\
\addlinespace[3pt]
\end{longtable}

\section*{2026-08-24 \textnormal{\small\textcolor{srlgrey}{--- 14 commits}}}
\addcontentsline{toc}{section}{2026-08-24}
\begin{longtable}{@{}p{0.055\textwidth}p{0.885\textwidth}@{}}
\endfirsthead\endhead
{\scriptsize\ttfamily fc0a08e} & \textbf{The surface was never 5 mm low. I was asking about the wrong place.}\newline{\small The map reported the support surface at 1.2448 m against a table top of 1.2500 -- 5.2 mm low -- with 78\% of all points below the true top. I wrote a plane-refinement to correct that bias before checking whether it was one. It was not. \code{Map.surface\_z} returned \code{offset / n\_z}: the plane's Z-INTERCEPT AT x = y = 0. A fitted support plane comes out slightly tilted, so its intercept is not its height anywhere a table exists -- and on this rig the origin is INSIDE THE WEARER, half a metre from the nearest tabletop.} \\
\addlinespace[3pt]
{\scriptsize\ttfamily 013d3d4} & \textbf{A map is a fact about the scene, and the scene moves}\newline{\small \code{pick\_from\_map} read the stored map and planned a grasp off it with NO FRESHNESS CHECK AT ALL. A map measured in a previous session looks exactly as authoritative as one measured a minute ago -- same file, same fields, same confident object list -- and the arm would have been driven at coordinates for a scene that no longer exists. The age is printed on every run now, and EXECUTING off a map older than 30 minutes is refused with the age, the command to re-measure, and the override.} \\
\addlinespace[3pt]
{\scriptsize\ttfamily b9d0614} & \textbf{It grasps from ABOVE now, and the angle is measured per object}\newline{\small The operator: "please try to pick the objects from above, not from inside the table or the side." They were right about what it looked like. The shipped anchor approaches at about -31 degrees -- from BELOW and behind -- so the wrist drove up out of the table toward the object. On the last recorded pick the wrist sat at z = 1.227 with the cube at 1.282: 55 mm UNDER it. docs/ENGINEERING_LOG.md says top-down is impossible: "top-down on a surface: 0 of 840 cells, any table height 0.70-1.10, either arm".} \\
\addlinespace[3pt]
{\scriptsize\ttfamily a2a7360} & \textbf{Re-measure only what changed: 18 minutes to 32 seconds, and the scene camera does the noticing}\newline{\small The operator: "if the table is the same then the scene should immediately understand and the robot should only look for the changes and grab that." THE FULL SWEEP IS THE RIGHT THING TO DO ONCE. It measures the support surface, learns which viewpoints the arm can reach, and learns which few of them see the whole table. None of that changes while the table stays put. What changes is the objects -- and sweeping the volume again to find that out is why the calibration was unusable.} \\
\addlinespace[3pt]
{\scriptsize\ttfamily 2a05963} & \textbf{Pick accuracy against ground truth: 11.3 mm mean, 4 of 4 inside the gate}\newline{\small WHY THE EXISTING TABLE DOES NOT ANSWER THIS. \code{accuracy\_table.py} reports 100\% grasp and 0.000 mm positioning error for every mode. Both are sound and neither says whether the gripper is on the OBJECT, because the arm is scored against THE COORDINATE IT WAS GIVEN. Hand it a coordinate 50 mm from the cube and it will hit that coordinate to 0.000 mm and the table will say 100\%.} \\
\addlinespace[3pt]
{\scriptsize\ttfamily e40cf6f} & \textbf{Shared autonomy: the confidently-wrong rate goes from 47.5\% to zero}\newline{\small \code{IntentEstimator} is a Bayesian posterior over WHICH OBJECT the operator is reaching for, and everything in modes 03 and 04 hangs off it being right. Its inputs are the object positions. That is the whole connection to calibration: it can only be as right about "which cube" as it is about "where the cubes are". THIS CANNOT BE MEASURED IN THE SIMULATION AS IT STANDS.} \\
\addlinespace[3pt]
{\scriptsize\ttfamily 4fec074} & \textbf{The table is found now, not assumed to be where the old one was}\newline{\small The operator asked whether this adapts to a table kept anywhere or only to the one in front. It was only the one in front, and the split is worth stating. ADAPTED ALREADY: the table's HEIGHT (measured, 1.2502 against a true 1.2500), its SIZE, and every object's position, size and count. DID NOT ADAPT: its POSITION. \code{BOUNDS} was a constant -- a box in front of the wearer, x 0.325..0.750, y 0.15..0.65 -- the height guess was the constant 1.25, and the camera's base heading pointed forward.} \\
\addlinespace[3pt]
{\scriptsize\ttfamily 82ebf84} & \textbf{A table beside the wearer, 300 mm lower, found to 0.3 mm -- and plainer words on the window}\newline{\small THE TABLE BESIDE THE WEARER was the untested case I had named. Tested now, by shifting the whole rendered scene (+0.35, -0.35, -0.30) while the task files stayed put -- so the table is at the wearer's SIDE, close in, and at a different height, and every declared coordinate is wrong: table found: surface z = 0.9503 m, tilt 0.03 deg, spanning x -0.758..0.995 y 0.019..0.540 left arm will sweep x 0.325..0.850 y 0.019..0.540 right arm will sweep x -0.758..-0.450 y 0.019..0.540 truth (after the shift) measured surface 0. [...]} \\
\addlinespace[3pt]
{\scriptsize\ttfamily 9b8c632} & \textbf{A table beside the wearer, and the last of the spurious objects}\newline{\small BESIDE, AND 300 mm LOWER. The untested case from the last commit, tested: the whole rendered scene shifted (+0.35, -0.35, -0.30) with the task files untouched, so the table is at the wearer's SIDE, close in, at a different height, and every declared coordinate is wrong.} \\
\addlinespace[3pt]
{\scriptsize\ttfamily 5f7db07} & \textbf{The whole session, filmed: scan, pick, quick check, and the window}\newline{\small \code{record\_session.py} films the SEQUENCE, because the sequence is the claim. A scan alone says the perception works; a pick alone says the arm works; the quick check means nothing except after a scan it can compare against. Filmed separately each would be true and the thing that matters -- a table measured once and then kept current in thirty seconds -- would not be shown at all.} \\
\addlinespace[3pt]
{\scriptsize\ttfamily 085f7c3} & \textbf{The gripper cameras were subscribed to topics nothing publishes}\newline{\small The operator: "on the gui i never saw gripper camera connected". They were right, and it could never have worked. The panel subscribed to \code{/<arm>\_wrist\_camera/image\_raw} and \code{/wrist\_mounted\_camera/<arm>/image}, and NOTHING IN THIS REPOSITORY PUBLISHES EITHER -- checked, no \code{create\_publisher} anywhere uses those names. Every camera producer and consumer here uses \code{/<arm>\_camera/color/image\_raw}: the mock, \code{env\_probe}, the calibration sweep, \code{verify\_colour\_vision}, \code{verify\_scan\_view}.} \\
\addlinespace[3pt]
{\scriptsize\ttfamily 8d5dc3c} & \textbf{Do not touch the splitter after RViz has been adopted}\newline{\small The operator reported the window not opening. It does open here -- X reports the main window at 1920x1060, Map State IsViewable, with RViz embedded and the connection checks cycling normally -- so what they hit was almost certainly my leavings: a GUI left running from my own testing, which \code{start\_gui.sh} correctly refuses to duplicate and which looks exactly like "not opening", and the wedged shared memory their connection panel was reporting as \code{shm=bad, daemon=bad}. The machine is clean now.} \\
\addlinespace[3pt]
{\scriptsize\ttfamily 5f315b8} & \textbf{Any shape, any colour: tested rather than reasoned about}\newline{\small The operator asked whether the table can be any shape and colour. Both had an argument -- the plane fit is geometry-only and the segmenter cuts on edges, not on particular colours -- and an argument is not a measurement. Two hooks on the mock make them measurable: SRL\_SCENE\_TABLE\_RGB recolours the table, SRL\_SCENE\_EXTRA\_BOX adds solids so it need not be a rectangle. COLOUR: NO EFFECT. Rendered dark grey (0.30, 0.30, 0.34), close to the background the scene is drawn on -- the hard case, not a token change.} \\
\addlinespace[3pt]
{\scriptsize\ttfamily daf7b91} & \textbf{The window opens behind your terminal, and the refusal only gave you a pid}\newline{\small The operator could not get the window to load. It loads: captured on the real WSLg display, the whole thing renders -- the map panel with every button legible, all four modes, RViz embedded and running at 31 fps, the E-STOP bar. X reports it at 1920x1060, Map State IsViewable, with the RViz window correctly reparented inside it at 549x819. Nothing is broken. TWO REASONS SOMEBODY SEES NOTHING ANYWAY, and both are now handled.} \\
\addlinespace[3pt]
\end{longtable}

\section*{2026-08-25 \textnormal{\small\textcolor{srlgrey}{--- 1 commit}}}
\addcontentsline{toc}{section}{2026-08-25}
\begin{longtable}{@{}p{0.055\textwidth}p{0.885\textwidth}@{}}
\endfirsthead\endhead
{\scriptsize\ttfamily 3d33f8a} & \textbf{The window was six pixels too wide, and WSLg answered by not drawing it}\newline{\small The operator saw a penguin icon in the taskbar that did nothing when clicked. Nothing was wrong with the window. Captured from the running process, every panel renders: the numbered RUN workflow, all four modes, the gripper camera panels, RViz embedded and reporting ready at 31 fps, the whole thing repainting at 3.0 ms/frame. \code{Gui.\_\_init\_\_} sized itself with \code{self.resize(1920, 1060)} -- a constant, never measured against a display. This screen is 1920 wide, so the right edge landed at 1926.} \\
\addlinespace[3pt]
\end{longtable}

